\documentclass[11pt]{article}

\usepackage[margin=1in]{geometry}
\usepackage{amsmath,amssymb,amsthm,mathtools}
\usepackage{enumitem}
\usepackage{hyperref}

\hypersetup{
  colorlinks=true,
  linkcolor=blue,
  citecolor=blue,
  urlcolor=blue
}

\newtheorem{definition}{Definition}
\newtheorem{proposition}{Proposition}
\newtheorem{assumption}{Assumption}
\theoremstyle{remark}
\newtheorem{remark}{Remark}

\newcommand{\Effects}{\mathcal{E}}
\newcommand{\LTL}{\mathrm{LTL}}
\newcommand{\LTLf}{\mathrm{LTL}_f}
\newcommand{\AP}{\mathrm{AP}}
\newcommand{\Close}{\mathrm{Close}}
\newcommand{\Lang}{\mathcal{L}}
\newcommand{\Nat}{\mathbb{N}}
\newcommand{\True}{\mathsf{true}}
\newcommand{\Frame}{\mathfrak{F}}
\newcommand{\Design}{\eta}
\newcommand{\Adm}{\mathrm{Adm}}

\title{Frame-Relative Counterfactual Causes for Explaining Synthesized
Temporal Behavior\\(with an automata-theoretic trace-level test and a
strategy-level reverse parity-game lift)}
\author{}
\date{}

\begin{document}
\maketitle

\begin{abstract}
Linear Temporal Logic (LTL) synthesis usually moves from a temporal
specification to an automaton, then to a parity game, then to a winning
strategy, and finally to a finite-state controller such as a Mealy machine.
This forward transformation is intentionally many-to-one: different LTL
specifications can yield the same controller, and different representations of
winning strategies can describe the same behavior. Therefore, a controller by
itself does not determine a unique original specification or a unique temporal
reason for its behavior.

This document describes a reverse use of parity-game structure. The goal is not
to reconstruct the original LTL formula and not to synthesize a new controller.
Instead, given a Mealy machine, an optional original LTL specification, a
reverse-causal frame, a temporal effect to explain, and an observed execution,
the goal is to identify a minimal temporal cause.
The input is best viewed as
\((M,\varphi,\Frame,E,\tau)\), where the Mealy machine \(M\) already supplies the
input/output partition, \(\varphi\) remains ordinary LTL specification
information, and \(\Frame\) is the causal frame used by the reverse procedure.
The central question is: among the non-vacuous temporal conditions whose removal
prevents the chosen effect from occurring on the closest admissible
counterfactual executions, which is the most preferred (by default the smallest,
most readable) under an explicitly chosen preference order? The
resulting construction treats the Mealy machine as a fixed strategy, builds an
automata-theoretic arena around the effect selected by the frame, and tests and
synthesizes candidate explanations by counterfactual reasoning over that arena.
The trace-level necessity test is, in every regime, a one-player shortest-edit
search plus non-emptiness check (the frozen machine has no free move); emptiness
alone decides only non-vacuity, while the counterfactual-dependence clause
requires optimizing to the closest removal and then evaluating the effect. A genuinely two-player
\emph{reverse parity game} appears for a distinct, strategy-level notion --- can
the controller drop a commitment and still \emph{guarantee} the effect against all
environments? --- which the frame reaches by re-opening the system's strategic
choices. This strategy-level notion additionally requires a synthesis artifact ---
the game \(G_\varphi\), the frozen strategy \(\sigma\), and a declared editable
vertex set \(V_{\Frame}\) --- which \(M\) alone does not determine; it is not a
function of \((M,\tau,E)\). We therefore describe the method as automata-theoretic counterfactual
extraction, with a reverse parity-game generalization for strategy-level causes,
and are careful throughout to keep the trace-level (one-player) and strategy-level
(two-player) notions distinct.
\end{abstract}

\section{Motivation}

The standard synthesis pipeline can be summarized as
\[
  \varphi_{\LTL}
  \longrightarrow
  A_{\varphi}
  \longrightarrow
  G_{\varphi}
  \longrightarrow
  \sigma
  \longrightarrow
  M .
\]
Here \(\varphi_{\LTL}\) is an LTL specification, \(A_{\varphi}\) is an
automaton recognizing the traces that satisfy the specification,
\(G_{\varphi}\) is a parity game encoding the interaction between the
environment and the system, \(\sigma\) is a winning strategy for the system,
and \(M\) is a finite-state implementation of that strategy.

This forward pipeline loses information. A Mealy machine records how a
controller reacts to inputs, but it does not record which LTL formula motivated
that behavior. Many specifications may be realized by the same machine. Also,
many different winning strategies may have implementations that are
behaviorally equivalent. Consequently, the reverse problem
\[
  M \longrightarrow \varphi_{\LTL}
\]
is not well-defined unless additional information is supplied.

The more useful reverse question is explanatory rather than reconstructive. We
assume that we are given
\[
  (M,\varphi,\Frame,E,\tau),
\]
where \(M=(S,s_0,I,O,\delta,\lambda)\) is the Mealy machine, which already fixes
the input/output partition with \(I\cap O=\emptyset\); \(\varphi\in\LTL(I\cup O)\)
is optional original specification information; \(\Frame\) is the reverse-causal
frame; \(E\) is the temporal effect to be explained, a single (possibly
conjunctive) \(\omega\)-regular trace property drawn from a space \(\Effects\) of
candidate effects; and \(\tau\) is the actual
observed execution, the trace. The task is to compute or check a temporal property as a
candidate cause, determine whether it belongs to the valid-cause set, and then
select a minimal cause. Appendix~\ref{app:phi} explains the relation
between \(\varphi\) and \(\Frame\), and Appendix~\ref{app:effect} explains the
possible choices of \(E\), including the TempoBench-compatible restriction.

To separate validity from minimality, let \(\mathcal{K}_{\Frame}\) be the space
of candidate temporal explanations allowed by the causal frame. Also let
\(E_{\Frame}\) be the effect selected by \(\Frame\); if no scoping is used,
then \(E_{\Frame}=E\). The set
\[
  \mathsf{Causes}_{M,\Frame,\tau}(E_{\Frame})
  \subseteq \mathcal{K}_{\Frame}
\]
contains all valid causes of \(E_{\Frame}\). An element
\[
  C_{\mathrm{valid}}\in \mathsf{Causes}_{M,\Frame,\tau}(E_{\Frame})
\]
is a valid cause when:
\begin{enumerate}[label=\arabic*.]
  \item \(C_{\mathrm{valid}}\) holds on the observed execution.
  \item \(E_{\Frame}\) holds on the observed execution.
  \item \textbf{Non-vacuity.} At least one admissible counterfactual execution
        violates \(C_{\mathrm{valid}}\); that is, removing the candidate is
        actually realizable under the frame.
  \item If \(C_{\mathrm{valid}}\) is removed in the closest admissible
        counterfactual executions, then \(E_{\Frame}\) no longer holds.
\end{enumerate}
Condition~3 is essential: without it, any property that no admissible
counterfactual can violate (for example a tautology, or any invariant that all
admissible traces already satisfy) would pass condition~4 vacuously.

It is worth stating once, plainly, which notion of ``cause'' conditions~1--4
formalize. They encode \emph{contingent counterfactual necessity}: \(C\) holds
on the actual execution (1--2), removing \(C\) is realizable under the frame
(3), and on the closest admissible executions that remove \(C\) the effect
fails (4). This is a \emph{counterfactual-necessity} reading of causation, not
a sufficiency criterion and not Halpern--Pearl actual causation, which would add
a minimal-sufficient-set requirement. We commit to this necessity reading
throughout, and the preference order below is the only place where additional
selection (primarily preferring smaller, readable causes; see
Section~\ref{sec:formal}) enters. To avoid an overloaded vocabulary we reserve
``logically weaker'' for a larger language \(\Lang(C)\) and ``insufficient'' for
a candidate that, on its own, does not force the effect; these are different
notions and we keep them lexically distinct. We avoid the unqualified phrase
``weakest necessary condition'': in standard logic the \emph{strongest} necessary
condition for \(E_{\Frame}\) (equivalently the weakest \emph{sufficient}
condition) is \(E_{\Frame}\) itself, while the weakest necessary condition is
\(\True\), so the phrase is doubly misleading for what we do; our sense is
counterfactual necessity under \(\preceq_\tau\).

A minimal cause is then a selected element
\[
  C^\star\in \mathsf{Causes}_{M,\Frame,\tau}(E_{\Frame})
\]
that is \(\triangleleft\)-minimal: no valid cause is strictly preferred to it
under the explicitly chosen preference order \(\triangleleft\). Several
incomparable minimal causes may exist, so we speak of \emph{a} minimal cause
rather than \emph{the} minimal cause; \(\triangleleft\) is assumed
well-founded, so that whenever the valid-cause set
\(\mathsf{Causes}_{M,\Frame,\tau}(E_{\Frame})\) is non-empty at least one
\(\triangleleft\)-minimal cause exists.

The intended output is not merely the statement that the controller wins. The
intended output is a temporal explanation of what had to remain true in order
for the chosen effect to occur.

\section{Related Work and Contributions}
\label{sec:related}

This work sits at the intersection of three lines of research, and our claims
should be read as deltas relative to them rather than as a fresh start.

\paragraph{Counterfactual and actual causation.} The counterfactual-necessity
core is the Lewis--Stalnaker tradition: a cause is what, on the closest worlds
where it is absent, removes the effect, with centering and a limit assumption on
the similarity order. Halpern--Pearl actual causation refines this with a
minimal-sufficient-set requirement and a contingency quantifier; we deliberately
adopt only the necessity half (Section~\ref{sec:overdetermination}) and are
explicit about what this excludes (preemption discrimination, sufficiency).

\paragraph{Temporal/$\omega$-regular causality.} The closest technical relatives
are automata-theoretic accounts of causes for $\omega$-regular effects in
reactive systems, notably the line of work synthesizing $\omega$-regular causes
for $\omega$-regular effects (downward-closed causes via game/automaton
constructions), which also underlies the ground truth of the TempoBench benchmark
we target. Relative to that work, we do \emph{not} claim a new causality-synthesis
algorithm; our contribution is the \emph{framing} layer below.

\paragraph{Reactive synthesis and parity games.} We use only standard facts:
LTL${}\to{}$deterministic parity automaton translation, parity games being
determined and solvable in quasi-polynomial time (and in
$\mathrm{UP}\cap\mathrm{co\text{-}UP}$), and HOA as the interchange format for
$\omega$-automata. These are background, not contributions.

\paragraph{Contributions.} Against this backdrop, the specific claims of this
document are: (1) an explicit \emph{causal frame} $\Frame$ that factors a temporal
causal query into background, effect scope, intervention model, candidate
vocabulary, similarity order, and preference order, together with canonical
default frames (Section~\ref{sec:canonical-frames}) that pin all components but
the declared vocabulary, grammar, and tie-breaks; (2) a clean separation of \emph{trace-existential}
necessity (one-player shortest-edit plus non-emptiness) from \emph{strategy-guarantee}
necessity (a two-player reverse parity game), with the latter formalized and given
a soundness statement (Section~\ref{sec:rpg}); (3) a precise, decidable
anti-circularity condition (Assumption~\ref{ass:anticirc}) with a mandatory effect
cone (Assumption~\ref{ass:cone}) for general effects, barring logical entailment
and point-position restatement; and (4) a finite-trace
specialization that aligns the framework with TempoBench-style TCE
(Appendix~\ref{app:tempobench-alignment}). The novelty is organizational and
definitional --- a reusable problem statement and its regimes --- not a new
hardness result or solver. We flag this explicitly so the contribution is judged
as such.

\section{Background: What a Parity Game Does}
\label{sec:background-parity}

A parity game is a mathematical model of an infinite interaction between two
players. In synthesis, the two players are usually:
\begin{itemize}
  \item the environment, which chooses inputs, and
  \item the system, which chooses outputs.
\end{itemize}

Formally, a parity game is a tuple
\[
  G = (V, V_0, V_1, T, \Omega),
\]
where \(V\) is the set of vertices, partitioned into the system vertices
\(V_0\) and the environment vertices \(V_1\) with \(V=V_0\cup V_1\) and
\(V_0\cap V_1=\emptyset\); \(T\subseteq V\times V\) is the set of edges; and
\(\Omega:V\to\Nat\) assigns a priority to each vertex. We require the edge
relation to be \emph{total}, i.e.
\[
  \forall v\in V\;\exists v'\in V.\;(v,v')\in T,
\]
so that every vertex has a successor and every play extends to an infinite path
(dead ends are excluded).
The game is played on this directed graph. Each vertex belongs to one of the two
players. When the token is on an environment vertex, the environment chooses the
next edge. When the token is on a system vertex, the system chooses the next
edge. Since reactive systems are meant to run forever, a play is an infinite
path through this graph:
\[
  v_0 v_1 v_2 \cdots .
\]

The word ``parity'' refers to the winning condition. Each vertex has a priority,
which is a natural number. During an infinite play, some priorities may appear
only finitely many times and some may appear infinitely often. In this document
we use the following convention:
\[
  \text{the system wins iff the least priority seen infinitely often is even.}
\]
This convention is arbitrary but standard; some texts use the greatest priority
seen infinitely often instead. What matters is that the convention is fixed.

Parity games are useful for LTL synthesis because LTL formulas can be translated
into deterministic parity automata. The automaton tracks whether the infinite
input/output trace satisfies the formula. The parity game combines this
automaton with choices made by the environment and the system. Solving the game
answers the following yes/no question: does there exist a system strategy
\(\sigma\) such that, for every environment behavior \(\rho\), the play that is
induced by \(\sigma\) and \(\rho\) satisfies \(\varphi\)? Formally:
\[
  \exists \sigma\;\forall \rho\;.\;
  \mathrm{Play}(\sigma,\rho)\models \varphi .
\]
If the answer is yes, the strategy can be implemented as a finite-state
controller, often represented as a Mealy machine.

\section{The Forward View}

Let
\[
  M = (S,s_0,I,O,\delta,\lambda)
\]
be a Mealy machine. The set \(S\) contains the internal states of the machine,
and \(s_0\in S\) is the initial state. The set \(I\) contains input atomic
propositions controlled by the environment, and \(O\) contains output atomic
propositions controlled by the system.

At time \(t\), the environment chooses a set
\[
  i_t \in 2^I
\]
of inputs that are true at that time. The Mealy machine responds with a set
\[
  o_t = \lambda(s_t,i_t) \in 2^O
\]
of outputs that are true at that time, and then updates its state by
\[
  s_{t+1} = \delta(s_t,i_t).
\]
The visible behavior at time \(t\) is the combined input/output valuation
\[
  a_t = i_t \cup o_t \in 2^{I\cup O}.
\]
The resulting trace is the infinite word
\[
  \tau = a_0 a_1 a_2 \cdots .
\]

This model explains why the forward direction loses explanatory detail. The
machine tells us which output follows from each state and input, but it does not
tell us which temporal property should be regarded as the reason for an
observed effect. For example, the same machine might satisfy both a request
response property and a weaker eventual-grant property. Looking only at the
machine does not determine which of those properties is the intended cause.

\section{The Reverse Question}
\label{sec:reverse-question}

The reverse construction proposed here is not
\[
  M \longrightarrow \varphi .
\]
Instead, it is
\[
  M + \varphi + \Frame + E + \tau
  \longrightarrow
  \text{minimal causal temporal property } C .
\]
In the rest of this document, \(\varphi\) keeps its ordinary synthesis meaning:
it is an LTL specification over \(I\cup O\). The reverse-causal construction
introduces a separate object, the causal frame \(\Frame\). The frame records how
the explanation problem is scoped: which effect is selected, which assumptions
are background, which propositions are relevant, which candidate causes are
allowed, and which counterfactual interventions are admissible. When the frame
is derived from an original specification, we write
\[
  \Frame=\operatorname{Frame}(\varphi;\Design),
\]
where \(\Design\) denotes the additional choices needed for causal explanation.

The formula \(\varphi\) may describe the original specification, but
\(\Frame\) supplies the reverse-causal structure, and \(E\) is the effect under
explanation, while \(\Frame\) may further restrict it to a scoped target
\(E_{\Frame}\) (Section~\ref{sec:formal}). Typical effects include:
\begin{itemize}
  \item an output eventually occurs after a request,
  \item a particular transition or state recurrence happens,
  \item the observed trace satisfies a chosen subproperty,
  \item the product of the machine and an automaton remains in an accepting
        region.
\end{itemize}

Choosing \(E\) is essential. If \(M\) realizes all of \(\varphi\) against every
environment, then asking what caused \(\varphi\) may be too broad. A sharper
question is: what caused this effect on this execution, or what rule embedded
in this strategy is responsible for this class of effects?

The reverse parity game asks a counterfactual explanation question:
\[
  \begin{gathered}
    \text{For a candidate cause } C,\\
    \text{do all closest counterfactual plays violating } C\\
    \text{make the effect } E \text{ fail?}
  \end{gathered}
\]
This is different from ordinary synthesis. Synthesis searches for a strategy
that forces a specification to hold. Reverse explanation starts with a fixed
strategy and asks which temporal condition is necessary for the chosen effect.

\section{Step-by-Step Construction}

Throughout this section, \(E\) denotes the frame-selected effect: when the frame
scopes the effect we mean \(E_{\Frame}\), and when no scoping is used
\(E_{\Frame}=E\). We write \(E\) for brevity and switch to the explicit
\(E_{\Frame}\) in the formal development of Section~\ref{sec:formal}.

\subsection{Step 1: Freeze the System Strategy}
\label{step:freeze}

First compile the effect \(E\), or the relevant part of the original
specification, into a deterministic parity automaton
\[
  A_E = (Q,q_0,2^{I\cup O},\delta_E,\Omega).
\]
The automaton reads one input/output valuation at a time. Its state records the
progress of the effect, and its priority function \(\Omega\) determines whether
the infinite trace satisfies \(E\). We use the same parity convention for
\(A_E\) as for games (Section~\ref{sec:background-parity}): a run is accepting
iff the least priority occurring infinitely often is even, and priorities are
assigned to states (state-based acceptance); a transition-based variant assigns
\(\Omega\) to edges instead and changes nothing essential below.

Now form the synchronous product
\[
  P = M \times A_E .
\]
A product state has the form
\[
  (s,q),
\]
where \(s\in S\) is a Mealy-machine state and \(q\in Q\) is an automaton state.
From product state \((s_t,q_t)\), one transition proceeds as follows:
\begin{enumerate}[label=\arabic*.]
  \item The environment chooses \(i_t\in 2^I\).
  \item The Mealy machine produces \(o_t=\lambda(s_t,i_t)\).
  \item The Mealy machine updates to \(s_{t+1}=\delta(s_t,i_t)\).
  \item The automaton updates to
        \[
          q_{t+1}=\delta_E(q_t,i_t\cup o_t).
        \]
\end{enumerate}

The important change is that the system is no longer a free player. Its choices
have already been fixed by \(M\). Consequently the \emph{trace-level} necessity
test --- searching the admissible set for the closest \(\neg C\)-traces and
asking whether they preserve \(E\) --- is a one-player shortest-edit search plus
non-emptiness/reachability check in
\emph{all three} intervention regimes, including strategy-structure: freezing
\(M\) removes the system's free move from this search. A genuine \emph{two-player}
parity game arises only for the distinct \emph{strategy-guarantee} question
(notion N2 of Section~\ref{sec:two-notions}): may the system \emph{re-open} its
commitments and still force \(E\) against every environment? That question is
posed over the original synthesis game \(G_\varphi\) with \(\sigma\) partially
re-opened, and we reserve the term \emph{reverse parity game} for it
(Section~\ref{sec:rpg}).

It is therefore cleanest to keep two distinct arenas in mind. The frozen product
\(P=M\times A_E\) introduced here is a \emph{one-player} arena: it underlies the
trace-level necessity test (N1) in all regimes, where the system has no free moves
and the question is a shortest-edit search plus non-emptiness/reachability. The genuinely
two-player arena is the original synthesis game \(G_\varphi\) with the winning
strategy \(\sigma\) partially re-opened (Step~\ref{step:search}, ``causal
dominion extraction''); it underlies the strategy-guarantee question (N2). These
are not the same object, and we are careful below not to present a two-player game
over the frozen product \(P\).

\subsection{Step 2: Choose the Actual Trace}
\label{step:trace}

Causality is evaluated relative to an observation. Let
\[
  \tau = (i_0\cup o_0)(i_1\cup o_1)(i_2\cup o_2)\cdots
\]
be the actual trace generated by \(M\). In implementations, \(\tau\) is often
stored as a lasso
\[
  uv^\omega,
\]
where \(u\) is a finite prefix and \(v\) is a finite loop that repeats forever.

A candidate cause must be true on the actual trace, and the effect must also be
true on the actual trace:
\[
  \tau\models C
  \qquad\text{and}\qquad
  \tau\models E .
\]
This condition prevents vacuous explanations. A property that was not present
in the observed execution cannot explain that execution.

\subsection{Step 3: Define the Allowed Counterfactuals}
\label{step:counterfactuals}

A counterfactual trace is a trace that differs from the actual trace in a
controlled way. The construction depends on which changes are allowed.
This step does not define \(\varphi\). It defines or instantiates
\(\mathcal{I}_{\Frame}\), the intervention component of the causal frame. The
same original LTL formula \(\varphi\) can support different reverse-causal
questions depending on whether \(\Frame\) allows input-only,
output-intervention, or strategy-structure counterfactuals.

\paragraph{Input-only counterfactuals.}
The environment may change inputs, but the Mealy machine must respond according
to its transition and output functions. Counterfactual traces therefore remain
inside the language \(\Lang(M)\) generated by \(M\). This is appropriate when
the question is: which environment behavior caused the machine to produce the
effect?

\paragraph{Output-intervention counterfactuals.}
The explanation game may allow direct changes to outputs. Such interventions
can break the ordinary behavior of \(M\), because the changed output may not be
the output prescribed by \(\lambda\). This version asks which system decisions
were necessary for the effect, but it no longer restricts attention to genuine
executions of the original machine.

\paragraph{Strategy-structure counterfactuals.}
The game may allow changes to strategy commitments, transitions, priorities, or
winning-region choices in the parity-game representation. This is the most
parity-game-native version. It asks which commitments of the winning strategy
prevent the environment from reaching a losing region. The resulting cause may
be a temporal rule such as
\[
  G(req \rightarrow X\,pending)
\]
or
\[
  G(pending \land ack \rightarrow X\,\neg pending).
\]
These are not just single events. They are rules about how the strategy must
continue to behave over time.

\paragraph{Admissible counterfactuals per regime.}
Each intervention model \(\mathcal{I}_{\Frame}\) fixes the admissible
counterfactual set \(\Adm_M^{\Frame}(\tau)\), the traces the reverse procedure
is allowed to consider. The three regimes instantiate it differently:
\begin{itemize}
  \item \textbf{Input-only:} counterfactuals must remain genuine machine
        behaviors, so
        \[
          \Adm_M^{\Frame}(\tau)=\Lang(M)\cap\Lang(B_{\Frame}).
        \]
  \item \textbf{Output-intervention:} outputs may be edited away from
        \(\lambda\), so admissible traces leave \(\Lang(M)\). A natural choice is
        \[
          \Adm_M^{\Frame}(\tau)=\mathrm{Edit}_{O}^{\Frame}(\Lang(M))\cap\Lang(B_{\Frame}),
        \]
        where \(\mathrm{Edit}_{O}^{\Frame}\) is the set of traces obtained from
        \(\Lang(M)\) by the output edits the frame permits. Schematically, fix a
        frame-permitted set \(P_{\Frame}\subseteq\Nat\times O\) of
        (timestep, output) positions that may be toggled; then
        \[
          \mathrm{Edit}_{O}^{\Frame}(\Lang(M))
          =\{\,\tau'\mid \exists\,\tau''\in\Lang(M).\ \tau' \text{ and } \tau''
          \text{ agree off } P_{\Frame}\,\},
        \]
        i.e. \(\tau'\) is obtained from some genuine machine behavior by flipping
        only outputs at positions in \(P_{\Frame}\). Inputs are untouched, so a
        frame that permits no positions recovers \(\Lang(M)\).
  \item \textbf{Strategy-structure:} the system's commitments are re-opened, so
        admissibility is defined over plays consistent with some strategy in the
        frame-permitted neighborhood \(\Sigma_{\Frame}\) of the frozen strategy:
        \[
          \Adm_M^{\Frame}(\tau)=
          \{\,\mathrm{Play}(\sigma',\rho)\mid \sigma'\in\Sigma_{\Frame},\ \rho\text{ an environment behavior}\,\}
          \cap\Lang(B_{\Frame}).
        \]
        Schematically, \(\Sigma_{\Frame}\) is a set of strategies that differ
        from the frozen strategy \(\sigma\) (the one implemented by \(M\)) on a
        frame-bounded set of game vertices, e.g. all \(\sigma'\) agreeing with
        \(\sigma\) outside a designated editable vertex set \(V_{\Frame}\subseteq
        V\). By default \(\Sigma_{\Frame}\) ranges over \emph{finite-memory}
        strategies on \(G_\varphi\) (sufficient for \(\omega\)-regular objectives
        on finite arenas), so a witness deviation can always be taken to be a
        finite-memory strategy realizing an accepting lasso
        (Section~\ref{sec:rpg}). We require \(\sigma\in\Sigma_{\Frame}\), which guarantees that the
        actual trace is itself admissible (\(\tau\in\Adm_M^{\Frame}(\tau)\)); see
        Assumption~\ref{ass:wf} below.
\end{itemize}
The input-only case is the default used in the rest of the document unless
stated otherwise. Crucially, a candidate that all admissible traces already
satisfy cannot be a non-vacuous cause in that regime; see the remark on
strategy-level causes in Sections~\ref{sec:trace-strategy}
and~\ref{sec:formal}.

Two of these three regimes do not in fact depend on the actual trace \(\tau\):
in the input-only and output-intervention cases \(\Adm_M^{\Frame}\) is a fixed
language. We nonetheless write \(\Adm_M^{\Frame}(\tau)\) uniformly because
frame-relative interventions \emph{can} restrict to perturbations of \(\tau\)
itself; the finite-horizon TempoBench specialization of
Appendix~\ref{app:tempobench-alignment}, which perturbs only \(\tau_{\le k}\),
is exactly such a case.

\subsection{Step 4: Define Similarity to the Actual Trace}
\label{step:similarity}

Counterfactual causality requires a notion of closeness. Let
\[
  \preceq_\tau
\]
be a preorder over traces, where
\[
  \tau' \preceq_\tau \tau''
\]
means that \(\tau'\) is at least as close to the actual trace \(\tau\) as
\(\tau''\) is.

The primary closeness order used throughout this document is the
\emph{minimal-intervention lexicographic} order. For a fixed actual trace
\(\tau\) and a trace \(\tau'\), let the \emph{change set} be
\[
  \mathrm{Diff}_\tau(\tau')=\{t\in\Nat\mid \tau'(t)\neq\tau(t)\},
  \qquad
  d(\tau,\tau')=|\mathrm{Diff}_\tau(\tau')|\in\Nat\cup\{\infty\},
\]
with the input-only variant \(d_I\) counting only positions whose \emph{input}
valuation differs. Fix once and for all an arbitrary total order \(\lhd_\Sigma\)
on the finite alphabet \(2^{I\cup O}\), and for a trace with finite change set
write \(\mathrm{prof}_\tau(\tau')\) for the \emph{difference profile}: the finite
list \(\langle(t_1,\tau'(t_1)),\ldots,(t_d,\tau'(t_d))\rangle\) of changed
positions in increasing position order, each tagged with its new value. Define
\[
  \tau'\preceq_\tau\tau''
  \quad\Longleftrightarrow\quad
  \bigl(d(\tau,\tau'),\ \mathrm{prof}_\tau(\tau')\bigr)
  \ \le_{\mathrm{lex}}\
  \bigl(d(\tau,\tau''),\ \mathrm{prof}_\tau(\tau'')\bigr),
\]
comparing first the change counts (with \(\infty\) above every natural number),
then, for equal finite counts, the difference profiles lexicographically
(position first, then value under \(\lhd_\Sigma\)). Traces with an \emph{infinite}
change set all share the top count \(\infty\); we do \emph{not} attempt to order
them among themselves and instead treat them as a single \(\preceq_\tau\)-maximal
class that is never selected as a closest removal when any finite-change removal
exists. Accordingly we define \(\Close_{\tau,M}^{\Frame}(\neg C)\) over the
\emph{finite-change} removals (those with \(d<\infty\)) whenever that set is
non-empty.

On the finite-change traces this \(\preceq_\tau\) is a genuine \emph{total order}
(any two are comparable, and equal lex-keys force equal traces because the profile
pins every changed position to its value and \(\tau\) pins the rest), so in
particular it is antisymmetric --- the hypothesis Proposition~\ref{prop:weakest}
uses. Its unique minimum is \(\tau\) itself, since \(d(\tau,\tau)=0\); this is the
centering condition (Assumption~\ref{ass:wf}(ii)). It is \emph{well-founded on the
finite-change traces}: the count key ranges in the well-ordered
\(\Nat\), and within a fixed finite count \(d\) the profiles are length-\(d\)
tuples over \(\Nat\times 2^{I\cup O}\) ordered lexicographically, which is
well-ordered, so there are no infinite strictly descending chains
(Assumption~\ref{ass:wf}(iii)). We stress that the \emph{lexicographic comparison
of infinite-change profiles would not be well-founded} --- it admits infinite
descending chains --- which is exactly why \(\Close\) is restricted to
finite-change removals; for effects whose only removals require infinitely many
edits, \(\Close\) is empty and the candidate fails non-vacuity, the correct
verdict under non-vacuity. This is a \textbf{scope boundary}, stated explicitly:
trace-level necessity (N1) is confined to \emph{co-safety / finite-change} causes.
A genuine \(\omega\)-liveness cause --- one whose removal requires infinitely many
edits, e.g.\ removing \(GF\,req\) from a lasso forces \(FG\,\neg req\) --- has
\(\Close=\emptyset\) and is reported as non-vacuously failing at the trace level;
such causes belong to the strategy level (N2), where ``dropping a commitment'' is
a single strategic deviation rather than infinitely many trace edits
(Section~\ref{sec:rpg}). We make the count the \emph{primary} key on principled grounds: it is the standard
\emph{minimal-intervention} / Hamming-distance notion for trace repairs (fewest
changed positions first), not a choice tuned to any one example. As an
illustration of its effect, in the toy example of
Section~\ref{sec:toy} ``issue the request one step later''
changes strictly more positions than the minimal removal, hence has a larger
count and is \(\preceq_\tau\)-above it, and so does not enter
\(\Close_{\tau,M}^{\Frame}(\neg req)\). The optional \emph{profile} tie-break that
orders same-count removals can make validity hinge on an implementation detail
(which minimum-count removal is selected), so reports that use it must declare it.
The recommended default is therefore the \emph{count-only coarsening} (drop the
profile tie-break), under which \(\Close\) is the full set of minimum-count
removals rather than a single trace and validity quantifies over \emph{all} of
them, removing the tie-break dependence; this is also the order the TempoBench
specialization of Appendix~\ref{app:tempobench-alignment} uses, since temporal
causal credit assignment asks for necessity over \emph{all} minimal-change
removals.

We deliberately do \emph{not} take the first-difference order as primary.
Comparing traces only by the earliest time they differ from \(\tau\) is the order
induced by the \emph{Baire metric} (an ultrametric) on the Cantor space
\((2^{I\cup O})^\omega\) of infinite words: two traces are closer the later they
first differ. That order is only a \emph{preorder} (all traces that first differ
at the same position are tied), it is not antisymmetric, and under it the
later-request counterfactual above would be \emph{tied} with the minimal removal
and would spuriously defeat the intended cause. The pure change-set
\emph{inclusion} order has the dual defect: the minimal removal (edit
\(\{(0,\emptyset),(1,\emptyset)\}\)) and the later-request trace (edit
\(\{(0,\emptyset),(1,\{req\}),(2,\{grant\})\}\)) have incomparable edit sets, so
both would be minimal and the cause would again fail. The count-based order above
avoids both pathologies; Proposition~\ref{prop:weakest} holds for any well-founded
\(\preceq_\tau\) (count-only or count-first), and the toy example is stated under
it.

\begin{remark}[Closest removals are well-behaved for lasso counterfactuals]
\label{rem:lasso-close}
When \(\tau\) is a lasso \(uv^\omega\), \(M\) is finite-state, and the regime is
input-only or output-intervention, every admissible \(\neg C\)-trace is itself a
lasso of bounded period (a path in the finite product
\(M\times A_{\neg C}\times A_{\preceq}\)), and the change count \(d\) is finite
exactly on the eventually-periodic perturbations that agree with \(\tau\) on its
tail. Hence the candidate set
\(\{\tau'\in\Adm_M^{\Frame}(\tau)\mid\tau'\models\neg C\}\) has a minimum-count
sublevel that is finite up to lasso representation. Under the count-only default
\(\Close_{\tau,M}^{\Frame}(\neg C)\) is the full set of minimum-count removals,
computable by a BFS/shortest-edit search over that product; under the optional
profile refinement (a total order) it is the unique \(\preceq_\tau\)-least
element, a single shortest-edit witness. Either way ``closest removal'' is finitely
represented and computable rather than a large, divergent set.
\end{remark}

\begin{remark}[Centering and the strategy-structure regime]
\label{rem:strat-centering}
Strategy-guarantee necessity (N2, Section~\ref{sec:two-notions}) is a
\emph{type-level} property of the controller \(\sigma\) --- ``can \(\sigma\) drop
a commitment and still guarantee \(E\)?'' --- and is centered on the actual
\emph{strategy} \(\sigma\), not on the actual trace \(\tau\). It therefore does
\emph{not} require trace-centering (Assumption~\ref{ass:wf}(ii)), and
Proposition~\ref{prop:rpg-sound} does not use it; the trace-level Lewis--Stalnaker
axiom is simply not the right notion at the strategy level. Trace-centering matters
only for the \emph{token-level} certificate \(C^{\flat}\)
(Proposition~\ref{prop:weakest}): if one wishes to run that construction
\emph{within} a strategy-structure frame (a separate, optional use), the
edited-vertex subset order ties all plays of \(\sigma\) with \(\tau\), so one must
refine it lexicographically (edited-vertex subset, then the trace-level
count-first order) to restore \(\tau\) as unique minimum. Absent that refinement,
Proposition~\ref{prop:weakest} is read only for the input-only/output regimes,
while the strategy-structure regime carries the validity content of
Proposition~\ref{prop:rpg-sound} (which needs no centering).
\end{remark}

Given a candidate cause \(C\), define the closest counterfactuals that remove
\(C\) as
\[
  \Close_{\tau,M}^{\Frame}(\neg C)
  =
  \min_{\preceq_\tau}
  \{
    \tau'\in \Adm_M^{\Frame}(\tau)
    \mid
    \tau'\models \neg C
  \}.
\]
This is the set of admissible counterfactual traces allowed by \(\Frame\) that
violate \(C\) and are minimal under the chosen closeness order, with
\(\Adm_M^{\Frame}(\tau)\) instantiated per regime as in Step~\ref{step:counterfactuals}.

We collect the standing well-formedness hypotheses that several later arguments
silently rely on.

\begin{assumption}[Well-formedness of the counterfactual model]
\label{ass:wf}
Fix \(M\), \(\Frame\), and the actual trace \(\tau\). We assume:
\begin{enumerate}[label=(\roman*)]
  \item \textbf{Actual-world admissibility.}
        \(\tau\in\Adm_M^{\Frame}(\tau)\); the actual trace is itself an
        admissible world. (In the input-only regime this means
        \(\tau\in\Lang(M)\cap\Lang(B_{\Frame})\), so in particular
        \(\tau\models B_{\Frame}\); in the strategy-structure regime it follows
        from \(\sigma\in\Sigma_{\Frame}\).)
  \item \textbf{Centering.} \(\preceq_\tau\) has \(\tau\) as its unique minimum:
        \(\tau\prec_\tau\tau'\) for every admissible \(\tau'\neq\tau\). The
        actual world is strictly closest to itself.
  \item \textbf{Well-foundedness (Limit Assumption).} \(\preceq_\tau\) is
        well-founded on \(\Adm_M^{\Frame}(\tau)\): there is no infinite strictly
        descending chain \(\cdots\prec_\tau\tau_2\prec_\tau\tau_1\).
        Consequently every non-empty admissible subset --- in particular
        \(\{\tau'\in\Adm_M^{\Frame}(\tau)\mid\tau'\models\neg C\}\) when it is
        non-empty --- has \(\preceq_\tau\)-minimal elements, so that
        \(\Close_{\tau,M}^{\Frame}(\neg C)\) is non-empty in that case; and
        moreover every admissible element lies above some \(\preceq_\tau\)-minimal
        element of any admissible set that contains it.
\end{enumerate}
\end{assumption}

Parts (i)--(ii) are the standard centering conditions of Lewis--Stalnaker
counterfactual semantics; they ensure that conditioning on \(\neg C\) is
evaluated against worlds genuinely reachable under the frame and that the actual
trace is the reference point of the similarity order. The two consequences drawn
in part (iii) --- existence of minimal removals and domination of each removal by
a minimal one --- are exactly what the non-vacuity clause and
Proposition~\ref{prop:weakest}(c) rely on. Well-foundedness holds automatically
when the candidate space and admissible set are finitely representable (for
example in the finite-horizon TempoBench-compatible setting of
Appendix~\ref{app:tempobench-alignment}), and it holds for the canonical
count-first lexicographic order of Step~\ref{step:similarity} \emph{when
\(\Close\) is restricted to finite-change removals} (the count key is then
well-ordered and, at each finite count, the difference profiles are; see
Remark~\ref{rem:lasso-close} for the lasso case). We emphasize that on
infinite-change traces the lexicographic profile comparison is \emph{not}
well-founded, so part~(iii) is genuinely a restriction to finite-change removals
(or, for orders other than the canonical one, an explicit hypothesis on
\(\preceq_\tau\)), not a free lunch.

\subsection{Step 5: Check a Candidate Cause (One-Player Test)}
\label{step:check}

Given a candidate \(C\), build automata for \(C\), \(\neg C\), and \(E\). The
counterfactual adversary tries to find a trace satisfying
\[
  \tau'
  \in
  \Close_{\tau,M}^{\Frame}(\neg C)
  \cap
  \Lang(E).
\]
The intersections with \(\Lang(M)\) and \(\Lang(\neg C)\) are not written here
because \(\Close_{\tau,M}^{\Frame}(\neg C)\subseteq\Adm_M^{\Frame}(\tau)\) and
already consists of traces satisfying \(\neg C\); in the input-only regime that
already entails \(\tau'\in\Lang(M)\). Such a trace is a counterexample to the
explanation: it says that \(C\) can be removed while \(E\) still holds on a
closest allowed counterfactual.

This trace-existential test (notion N1 of Section~\ref{sec:two-notions}) is a
\emph{one-player} shortest-edit search plus non-emptiness check in \emph{every}
regime: once \(M\) is
frozen the system has no free move, and the adversary merely searches the
admissible language for the closest \(\neg C\)-traces and asks whether any
preserves \(E\). The genuinely
two-player reverse parity game arises only for the distinct
\emph{strategy-guarantee} notion (N2), treated in Section~\ref{sec:rpg}, not for
this trace test. The explainer wins when the adversary cannot find such a trace. Equivalently, \(C\) passes the
counterfactual test when
\[
  \Close_{\tau,M}^{\Frame}(\neg C) \neq \emptyset
  \quad\text{and}\quad
  \Close_{\tau,M}^{\Frame}(\neg C) \cap \Lang(E) = \emptyset .
\]
The first conjunct is the non-vacuity requirement: there must actually be a
closest way to remove \(C\). Thus a candidate cause is valid exactly when
removing it is realizable and destroys the effect on all closest allowed
counterfactual executions.

\subsection{Step 6: Search for a Minimal Cause}
\label{step:search}

Checking one candidate is not enough. The real goal is to find the best
candidate according to an explicit preference order. Three useful approaches
are available.

\paragraph{Automata-theoretic cause synthesis.}
When the effect is \(\omega\)-regular and the similarity relation is regular
enough to be represented automata-theoretically, the cause can itself be
represented by an automaton. Write
\[
  A=\Adm_M^{\Frame}(\tau),
  \qquad
  F=\{\tau'\in A\mid \tau'\not\models E_{\Frame}\}
\]
for the admissible set and the admissible \emph{effect-failures}, and recall
\(\uparrow_{\preceq_\tau}(F)=\{\tau''\mid\exists\,\tau'\in F.\ \tau'\preceq_\tau\tau''\}\).
The direction of the preorder must be read carefully: because
\(\tau'\preceq_\tau\tau''\) means \(\tau'\) is \emph{closer} to the actual trace
\(\tau\), the upward closure \(\uparrow_{\preceq_\tau}(F)\) is the set of
traces that are \emph{at least as far} from \(\tau\) as some effect-failure.
``Above'' means further away, not closer or better. The synthesized
cause \(C^{\flat}\) is \emph{any} \(\omega\)-regular property whose intersection
with the admissible set is the failure-free neighborhood,
\[
  \Lang(C^{\flat}) \cap A
  =
  \overline{\uparrow_{\preceq_\tau}(F)} \cap A
  =
  A\setminus\uparrow_{\preceq_\tau}(F),
\]
i.e. \(\tau'\models C^{\flat}\) (within \(A\)) exactly when no admissible
effect-failure is at least as close to \(\tau\) as \(\tau'\). Note that this pins
\(C^{\flat}\) only on \(A\); off \(A\) it is unconstrained, so ``\(C^{\flat}\)''
names an equivalence class of properties agreeing on \(A\), not a unique formula.
We speak of \emph{a} canonical certificate, and the definite article below is
shorthand for ``any representative.''

\begin{proposition}[Nearest-failure-anchored cause]
\label{prop:weakest}
Suppose Assumption~\ref{ass:wf} holds, that \(\preceq_\tau\) is a partial order
on \(A\) (antisymmetric, as the count-first lexicographic profile refinement of
Step~\ref{step:similarity} is, being total), that the actual effect holds
(\(\tau\models E_{\Frame}\)), and that there is at least one admissible
effect-failure, \(F\neq\emptyset\). Then:
\begin{enumerate}[label=(\alph*)]
  \item \(C^{\flat}\) satisfies the four valid-cause clauses of
        \(\mathsf{Causes}_{M,\Frame,\tau}(E_{\Frame})\) (actuality, the two effect
        clauses, and non-vacuity); hence \(C^{\flat}\) is a valid cause
        \emph{provided it is admitted to} \(\mathcal{K}_{\Frame}\). As a
        distance-defined certificate, \(C^{\flat}\) need not itself lie in
        \(\mathcal{K}_{\Frame}\) (it may be unreadable or fail effect-disjointness),
        so we do \emph{not} assert \(C^{\flat}\in\mathsf{Causes}\) unconditionally;
        \(C^{\flat}\) is a witness object, not necessarily the returned answer
        \(C^{\star}\);
  \item its closest removals are exactly the globally nearest admissible
        effect-failures,
        \[
          \Close_{\tau,M}^{\Frame}(\neg C^{\flat})
          =\min\nolimits_{\preceq_\tau}(F);
        \]
  \item among valid causes whose closest removals are the globally nearest
        effect-failures, \(C^{\flat}\) is the \emph{smallest-language} one:
        \(\Lang(C^{\flat})\cap A\subseteq\Lang(C)\cap A\) for every such \(C\).
        Equivalently --- and this is not a contradiction --- its admissible models
        \(\Lang(C^{\flat})\cap A\) form the \emph{largest} failure-free
        neighborhood of \(\tau\) (the set of admissible traces strictly closer to
        \(\tau\) than every admissible effect-failure), because the complement
        \(\uparrow_{\preceq_\tau}(F)\) it excludes is the \emph{smallest}
        failure-containing up-set; the cause is language-minimal precisely because
        the failures it must admit are pushed as far out as possible.
\end{enumerate}
\end{proposition}

\begin{proof}
Within \(A\) we have \(\neg C^{\flat}\equiv\uparrow_{\preceq_\tau}(F)\), so the
admissible \(\neg C^{\flat}\)-traces are exactly \(\uparrow_{\preceq_\tau}(F)\).

\emph{(b) and minimality of failures.} We first show
\(\min_{\preceq_\tau}(\uparrow(F))=\min_{\preceq_\tau}(F)\). Since
\(F\subseteq\uparrow(F)\), take \(m\in\min(F)\); if some
\(x\in\uparrow(F)\) had \(x\prec_\tau m\), then \(x\in\uparrow(F)\) gives
\(f\preceq_\tau x\) for some \(f\in F\), whence \(f\preceq_\tau m\); were
\(m\preceq_\tau f\) we would get \(m\preceq_\tau x\), contradicting
\(x\prec_\tau m\), so \(f\prec_\tau m\), contradicting \(m\in\min(F)\). Thus
\(m\in\min(\uparrow(F))\), giving \(\min(F)\subseteq\min(\uparrow(F))\).
Conversely, if \(m\in\min(\uparrow(F))\) then \(f\preceq_\tau m\) for some
\(f\in F\subseteq\uparrow(F)\); minimality of \(m\) forces \(m\preceq_\tau f\),
and antisymmetry gives \(m=f\in F\), so \(m\in\min(F)\). Hence the two minimal
sets coincide, and \(\Close_{\tau,M}^{\Frame}(\neg C^{\flat})
=\min(\uparrow(F))=\min(F)\), which proves (b).

\emph{(a).} By the Limit Assumption and \(F\neq\emptyset\),
\(\min(F)\neq\emptyset\), so \(\Close_{\tau,M}^{\Frame}(\neg C^{\flat})\neq\emptyset\)
(non-vacuity). By (b) every closest removal lies in \(F\), so it fails
\(E_{\Frame}\); hence \(\Close_{\tau,M}^{\Frame}(\neg C^{\flat})\cap\Lang(E_{\Frame})=\emptyset\).
Actuality holds: \(\tau\models E_{\Frame}\) by hypothesis, and \(\tau\models C^{\flat}\)
because by centering (Assumption~\ref{ass:wf}(ii)) no \(f\in F\) (each
\(f\neq\tau\)) satisfies \(f\preceq_\tau\tau\), so \(\tau\notin\uparrow(F)\). All
four valid-cause clauses hold; whether \(C^{\flat}\in\mathsf{Causes}\) then depends
only on whether \(C^{\flat}\in\mathcal{K}_{\Frame}\), as noted in~(a).

\emph{(c).} Let \(C\) be any valid cause with
\(\Close_{\tau,M}^{\Frame}(\neg C)=\min(F)\), and suppose
\(\tau'\in\Lang(C^{\flat})\cap A\), i.e. \(\tau'\notin\uparrow(F)\): no
effect-failure is at least as close as \(\tau'\). If \(\tau'\notin\Lang(C)\),
then \(\tau'\) is an admissible \(\neg C\)-trace, so by well-foundedness
(Assumption~\ref{ass:wf}(iii)) it lies above some \(\preceq_\tau\)-minimal
\(\neg C\)-trace, i.e.\
\(\tau'\in\uparrow_{\preceq_\tau}(\Close(\neg C))=\uparrow(\min(F))
\subseteq\uparrow(F)\) (some closest removal is at least as close as
\(\tau'\)), contradicting \(\tau'\notin\uparrow(F)\). Hence
\(\Lang(C^{\flat})\cap A\subseteq\Lang(C)\cap A\); since \(C\) was arbitrary in
this class, \(C^{\flat}\) has the smallest language among anchored valid causes,
and \(\Lang(C^{\flat})\cap A=A\setminus\uparrow(F)\) is exactly the set of
admissible traces strictly closer to \(\tau\) than every admissible
effect-failure.
\end{proof}

\begin{remark}[Why ``the weakest valid cause'' is the wrong target]
\label{rem:weakest-degenerate}
It is tempting to seek the \emph{extensionally maximal} (largest-language) valid
cause --- we avoid calling it ``weakest'' to prevent clash with Dijkstra's
weakest precondition --- but that object is degenerate and non-explanatory. Removing a single
nearest failure \(f_0\in\min_{\preceq_\tau}(F)\) already gives a valid cause with
language \(A\setminus\{f_0\}\): its only closest removal is \(f_0\), which fails
\(E_{\Frame}\). Such causes are essentially maximal yet say nothing, the choice of
\(f_0\) is arbitrary, and the valid-cause set therefore has no useful greatest
element. The construction does \emph{not} return this; the \emph{answer} is the
\(\triangleleft\)-minimal readable \(C^{\star}\) (below), while \(C^{\flat}\) is a
\emph{certificate} whose language is the maximal failure-free neighborhood of
\(\tau\) (Proposition~\ref{prop:weakest}(c)). \(C^{\flat}\) is a \(\tau\)-specific,
distance-defined object --- ``stay strictly closer to \(\tau\) than every way of
losing the effect'' --- determined only on \(A\), generally \emph{not} a short or
human-readable temporal formula, and not necessarily a member of
\(\mathcal{K}_{\Frame}\); readable, preference-order-minimal explanations are
obtained instead by the search routes below, which optimize \(\triangleleft\) over
\(\mathcal{K}_{\Frame}\). The certificate stays \(\omega\)-regular when
\(\uparrow_{\preceq_\tau}(F)\) is \(\omega\)-regular --- which holds for the
count-first order via the shortest-edit construction of
Remark~\ref{rem:lasso-close} on lassos, and otherwise must be posited as an
explicit automaticity hypothesis on \(\preceq_\tau\) (we do not claim the order
\emph{relation} is automatic in general; see Section~\ref{sec:complexity}). Antisymmetry of \(\preceq_\tau\) is what the proof uses;
\emph{totality} is sufficient but not necessary, so the guarantee holds for the
count-first lexicographic profile refinement of Step~\ref{step:similarity} (which is
total, hence antisymmetric) as well as for any merely antisymmetric partial
order; under the count-only default it holds via the nearest-layer invariance
hypothesis below. For a non-antisymmetric preorder such as the first-difference (Baire)
order the same argument goes
through only under the extra hypothesis that \(E_{\Frame}\) is invariant under
\(\preceq_\tau\)-equivalence on the nearest layer (all traces tied with a nearest
failure also fail \(E_{\Frame}\)); without that hypothesis the conclusion can
fail, which is one reason we do not adopt the Baire order as primary.
\end{remark}

\paragraph{Counterexample-guided search over LTL sketches.}
Another option is to restrict explanations to a readable grammar, such as
\[
\begin{aligned}
  C ::= {}& p
  \mid \neg C
  \mid C\land C
  \mid X C
  \mid F C
  \mid G C
  \mid C\,U\,C ,
\end{aligned}
\]
or to a smaller template family such as
\[
  G(a\rightarrow F b),\qquad
  F(a\land X b),\qquad
  G(a\rightarrow X b),\qquad
  GF a,\qquad
  FG a .
\]
The loop is:
\begin{enumerate}[label=\arabic*.]
  \item propose a candidate \(C\),
  \item check it in the reverse parity game,
  \item if it fails, extract a counterfactual trace \(\tau'\),
  \item refine the candidate,
  \item minimize according to the chosen ordering.
\end{enumerate}
This approach is practical when the output should be a human-readable LTL
formula rather than only an automaton.

\paragraph{Causal dominion extraction.}
A more parity-game-native approach begins with the original parity game
\(G_\varphi\) and the winning strategy \(\sigma\). Restrict the game to plays
consistent with \(\sigma\), then ask which strategy commitments prevent entry
into losing dominions. The extracted cause may be a set of vertices
\[
  D_C\subseteq V
\]
or a set of guarded transitions
\[
  T_C\subseteq T .
\]
The temporal explanation is then obtained by translating those commitments into
conditions over traces. For example, a commitment may become the rule
\[
  G(req\rightarrow X\,pending),
\]
because without that rule the environment can force the play into a region where
requests remain unresolved.

\section{Trace Causes and Strategy Causes}
\label{sec:trace-strategy}

There are two related but distinct explanation targets.

\paragraph{Trace-level cause.}
A trace-level cause explains one observed execution. Given an actual trace
\(\tau\), it finds a property \(C_\tau\) that caused \(E\) on that trace. For
example, if the actual trace contains a request followed immediately by a grant,
then a trace-level explanation may be
\[
  C_\tau = req \land X\,grant .
\]
This says what happened on that execution.

\paragraph{Strategy-level cause.}
A strategy-level cause explains why the controller as a whole realizes a
property. Given \(M\), it finds a temporal rule \(C_M\) that is satisfied by all
traces of the machine and is responsible for the effect or specification. A
typical strategy-level explanation is
\[
  C_M = G(req\rightarrow F\,grant).
\]
This says what the controller must keep doing across all relevant executions.

The strategy-level version is more ambitious. It is not merely asking which
event caused another event. It asks which temporal invariant, recurrence rule,
or response rule embedded in the controller is responsible for success.

\paragraph{Strategy-level causes require a non-input-only frame.}
The following makes precise why a strategy-level cause cannot be tested under an
input-only frame.

\begin{proposition}[Strategy-level causes need a re-opening frame]
\label{prop:strategy-frame}
Let \(C_M\) be a strategy-level cause, i.e. an invariant of \(M\) with
\(\Lang(M)\subseteq\Lang(C_M)\). If \(\Frame\) permits only input-only
counterfactuals, then \(C_M\) is never a valid cause:
\(C_M\notin\mathsf{Causes}_{M,\Frame,\tau}(E_{\Frame})\) for every \(\tau\) and
\(E_{\Frame}\).
\end{proposition}

\begin{proof}
Under the input-only frame
\(\Adm_M^{\Frame}(\tau)=\Lang(M)\cap\Lang(B_{\Frame})\subseteq\Lang(M)
\subseteq\Lang(C_M)\). Hence
\(\{\tau'\in\Adm_M^{\Frame}(\tau)\mid\tau'\models\neg C_M\}=\emptyset\), so
\(\Close_{\tau,M}^{\Frame}(\neg C_M)=\emptyset\) and the non-vacuity clause of
\(\mathsf{Causes}_{M,\Frame,\tau}(E_{\Frame})\) (Section~\ref{sec:formal}) fails.
\end{proof}

Thus \(C_M\) is correctly rejected as vacuous rather than passing the necessity
test trivially, and evaluating a strategy-level cause \emph{requires} a frame
whose intervention model \(\mathcal{I}_{\Frame}\) re-opens the machine's
commitments so that some admissible counterfactual can violate \(C_M\). The
principled regime for this is \emph{strategy-structure} intervention: its
admissible traces are the plays of genuine alternative strategies
\(\sigma'\in\Sigma_{\Frame}\), so a violation of \(C_M\) corresponds to a real
strategic deviation and the test is a counterfactual over strategies. An
\emph{output-intervention} frame also produces traces outside \(\Lang(M)\) and
can therefore falsify \(C_M\), but its edits need not correspond to any
strategy; it is best understood as a separate output/event-causality regime that
asks which individual output decisions were necessary, not as a test of the
strategic rule \(C_M\). We accordingly treat strategy-level causes as living in
strategy-structure frames, and read output-intervention results as output
causality. Input-only counterfactuals can meaningfully test only trace-level
causes.

\section{A Formal Definition}
\label{sec:formal}

Let \(M\) be a Mealy machine and let \(\tau\in\Lang(M)\) be the actual trace.
The formula \(\varphi\in\LTL(I\cup O)\) may be the original LTL specification,
with semantics \(\Lang(\varphi)\subseteq(2^{I\cup O})^\omega\). The
reverse-causal object is the frame
\[
  \Frame
  =
  (B_{\Frame},
   E_{\Frame},
   S_{\Frame},
   O_{\Frame},
   \mathcal{K}_{\Frame},
   \mathcal{I}_{\Frame},
   {\preceq},
   \triangleleft).
\]
Here \(E_{\Frame}\) is the effect selected by the frame and \(S_{\Frame}\) is the
set of selected effect indices (used when \(E=E_1\land\cdots\land E_m\) is
decomposed; see Appendix~\ref{app:phi}). These two components are not
independent: \(E_{\Frame}=\bigwedge_{r\in S_{\Frame}}E_r\) is determined by
\(S_{\Frame}\) and the decomposition of \(E\), so listing both is a deliberate
convenience (one may drop \(E_{\Frame}\) and keep \(S_{\Frame}\), or vice versa).
Further, \(B_{\Frame}\) is the background constraint supplied by the frame,
\(O_{\Frame}\subseteq O\) is the set of relevant outputs,
\(\mathcal{K}_{\Frame}\) is the candidate-cause space, and
\(\mathcal{I}_{\Frame}\) is the allowed intervention model (which fixes
\(\Adm_M^{\Frame}\) as in Step~\ref{step:counterfactuals}). Environment assumptions \(A_{\mathrm{env}}\)
and system guarantees \(G_{\mathrm{sys}}\), when present, are folded into
\(B_{\Frame}\). If there are no background constraints, take
\(B_{\Frame}=\True\). The frame also supplies the family of similarity
preorders \({\preceq}=\{\preceq_\tau\}_\tau\) (one per possible actual trace; the
operative trace \(\tau\) is a separate input that selects \(\preceq_\tau\)) and
the candidate-minimality order \(\triangleleft\), where \(C'\triangleleft C\)
means that \(C'\) is preferred to \(C\).

First define the closest admissible counterfactuals that remove a candidate
cause:
\[
  \Close_{\tau,M}^{\Frame}(\neg C)
  =
  \min_{\preceq_\tau}
  \{
    \tau'\in \Adm_M^{\Frame}(\tau)
    \mid
    \tau'\models \neg C
  \}.
\]
The set \(\Adm_M^{\Frame}(\tau)\) contains the admissible counterfactual traces
allowed by \(\Frame\). In the ordinary input-only case,
\[
  \Adm_M^{\Frame}(\tau)
  =
  \Lang(M)\cap\Lang(B_{\Frame}),
\]
which does not in fact depend on \(\tau\); as noted in
Step~\ref{step:counterfactuals}, the \(\tau\) argument is retained uniformly only
because frame-relative interventions (notably the TempoBench specialization of
Appendix~\ref{app:tempobench-alignment}) may restrict to perturbations of
\(\tau\).

The set of valid causes is:
\[
\begin{aligned}
  \mathsf{Causes}_{M,\Frame,\tau}(E_{\Frame})
  =
  \{\,C\in\mathcal{K}_{\Frame}
  \mid{}&
  \tau\in\Lang(C),\ 
  \tau\in\Lang(E_{\Frame}),\\
  &
  \Close_{\tau,M}^{\Frame}(\neg C)\neq\emptyset,\\
  &
  \Close_{\tau,M}^{\Frame}(\neg C)
  \cap \Lang(E_{\Frame})=\emptyset
  \,\}.
\end{aligned}
\]
The third conjunct is the non-vacuity requirement: there must be at least one
admissible counterfactual that removes \(C\). Without it, any \(C\) that no
admissible trace violates, in particular a tautology or any invariant of \(M\)
under an input-only frame, would satisfy the emptiness condition vacuously.
Thus \(C_{\mathrm{valid}}\) is a valid cause exactly when
\[
  C_{\mathrm{valid}}
  \in
  \mathsf{Causes}_{M,\Frame,\tau}(E_{\Frame}).
\]

The candidate space \(\mathcal{K}_{\Frame}\) must itself be constrained, or the
necessity test admits circular ``causes'' that re-assert the effect. The key
distinction is between \emph{causing \(E\) through \(M\)} --- which is exactly
what a genuine cause does, and must be allowed --- and \emph{restating \(E\)} ---
asserting (part of) the effect at its own evaluation positions, which is
circular. The condition below bans the latter without touching the former.

\begin{assumption}[Anti-circularity of the candidate space]
\label{ass:anticirc}
Every candidate \(C\in\mathcal{K}_{\Frame}\) is \emph{effect-disjoint}, in two
checkable senses:
\begin{enumerate}[label=(\roman*)]
  \item \textbf{No restatement (semantic, point/bounded effects).} For a point
        effect \(E_{\Frame}=X^k o\), \(C\) does not constrain the value of \(o\) at
        position \(k\): \(\Lang(C)\) is invariant under flipping \(o\) at position
        \(k\) (i.e.\ \(\tau'\in\Lang(C)\iff\tau'[k{:}o{\mapsto}\neg o]\in\Lang(C)\),
        flipping only \(o\) at time \(k\)). This is the clean, semantic form of
        ``\(C\) does not assert the effect at its own position''; it admits a
        cooldown cause ``\(o\) at \(t=5\)'' for effect ``\(o\) at \(t=10\)''
        (\(t\neq k\)) and any inputs, while barring \(X^k o\), \(X^k\neg o\), and
        trivial re-encodings such as \(X^k(o\lor\bot)\). For a bounded conjunctive
        effect the masking is applied at each constrained position. \emph{For a
        general \(\omega\)-regular effect} (\(GF\,grant\), \(G(req\to F\,grant)\),
        parity acceptance) there is no single ``effect position'' to mask, so guard
        (ii) and a \emph{mandatory} declared effect cone
        (Assumption~\ref{ass:cone}) apply instead.
  \item \textbf{No pure-logic entailment (all effects).} \(C\) does not force the
        effect as a matter of logic \emph{alone}, independent of \(M\):
        \(\Lang(C)\not\subseteq\Lang(E_{\Frame})\). This is non-inclusion of two
        \(\omega\)-regular languages, decidable by testing emptiness of
        \(\Lang(C)\cap\overline{\Lang(E_{\Frame})}\). Crucially the test is
        \emph{not} relativized to \(\Adm_M^{\Frame}(\tau)\): a candidate that
        forces \(E_{\Frame}\) only \emph{through the machine dynamics} (as every
        real cause does) is \emph{not} excluded; only a candidate that forces
        \(E_{\Frame}\) by pure logic, i.e.\ that is already a logical strengthening
        of the effect, is. This clause is robust to trivial semantic re-encodings
        (e.g.\ \(X^k(o\lor\bot)\equiv X^k o\) is barred) that a purely syntactic
        guard could miss, and it does the work for general effects, where~(i) is
        silent.
\end{enumerate}
We call this ``precise and decidable'' in the scoped sense that~(i) is a decidable
semantic invariance check for point/bounded effects, (ii) is a decidable
\(\omega\)-regular inclusion test for all effects, and the effect cone
(Assumption~\ref{ass:cone}) is a declared finite \(\mathrm{AP}\)-subset; we do
\emph{not} claim a uniform syntactic ``effect-position'' notion for arbitrary
temporal effects, which is why (i) is scoped to point/bounded effects. Guards (i)
and (ii) bar \emph{logical entailment} and \emph{point-position restatement}; they
do not by themselves exclude system-relative aliases or near-restatements of a
general effect, which is why the effect cone below is mandatory for such effects.
\end{assumption}

\begin{assumption}[Effect cone (mandatory for general effects)]
\label{ass:cone}
For every non-point effect \(E_{\Frame}\), the frame declares a finite
\(\mathrm{AP}\)-subset \(\mathsf{Cone}(E_{\Frame})\) of propositions that may
appear in effect-shaped subformulas (outputs and acceptance atoms tied to \(E\)).
Every \(C\in\mathcal{K}_{\Frame}\) must satisfy
\(\mathrm{AP}(C)\cap\mathsf{Cone}(E_{\Frame})=\emptyset\) unless \(C\) is
explicitly marked, in \(\Frame_{\mathrm{strat}}\), as an internal-state observable
(e.g.\ \(\mathit{at\_}s\)) or a declared system-guarantee invariant. This blocks
vacuous near-restatements that name the effect atom --- e.g.\ a bare recurrence
\(GF\,grant\), or a guarded re-assertion of \(grant\) --- for effect \(GF\,grant\)
when \(\mathsf{Cone}=\{grant\}\), while still allowing input causes and, under a
declared regime, strategy-level state observables and system-guarantee invariants
(e.g.\ the controller commitment \(G(req\to F\,grant)\), admitted only by such an
explicit declaration, never at the trace level).
\end{assumption}

This resolves what would otherwise be a fatal circularity in the wrong place. For
the general effect \(F\,grant\), guard~(i) is silent and~(ii) is operative: the
input cause \(req\) satisfies \(\Lang(req)\not\subseteq\Lang(F\,grant)\) (a trace
with \(req\) at time \(0\) and no \(grant\) ever is a model of \(req\) but not of
\(F\,grant\)), so it is admitted. That \(req\) forces \(F\,grant\) \emph{on
\(\Lang(M)\)} is precisely why it is a cause, and the de-relativized~(ii) does not
penalize it. Likewise the TempoBench cause \(\Gamma^\star\)
(Appendix~\ref{app:tempobench-alignment}), a set of input literals that pins \(o\)
at \(k\) \emph{through the functional machine}, satisfies
\(\Lang(C_{\Gamma^\star})\not\subseteq\Lang(X^k o)\) and is admitted (for the point
effect \(X^k o\), guard~(i) masks only \(o@k\), which \(\Gamma^\star\), being over
inputs, does not touch); and the strategy-level invariant \(G(req\rightarrow
F\,grant)\) names the cone atom \(grant\), so when \(\mathsf{Cone}(F\,grant)=
\{grant\}\) it is admitted \emph{only} as a declared system-guarantee invariant
under \(\Frame_{\mathrm{strat}}\) (never at the trace level); given that
declaration it is not a logical restatement either, since a request-free,
grant-free trace satisfies it vacuously while failing the effect. What is
\emph{excluded} is genuine restatement: \(X\,grant\) and
\(req\land X\,grant\) for effect \(F\,grant\) entail \(F\,grant\) logically,
failing~(ii); a bare recurrence \(GF\,grant\) or any candidate naming a cone atom
without a declaration is barred by the cone; and \(X^k o\), \(X^k(o\lor\bot)\) for
a point effect \(X^k o\) are barred by~(i) (and~(ii)). Thus
Assumption~\ref{ass:anticirc} together with the mandatory effect cone bars logical
entailment, point-position restatement, and undeclared use of effect-cone atoms,
independently of the preference order, while admitting all of the document's worked
causes; it does not claim to exclude every system-relative alias, which is why a
near-restatement naming a cone atom requires an explicit declaration.

\begin{definition}[Minimal temporal cause]
\label{def:mincause}
A candidate \(C^\star\in\mathcal{K}_{\Frame}\) is a minimal temporal
cause of \(E_{\Frame}\) on \(\tau\) relative to \(M\), \(\Frame\), and
\(\triangleleft\) iff:
\begin{enumerate}[label=\arabic*.]
  \item \textbf{Validity.}
        \[
          C^\star\in
          \mathsf{Causes}_{M,\Frame,\tau}(E_{\Frame}).
        \]
  \item \textbf{Minimality.} There is no
        \(C'\in\mathsf{Causes}_{M,\Frame,\tau}(E_{\Frame})\) such that
        \[
          C'\triangleleft C^\star .
        \]
\end{enumerate}
\end{definition}

The condition is \(\triangleleft\)-minimality, not a global minimum: when
\(\triangleleft\) is a partial order, several incomparable minimal causes may
exist, and ``minimal cause'' names any one of them. We assume
\(\triangleleft\) is well-founded (which holds automatically when
\(\mathcal{K}_{\Frame}\) is finite), so that whenever
\(\mathsf{Causes}_{M,\Frame,\tau}(E_{\Frame})\) is non-empty a minimal element
exists.

The preference order \(\triangleleft\) must be chosen deliberately, and its
direction must be stated. Recall that validity already encodes necessity
(Section~\ref{sec:formal}); \(\triangleleft\) only selects \emph{among} necessary
causes. An earlier draft of this work took semantic permissiveness (larger
\(\Lang(C)\), ``logically weaker'') as the \emph{primary} key, but that choice is
defective: if \(A\) and \(B\) are each valid necessary conditions and \(A\lor B\)
is too, then \(A\lor B\) is strictly weaker than either, so a permissiveness-first
order drifts toward ever-larger disjunctions of necessary conditions and would
even prefer \(F\,req\) to \(req\). We therefore make \emph{syntactic readability
the primary key and permissiveness only a tie-break}. The canonical order is
\[
  \bigl(\text{smaller formula size},\ \text{smaller temporal depth},\
  \text{more permissive (larger }\Lang(C)\text{)}\bigr),
\]
read lexicographically. For ``formula size'' and ``temporal depth'' to be
well-defined and reproducible, we fix once a candidate grammar (the readable LTL
grammar of Step~\ref{step:search}) and a normal form (negation pushed to literals,
a fixed associativity/commutativity ordering on \(\land,\lor\), no redundant
\(X\)-nesting), and measure size as the node count of the syntax tree in that
normal form; without such a fixed representation ``smallest formula'' is not
reproducible, so the grammar and size measure are part of the frame. Size and
depth are minimized first; among candidates of equal size and depth, the more
permissive (weaker) one is preferred, since at fixed syntactic complexity
committing to less is the better explanation. We do not claim this particular
lexicographic priority (size \(\succ\) depth \(\succ\) permissiveness) is the
uniquely correct one --- it is a deliberate, declared design choice, and the frame
may substitute another well-founded order; what matters is that it be fixed and
stated. This default has three consequences. (i) It avoids overfitting: large conjunctions such as ``the
entire actual trace occurred'' lose on size. (ii) It avoids disjunctive drift:
\(A\lor B\) is syntactically larger than \(A\), so it is chosen only when no
smaller necessary cause exists --- which is exactly the overdetermined case
(Section~\ref{sec:overdetermination}), where the disjunction is the correct answer.
(iii) The degenerate ``weakest'' object of
Remark~\ref{rem:weakest-degenerate} (which excludes a single nearest
counterfactual) has no small readable formula, so it never wins. To keep these
guarantees we require of \(\mathcal{K}_{\Frame}\) that the trace-exclusion
artifact \(A\setminus\{f_0\}\) be inexpressible (which holds for the readable
grammar and template families of Step~\ref{step:search}); we make this an
explicit condition on admissible candidate classes. Throughout we say ``more
permissive'' for larger language and reserve ``insufficient'' for a candidate
that does not by itself force the effect; these are kept lexically distinct.

If the necessary condition lies outside \(\mathcal{K}_{\Frame}\) --- for
instance \(a\lor b\) when \(\mathcal{K}_{\Frame}\) admits only literals --- then
no valid cause exists in the class and the method reports
\(\mathsf{Causes}_{M,\Frame,\tau}(E_{\Frame})=\emptyset\); the remedy is to
enlarge \(\mathcal{K}_{\Frame}\), not to return a spurious singleton.

Three extremal objects appear in this document, and they must not be conflated.
The following table fixes their roles once and for all.
\[
\renewcommand{\arraystretch}{1.3}
\begin{array}{lll}
  \textbf{Object} & \textbf{What it is} & \textbf{Role}\\
  \hline
  C^{\star}\ (\triangleleft\text{-minimal}) &
    \text{syntactically smallest readable valid cause in }\mathcal{K}_{\Frame}
    \text{ (permissiveness as tie-break)} &
    \text{the answer returned}\\
  C^{\flat}\ (\text{Prop.~\ref{prop:weakest}}) &
    \text{smallest-language anchored cause} &
    \text{a certificate, not the answer}\\
  \text{extensionally maximal} &
    \text{largest-language valid cause }(A\setminus\{f_0\}) &
    \text{degenerate, excluded}
\end{array}
\]
The \emph{answer} is \(C^{\star}\): among readable, effect-disjoint candidates,
the order prefers the smallest formula size, then temporal depth, then --- only as
a tie-break --- the more permissive (larger-language) necessary condition. The
\emph{certificate} \(C^{\flat}\) is a different, \(\tau\)-specific distance object
--- the maximal failure-free neighborhood --- that witnesses \emph{why} the test
passes but is generally not a readable formula; it is reported alongside
\(C^{\star}\), not in place of it. The \emph{extensionally maximal} valid cause
(drop a single nearest failure) is rejected outright as non-explanatory: it is
both unreadable and, being huge, loses on size (Remark~\ref{rem:weakest-degenerate}).
There is thus no tension between the size-minimizing \emph{answer}, the
smallest-language \emph{certificate}, and the excluded largest-language artifact:
they are three distinct objects with three distinct purposes.

\subsection{Canonical Frames}
\label{sec:canonical-frames}

The frame \(\Frame\) has many components, and a fair objection is that a
sufficiently free choice of \(\Frame\) could make almost any property a
``cause,'' so that the cause is injected by the analyst rather than extracted
from the system. We answer this not by forbidding configuration but by fixing
a small number of \emph{canonical default frames}; results should be reported
against one of these unless a deviation is explicitly justified (the honesty
requirement of Appendix~\ref{app:phi}).

\begin{description}[leftmargin=2.4cm,style=nextline]
  \item[Standard trace causality \(\Frame_{\mathrm{trace}}\)]
  \(B_{\Frame}=\True\) (or the stated environment assumptions), input-only
  interventions (\(\Adm=\Lang(M)\)), the count-only order
  \(\preceq_\tau\) of Step~\ref{step:similarity} by default (the profile tie-break
  is used only when explicitly declared), effect-disjoint candidates
  (Assumption~\ref{ass:anticirc}) over the observable vocabulary \(I\cup O\), and
  \(\triangleleft\) preferring smaller formula size, then smaller temporal depth,
  then the more permissive necessary condition. This is the default for explaining
  a single observed execution; the toy example uses the optional profile refinement
  for a unique witness.
  \item[Standard strategy causality \(\Frame_{\mathrm{strat}}\)]
  as above but with strategy-structure interventions over a stated editable
  vertex set \(V_{\Frame}\) and the edited-vertex subset order refined by trace
  distance (so \(\tau\) remains the unique \(\preceq_\tau\)-minimum; see
  Remark~\ref{rem:strat-centering}), and effect-disjointness under the mandatory
  effect cone (Assumption~\ref{ass:cone}) with the strategy-level exception for
  declared system-guarantee invariants and internal-state observables, so that
  strategy-level invariants
  such as \(G(req\rightarrow F\,grant)\) become testable via the reverse parity
  game (Section~\ref{sec:rpg}).
  \item[TempoBench-style \(\Frame_{\mathrm{TB}}\)]
  the finite-horizon, point-effect, input-only frame of
  Appendix~\ref{app:tempobench-alignment}.
\end{description}

Two safeguards make these defaults principled rather than tunable to taste.
First, every component except the candidate vocabulary, grammar, similarity
tie-break, and effect cone is fixed once the regime
is chosen, and the vocabulary is constrained by effect-disjointness, so the only
remaining freedom is which observable propositions are eligible --- a modeling
choice that must be declared. Second, under a fixed canonical frame \emph{including
those declarations} the valid-cause
set \(\mathsf{Causes}_{M,\Frame,\tau}(E_{\Frame})\) is determined by \(M\),
\(\tau\), and \(E_{\Frame}\) alone, so the cause is a function of the system and
observation given that declaration, not a free parameter. (Determinacy is thus
conditional on the declared vocabulary, grammar, count-only vs.\ profile
\(\preceq_\tau\), and effect cone, which must be stated.) Frame-relativity then means results are
\emph{conditional on a declared, reproducible setting}, exactly as background
conditions are held fixed in ordinary causal explanation, rather than
cherry-picked per query.

\section{Toy Example}
\label{sec:toy}

Consider the response specification
\[
  \varphi = G(req\rightarrow F\,grant).
\]
It says that whenever a request occurs, a grant must eventually occur. Suppose
the Mealy machine has two states, with initial state \(s_0=idle\):
\[
  idle
  \qquad\text{and}\qquad
  pending .
\]
We fix the machine completely so that later claims about its invariants are
well-defined. The transition and output functions are total:
\[
\begin{array}{rclcl}
  \delta(idle,req)&=&pending, &\quad& \lambda(idle,\cdot)=\emptyset,\\
  \delta(idle,\neg req)&=&idle, &\quad& \lambda(pending,\cdot)=\{grant\},\\
  \delta(pending,\cdot)&=&idle. & &
\end{array}
\]
In words: a request seen in \(idle\) moves the machine to \(pending\) at the next
step; in \(pending\) the machine emits \(grant\) (regardless of input) and
returns to \(idle\); a request that arrives while in \(pending\) is not latched
(the machine still returns to \(idle\)). The output depends only on the state
(\(grant\) holds exactly in \(pending\)), so this toy machine is effectively a
Moore machine; a Mealy machine would in general also let the output depend on the
current input.
Now consider an actual trace beginning with a request followed by a grant,
written as valuations over \(I\cup O\):
\[
  \{req\},\quad \{grant\},\quad \emptyset,\quad \emptyset,\quad \ldots
\]
So the machine is in \(idle\) at time \(0\), enters \(pending\) at time \(1\)
(where \(grant\) holds), and returns to \(idle\) at time \(2\). Let the effect
be
\[
  E = F\,grant .
\]

\paragraph{The minimal trace-level cause is \(req\).}
Take the default input-only frame, so admissible counterfactuals are exactly
machine behaviors, \(\Adm_M^{\Frame}(\tau)=\Lang(M)\), and use the canonical
count-first lexicographic closeness order of Step~\ref{step:similarity}. Consider
the candidate \(req\) (the proposition \(req\) holds at time \(0\)). It holds on
\(\tau\), and the effect \(F\,grant\) holds on \(\tau\). The closest removal is
the behavior with the smallest change count: it sets \(\neg req\) at time \(0\)
and changes nothing else that it need not. Since the machine starts in \(idle\)
and now receives no request at time \(0\), it stays in \(idle\); with the
remaining inputs unchanged (all \(\emptyset\)) it never enters \(pending\) and
never emits \(grant\). Its difference profile is
\(\langle(0,\emptyset),(1,\emptyset)\rangle\), so \(d=2\) (time~\(1\) loses the
\(grant\) the original trace had), and \(F\,grant\) fails on it. We check that
\emph{no} admissible \(\neg req\)-removal of count \(\le 2\) preserves the effect:
to emit \(grant\) the machine must enter \(\mathit{pending}\), which requires a
\(req\) in \(idle\) at some time \(j\ge 1\); the cheapest such trace sets
\(\neg req\) at \(0\), \(req\) at \(1\), and thereby also changes the output at
time~\(2\), giving profile
\(\langle(0,\emptyset),(1,\{req\}),(2,\{grant\})\rangle\) with \(d=3>2\). Hence
every count-minimal \(\neg req\)-removal fails \(F\,grant\), and since
\(\preceq_\tau\) is total the unique closest removal is the \(d=2\) behavior
above. (Under the first-difference/Baire order the later-request behavior would
be \emph{tied} with the minimal removal, and under the pure change-set
\emph{inclusion} order it would be incomparable to it; in both cases it would
enter \(\Close\) and \(req\) would spuriously fail. The count-first order excludes
it; see Step~\ref{step:similarity}.) Hence \(req\) is non-vacuous and
counterfactually necessary: by the definition of Section~\ref{sec:formal},
\(req\) \emph{is} a valid cause.
Because the preference order minimizes syntactic size (and \(req\), a single
atom, is the smallest readable effect-disjoint candidate that holds on \(\tau\)),
the candidate \(req\) is preferred to the larger \(req\land X\,grant\). Therefore,
under this frame, the formally correct minimal trace-level cause is
\[
  C_\tau = req,
\]
not \(req\land X\,grant\). The stronger candidate \(req\land X\,grant\) is in fact
excluded from \(\mathcal{K}_{\Frame}\) before the order is even consulted: it
entails \(F\,grant\) by logic alone (\(\Lang(req\land X\,grant)\subseteq
\Lang(F\,grant)\), since \(grant\) at time \(1\) already gives \(F\,grant\)), so
it fails clause~(ii) of Assumption~\ref{ass:anticirc}; its conjunct
\(X\,grant\) is a restatement of part of the effect. By contrast \(req\) does
\emph{not} entail \(F\,grant\) by logic alone (only through \(M\)), so it is
admitted. Even were \(req\land X\,grant\) admitted, it would be dispreferred to
\(req\): it is syntactically larger, and on the size-minimizing order of
Section~\ref{sec:formal} a single atom beats a four-symbol conjunction.

\paragraph{Why \(req\) can feel uninformative, and what to do about it.}
The intuition that ``a request alone does not explain \emph{why} the machine
granted'' is real, but it is a critique of the input-only \emph{frame}, not of
the candidate \(req\). An input-only frame fixes the machine's transition and
output rules as unchallengeable background facts; relative to those facts, the
only thing that had to be true on the trace for the grant to appear is that a
request arrived, so \(req\) is exactly the right answer. The machine's response
rule does not appear in the cause precisely because the frame has placed it in
the background.

To make that response rule itself a candidate explanation, one must move to a
frame that re-opens it. At the strategy level, the explanatory target is a rule
such as
\[
  C_M = G(req\rightarrow F\,grant),
\]
or, if the internal state \(pending\) is exposed for explanation,
\[
  G\bigl((\mathit{at\_idle}\land req)\rightarrow X\,pending\bigr)
  \land
  G(pending\rightarrow grant).
\]
We must guard the first conjunct by \(\mathit{at\_idle}\): a request arriving in
\(pending\) is not latched on this machine (Section~\ref{sec:toy}), so the
unguarded \(G(req\rightarrow X\,pending)\) is \emph{not} an invariant of \(M\),
whereas the guarded form is. Note also that \(grant\) holds in the same step as
\(pending\) (the machine is Moore-style), so the second conjunct is
\(G(pending\rightarrow grant)\), not \(G(pending\rightarrow X\,grant)\). Together
the two invariants entail
\(G\bigl((\mathit{at\_idle}\land req)\rightarrow X\,grant\bigr)\); hence on every
trace whose requests arrive only in \(idle\) --- which includes the observed
trace and its input-only counterfactuals --- they imply
\(G(req\rightarrow F\,grant)\), the response specification. By
Proposition~\ref{prop:strategy-frame}, such a \(C_M\) is an invariant of \(M\)
and can be tested only under a strategy-structure frame: under input-only
counterfactuals it is vacuous and so fails non-vacuity. Note that \(C_M\) names
the cone atom \(grant\); it is admissible only because \(\Frame_{\mathrm{strat}}\)
declares it a system-guarantee invariant (Assumption~\ref{ass:cone}), not at the
trace level. Given that declaration, \(C_M\) also does not entail \(F\,grant\) by
pure logic (a request-free trace satisfies \(C_M\) vacuously while failing
\(F\,grant\)), so clause~(ii) of Assumption~\ref{ass:anticirc} holds as well.

The reverse game thus distinguishes these targets by changing the frame. Asking
about the observed trace under the input-only frame yields the trace-level cause
\(req\). Asking which strategic rule makes the controller reliably satisfy the
response effect requires a strategy-structure frame and yields the global
response property. The two answers are not competitors at one level; they are the
correct outputs of two different frames.

\section{Overdetermination, Preemption, and the Role of the Background}
\label{sec:overdetermination}

Because validity is pure counterfactual necessity (Section~\ref{sec:formal}), the
method inherits the well-known limitations of necessity under
\emph{overdetermination} and \emph{preemption}, and a reviewer familiar with
Halpern--Pearl actual causation will expect these to be addressed head-on. We do
so here, and explain why the necessity reading is a deliberate design choice
rather than an oversight.

\paragraph{Overdetermination.} Suppose two inputs \(a\) and \(b\) each, on their
own, force the effect \(E=X\,grant\): the machine grants at time \(1\) whenever
\(a\) holds at time \(0\), and also whenever \(b\) holds at time \(0\), and the
actual trace has both. Then neither \(a\) nor \(b\) is counterfactually
necessary: removing \(a\) alone leaves \(b\), which still forces the grant, so the
closest \(\neg a\)-removal preserves \(E\), and symmetrically for \(b\).
Consequently neither singleton is a valid cause. This is not a defect of the
construction but the correct report under necessity: \emph{no single input was
necessary}. What \emph{is} valid is the disjunction's failure set, captured by the
cause
\[
  C = a\lor b,
\]
whose closest removal sets \emph{both} to false and does destroy \(E\). The
method therefore returns the disjunctive necessary condition \(a\lor b\) ---
\emph{provided \(\mathcal{K}_{\Frame}\) admits disjunctions}. If the candidate
class is restricted (e.g.\ to single literals, as in the TempoBench-style frame),
then \(a\lor b\notin\mathcal{K}_{\Frame}\), no singleton is valid, and the method
correctly reports \(\mathsf{Causes}=\emptyset\) (Section~\ref{sec:formal}); the
remedy is to enlarge \(\mathcal{K}_{\Frame}\) to include the disjunction, not to
return an arbitrary singleton. In this symmetric example the certificate
\(C^{\flat}\) of Proposition~\ref{prop:weakest} coincides with \(a\lor b\) on the
admissible set, but we caution that this identification is special to the example:
in general \(C^{\flat}\) is the distance-defined failure-free neighborhood and
need not equal any particular readable disjunction. We state this behavior
explicitly so that overdetermined effects are reported as disjunctive causes (when
expressible), never as an arbitrary choice between symmetric inputs. This does not
conflict with the readability-first preference order
(Section~\ref{sec:formal}), which penalizes disjunctions on size: \(a\lor b\) is
dispreferred to any smaller valid singleton, so it is returned \emph{only} when no
smaller valid cause exists --- exactly the overdetermined case, where there
genuinely is none. (We return a \emph{necessary condition}, not ``the actual
cause''; necessity does not by itself adjudicate which disjunct is the actual
cause.)

\paragraph{Preemption.} If \(a\) forces \(E\) and \(b\) would have forced \(E\)
only had \(a\) been absent, then again \(a\) alone is not necessary on a trace
where both are present; the necessary condition is once more
\(a\lor b\). Distinguishing the preempting cause \(a\) from the preempted backup
\(b\) requires a sufficiency or contrastive notion, which necessity does not
provide. We do not claim to make this distinction; a frame that wants it must
import a Halpern--Pearl-style contingency (holding \(b\) fixed at its actual
value while varying \(a\)), which can be expressed as a background constraint
\(B_{\Frame}\) that pins \(b\), recovering \(a\) as the necessary cause
\emph{relative to that frame}.

\paragraph{Why necessity, and why this is not vacuous.} Two objections must be
separated. First, ``necessity flags trivial conditions such as `the system was
powered on'.'' It does not, because of two safeguards already in the definition:
non-vacuity (Section~\ref{sec:formal}) discards any condition that no admissible
counterfactual can violate, and the background \(B_{\Frame}\) absorbs standing
operating assumptions into the setting rather than the cause (the lamp/switch
discussion of Appendix~\ref{app:phi}). A universally-true operating assumption is
either folded into \(B_{\Frame}\) or fails non-vacuity, so it is never returned.
Second, ``why not sufficiency, as in Halpern--Pearl?'' Halpern--Pearl actual
causation adds a minimal-sufficient-set requirement with an explicit
quantifier over \emph{contingencies} (settings of the other variables): one must
exhibit a witness set whose fixing, under some contingency, both preserves the
effect actually and removes it counterfactually. The complexity of deciding HP
actual causation is, even in the Boolean setting, sensitive to the exact
definition: it is \(\mathrm{NP}\)-complete for the original HP definition in
binary (Boolean) models, \(\Sigma_2^{P}\)-complete for the original definition in
general models, and \(\mathrm{D}_2^{P}\)-complete for the updated definition
(Eiter--Lukasiewicz; Aleksandrowicz--Chockler--Halpern--Ivrii). Lifting any of
these to \(\omega\)-regular effects layers the contingency quantifier on top of
language operations, so HP-style sufficiency is at least as expensive as --- and
conceptually heavier than --- our necessity test (a one-player shortest-edit plus
non-emptiness check, or, at the strategy level, a parity solve followed by a
non-emptiness check; Section~\ref{sec:complexity}). We make only the modest claim that
necessity is \emph{no harder} (and uniform and decidable), not the stronger and
unproved claim that sufficiency is \emph{strictly} harder over \(\omega\)-regular
effects; a precise separation would require fixing an HP-over-\(\omega\)-regular
semantics, which we do not develop here. We therefore target necessity by design, report
overdetermined effects as disjunctive causes, and expose contingency through
\(B_{\Frame}\) when a sharper, sufficiency-like answer is wanted. This keeps the core notion uniform and decidable while leaving
a clearly marked door to contrastive refinements.

\section{Why the Parity Condition Matters}

Many temporal properties cannot be explained by a single finite event. Examples
include:
\[
  GF\,a,\qquad
  FG\,stable,\qquad
  G(req\rightarrow F\,ack).
\]
These formulas refer to recurrence, eventual stability, and repeated response
obligations. A parity automaton is designed to reason about exactly this kind of
infinite behavior.

In the product \(M\times A_E\), the cause may correspond to preserving an
accepting recurrence pattern. For example, the controller may win because every
request eventually drives the product into a strongly connected component whose
recurring priority is accepting. Alternatively, the cause may be an invariant
that prevents the play from entering a losing dominion after a request becomes
pending.

Crucially, the parity machinery is genuinely needed only in the
\emph{strategy-structure} regime. In the input-only and output-intervention
regimes the system has no free move, so (as Step~\ref{step:freeze} explains) the
counterfactual test is a one-player shortest-edit search plus non-emptiness check,
and no two-player game is solved. The remainder of this section makes the two-player object precise
for the strategy-structure regime, where the title's ``reverse parity game'' is
the operative construction.

\section{The Reverse Parity Game}
\label{sec:rpg}

This section gives the two-player object that the input-only and
output-intervention regimes specialize away. (``Reverse'' names the
\emph{explanatory} direction --- running the synthesis machinery backwards, from a
frozen controller to the conditions it needed --- not a retrograde or
backward-time analysis; the game itself is a standard forward parity game on a
partial-commitment arena.) Throughout, fix the synthesis game
\[
  G_\varphi=(V,V_0,V_1,T,\Omega)
\]
of Section~\ref{sec:background-parity}, the winning strategy \(\sigma\)
implemented by \(M\), a strategy-level candidate cause \(C\) with deterministic
parity monitor \(A_{\neg C}=(Q_{\neg C},q^{0}_{\neg C},2^{I\cup O},\delta_{\neg
C},\Omega_{\neg C})\) for its negation, and the effect monitor \(A_E\). The frame
designates an \emph{editable vertex set}
\[
  V_{\Frame}\subseteq V_0
\]
of system vertices whose moves may be re-opened; this is the same
\(V_{\Frame}\) that defines the strategy neighborhood \(\Sigma_{\Frame}\) in
Step~\ref{step:counterfactuals} (\(\sigma'\in\Sigma_{\Frame}\) iff \(\sigma'\)
agrees with \(\sigma\) on \(V_0\setminus V_{\Frame}\)).

\subsection{Two Distinct Necessity Notions}
\label{sec:two-notions}

Before defining the game we separate two notions that the earlier
trace-based definition and the game would otherwise conflate. Fix the
strategy-structure admissible set
\(\Adm_M^{\Frame}(\tau)=\{\mathrm{Play}(\sigma',\rho)\mid\sigma'\in\Sigma_{\Frame},\
\rho\}\cap\Lang(B_{\Frame})\).
\begin{description}[leftmargin=2.2cm,style=nextline]
  \item[Trace-existential necessity (N1)]
  the literal instance of \(\mathsf{Causes}\) (Section~\ref{sec:formal}): the
  \emph{closest} admissible trace satisfying \(\neg C\) must fail \(E\). Here the
  re-opened system and the environment together pick a \emph{single}
  counterfactual trace, so deciding N1 is an \emph{existential} (one-player)
  shortest-edit search plus non-emptiness check over the admissible set, exactly as in
  Step~\ref{step:check}; it is not a genuine two-player game.
  \item[Strategy-guarantee necessity (N2)]
  the genuinely two-player, strategy-level question: can the system
  \emph{drop the commitment} \(C\) and still \emph{guarantee} \(E\) against every
  \emph{admissible} environment? ``Admissible'' means satisfying the background
  \(B_{\Frame}\) (e.g.\ a fairness assumption); quantifying over \emph{all}
  environments, including \(B_{\Frame}\)-violating ones, would let an unfair
  environment trivially falsify any liveness commitment, so both clauses are
  relativized to \(B_{\Frame}\). Assuming the frozen \(\sigma\) secures the
  guarantee and satisfies \(C\) (the strategy-level actuality requirement), \(C\)
  is \emph{strategy-necessary} iff there is \textbf{no}
  \(\sigma'\in\Sigma_{\Frame}\) that simultaneously
  (i)~secures \(E\) on admissible plays,
  \(\forall\rho.\ \mathrm{Play}(\sigma',\rho)\models(B_{\Frame}\rightarrow E)\),
  and (ii)~is not \(C\)-committed on an admissible play,
  \(\exists\rho.\ \mathrm{Play}(\sigma',\rho)\models B_{\Frame}\wedge\neg C\).
  (When \(B_{\Frame}=\True\) these are \(\forall\rho.\,E\) and
  \(\exists\rho.\,\neg C\).) N2 is centered on the actual \emph{strategy}
  \(\sigma\), not on \(\tau\) (cf.\ Remark~\ref{rem:strat-centering}).
\end{description}
N1 and N2 are different questions --- token-level (this trace) versus type-level
(this controller's guarantee), in the philosophy-of-causation terminology --- and
we give them separate treatments. The reverse parity game decides N2; N1 is
decided by the one-player check already described. We do \emph{not} claim a single
parity game decides N1, and the strategy-structure ``closest-removal'' reading of
\(\mathsf{Causes}\) (with the edited-vertex subset order) is an existential
optimization layered on top of the one-player check, not the game below.

\subsection{The Arena, Players, and Winning Condition}

\begin{definition}[Reverse parity game]
\label{def:rpg}
The reverse parity game
\[
  G^{\mathrm{rev}}_{M,\Frame,C,E}
  =(V^{r},V^{r}_{0},V^{r}_{1},T^{r},\Omega^{r})
\]
has vertices \(V^{r}=V\times Q_B\times Q_E\times Q_{\neg C}\), each pairing a game
vertex with the states of deterministic monitors \(A_B\) (for the background
\(B_{\Frame}\)), \(A_E\), and \(A_{\neg C}\). The vertex set is \emph{partitioned}
into the \(\exists\)-player's vertices \(V^{r}_{0}\) and the \(\forall\)-player's
vertices \(V^{r}_{1}\):
\begin{itemize}
  \item \(V^{r}_{0}\) (\(\exists\), the re-opened system): vertices whose
        \(V\)-component lies in the editable set \(V_{\Frame}\), where \(\exists\)
        may take any \(G_\varphi\)-move;
  \item \(V^{r}_{1}\) (\(\forall\)): all other vertices, namely those whose
        \(V\)-component lies in \(V_1\) (environment input choices) \emph{and}
        those whose \(V\)-component lies in the frozen set
        \(V_0\setminus V_{\Frame}\). Frozen vertices are \emph{single-successor}
        (their only outgoing edge is \(\sigma\)'s prescribed move), so assigning
        them to \(\forall\) keeps the tuple a genuine partition while leaving
        play unaffected --- their owner has no choice.
\end{itemize}
We assume the synthesis game carries an explicit \emph{labeling}
\(\ell:T\to 2^{I\cup O}\) of each \(G_\varphi\)-edge by the input/output valuation
produced on that move (the standard labeling of the synthesis game, where a
round consists of an environment input choice followed by the system output);
this is the label the monitors read. Edges \(T^{r}\) are then induced by \(T\) on
the \(V\)-component together with the synchronous updates of each monitor on
\(\ell(e)\) along each \(G_\varphi\)-edge \(e\); the edge relation is total. The
\(\exists\)-player's \emph{objective} is the assume-guarantee condition
\(B_{\Frame}\rightarrow E\): \(\exists\) wins a play \(\pi\) iff
\(\pi|_{I\cup O}\models B_{\Frame}\rightarrow E\). This is \(\omega\)-regular, so
\(\Omega^{r}\) is a parity index for it and \(G^{\mathrm{rev}}_{M,\Frame,C,E}\) is
a genuine parity game; \(A_B\) and \(A_{\neg C}\) are also read off as
\emph{auxiliary} acceptance conditions used in the witness search below.
\end{definition}

\subsection{Soundness}

The decision procedure must be stated with care, because N2 is a
\emph{conjunctive} strategy-existence question --- \(\exists\sigma'\) that both
secures \(B_{\Frame}\rightarrow E\) (a \(\forall\rho\) condition) and admits an
admissible \(\neg C\) play (an \(\exists\rho.\ B_{\Frame}\wedge\neg C\) condition).
\textbf{Caution:} for a parity/B\"uchi objective, staying inside the winning
region forever does \emph{not} imply winning (a play may remain in \(W_\exists\)
yet violate the parity condition), so we do \emph{not} use a ``trap''
multi-strategy of ``all moves staying in \(W_\exists\)''; instead we re-impose the
objective explicitly in the witness search.

\begin{proposition}[Reverse-game soundness for strategy-guarantee necessity]
\label{prop:rpg-sound}
Assume the frozen \(\sigma\in\Sigma_{\Frame}\) secures \(B_{\Frame}\rightarrow E\)
and satisfies \(C\) (the strategy-level actuality requirement). Let \(W_\exists\)
be the \(\exists\)-winning region of \(G^{\mathrm{rev}}_{M,\Frame,C,E}\) for
objective \(B_{\Frame}\rightarrow E\), and let \(\mathcal{A}[W_\exists]\) be the
sub-arena obtained by restricting \(G^{\mathrm{rev}}\) to \(W_\exists\) (keeping
only edges that stay in \(W_\exists\)). Then \(C\) is \emph{not}
strategy-necessary (N2 fails) iff \(\mathcal{A}[W_\exists]\) contains a play
\(\pi\) with
\[
  \pi|_{I\cup O}\models B_{\Frame}\wedge E\wedge\neg C .
\]
Equivalently, \(C\) is strategy-necessary iff the product
\(\mathcal{A}[W_\exists]\times A_B\times A_E\times A_{\neg C}\) has no accepting
lasso for the conjunction \(B_{\Frame}\wedge E\wedge\neg C\) --- a one-player
\(\omega\)-regular non-emptiness check (not mere reachability of a vertex).
\end{proposition}

\begin{proof}
Let \(\Pi_{\mathrm{win}}\) be the plays consistent with \emph{some}
\(\sigma'\in\Sigma_{\Frame}\) that secures \(B_{\Frame}\rightarrow E\). We claim
\(\Pi_{\mathrm{win}}=\{\pi:\pi\text{ stays in }W_\exists\text{ and }\pi\models
B_{\Frame}\rightarrow E\}\). (\(\subseteq\)) A securing strategy keeps every play
in \(W_\exists\) and wins it, so its plays satisfy \(B_{\Frame}\rightarrow E\).
(\(\supseteq\)) Given such a \(\pi\), we may take \(\pi=uv^\omega\) to be an
accepting \emph{lasso}: \(\mathcal{A}[W_\exists]\) is finite and the winning
condition \(\omega\)-regular, so whenever a qualifying play exists a lasso one does.
Define \(\sigma'\) to follow the (ultimately periodic) moves of \(\pi\)
while the history matches \(\pi\), and a fixed \emph{finite-memory} winning strategy
from \(W_\exists\)
after any environment deviation: every \(\exists\)-move on \(\pi\) goes to
\(W_\exists\) (\(\pi\) stays there), and the environment cannot escape
\(W_\exists\). Following a lasso needs only finite memory (the period), and
finite-memory winning strategies exist for parity objectives on finite arenas, so
\(\sigma'\) is finite-memory; hence \(\sigma'\in\Sigma_{\Frame}\) secures
\(B_{\Frame}\rightarrow E\) and has \(\pi\) as a play. Now \(C\) is not strategy-necessary iff some
\(\sigma'\) secures \(B_{\Frame}\rightarrow E\) (i) and admits a play with
\(B_{\Frame}\wedge\neg C\) (ii), i.e.\ iff \(\Pi_{\mathrm{win}}\) contains a play
with \(B_{\Frame}\wedge\neg C\). By the claim such a play stays in \(W_\exists\)
and satisfies \(B_{\Frame}\rightarrow E\); with \(B_{\Frame}\) this gives \(E\), so
it satisfies \(B_{\Frame}\wedge E\wedge\neg C\) and lies in
\(\mathcal{A}[W_\exists]\). Conversely any \(\mathcal{A}[W_\exists]\)-play with
\(B_{\Frame}\wedge E\wedge\neg C\) satisfies \(B_{\Frame}\rightarrow E\) and stays
in \(W_\exists\), hence lies in \(\Pi_{\mathrm{win}}\) and witnesses (i)+(ii). The
lasso/non-emptiness characterization is standard for deterministic
\(\omega\)-automata over a finite arena.
\end{proof}

This decision is sound \emph{and} complete and uses only (a)~a two-player parity
solve for the winning region \(W_\exists\) and (b)~a one-player
\(\omega\)-regular non-emptiness check on \(\mathcal{A}[W_\exists]\). It does
\emph{not} rely on a maximal-permissive-strategy construction, sidestepping the
fact that for parity objectives the ``moves staying in \(W_\exists\)''
multi-strategy is \emph{not} winning (a Bernet--Janin--Walukiewicz most-permissive
\emph{winning} strategy needs memory); the crucial point is that the objective
\(B_{\Frame}\rightarrow E\) is re-imposed in the witness search, so non-winning
plays that merely stay in \(W_\exists\) are excluded. Step~(ii) is a non-emptiness
question for the monitor's \(\omega\)-acceptance condition, \emph{not} reachability
of an accepting vertex. Minimizing the edited-vertex set \(V_{\Frame}\) (the
\(\subseteq\)-closest deviation reading) is a \emph{separate} optimization over
subsets of \(V_0\) layered on top of this decision; it is not encoded in a single
parity game, and we do not claim it is (see Section~\ref{sec:complexity} for its
cost). The input-only and output-intervention regimes are the degenerate case
\(V_{\Frame}=\emptyset\), where \(V^{r}_{0}=\emptyset\), the game has a single
(environment) player, and the whole question collapses to the one-player
shortest-edit plus non-emptiness check of Step~\ref{step:check} for
trace-existential necessity (N1).

\subsection{From Game Commitments to a Trace Cause}
\label{sec:dominion-translation}

Definition~\ref{def:rpg} also makes precise the ``causal dominion extraction'' of
Step~\ref{step:search}. Suppose \(C\) is strategy-necessary (no
\(\mathcal{A}[W_\exists]\)-play satisfies \(B_{\Frame}\wedge E\wedge\neg C\),
Proposition~\ref{prop:rpg-sound}). The witness is the set of frozen commitments
that block \(\exists\): let
\[
  D_C\subseteq V_{\Frame}
\]
be the editable vertices at which \(\sigma\)'s prescribed move is the one whose
re-opening would otherwise let \(\exists\) secure \(B_{\Frame}\rightarrow E\) via a
\(\neg C\)-deviation (i.e.\ the commitments without which
Proposition~\ref{prop:rpg-sound} would find an \(E\)-securing, \(C\)-violating
play in \(\mathcal{A}[W_\exists]\)), and let
\[
  T_C=\{(v,\sigma(v))\mid v\in D_C\}\subseteq T
\]
be the corresponding \(\sigma\)-transitions. When the commitment is a
\emph{safety} rule (take \(\sigma\)'s move at every visit to \(D_C\)), the induced
trace-level cause is the \(\omega\)-regular property
\[
  C
  \;\equiv\;
  G\!\!\bigwedge_{v\in D_C}\!\!
  \bigl(\mathit{at\_}v \rightarrow \mathit{move}_{\sigma}(v)\bigr),
\]
read over traces with the relevant game positions exposed (Appendix's
\(\mathit{at\_s}\) convention): ``whenever the play is at a committed position
\(v\), it takes \(\sigma\)'s prescribed move,'' with negation
\(\neg C\equiv F\bigvee_{v\in D_C}(\mathit{at\_}v\wedge\neg\mathit{move}_\sigma(v))\).
When the commitment is a \emph{liveness} rule --- as in the recurrence example
below, where \(C=GF\,\mathit{at\_serve}\) --- the induced cause instead has the
form \(GF\bigvee_{v\in D_C}\mathit{at\_}v\) (``keep returning to a committed
position''), with negation \(FG\) of its complement. In both cases breaking the
commitment in \(T_C\) is exactly what an admissible \(\neg C\)-deviation does;
this is the formal content of the slogan ``a commitment becomes a rule such as
\(G((\mathit{at\_idle}\wedge req)\rightarrow X\,pending)\).''

\paragraph{Minimal strategy cause.} The strategy level has, so far, only a
\emph{validity} notion (N2), in contrast to the trace level's validity-plus-%
\(\triangleleft\)-minimality (Definition~\ref{def:mincause}). We complete the
symmetry: a \emph{minimal strategy cause} is a \(\subseteq\)-minimal editable set
\(V_{\Frame}\subseteq V_0\) such that the corresponding commitment \(C\) is
strategy-necessary (Proposition~\ref{prop:rpg-sound}); equivalently, a minimal
causal dominion \(D_C\). Unlike the trace-level order, this minimization is not a
single parity solve: it is a search over vertex subsets with the parity decision
as an oracle (Section~\ref{sec:complexity}). We therefore state plainly that the
strategy level delivers \emph{validity} via the reverse parity game and
\emph{minimality} only via this separate, combinatorially harder optimization,
which we do not claim to solve in polynomial time.

\subsection{A Worked Liveness Example}
\label{sec:liveness-example}

The trace-level toy example never exercises the parity condition because its
effect \(F\,grant\) is a \emph{reachability} (co-safety) property --- satisfied as
soon as a grant appears, so its monitor needs no recurrence and the test reduces
to reachability. We now give a recurrence effect that genuinely needs parity. Take
a two-state machine over input \(req\) and output \(grant\) with
\[
  \delta(\mathit{serve},\cdot)=\mathit{wait},\quad
  \delta(\mathit{wait},req)=\mathit{serve},\quad
  \delta(\mathit{wait},\neg req)=\mathit{wait},
\]
\[
  \lambda(\mathit{serve},\cdot)=\{grant\},\quad
  \lambda(\mathit{wait},\cdot)=\emptyset .
\]
Thus, being Moore, the machine emits \(grant\) exactly \emph{while} in
\(\mathit{serve}\) (not the step after), and it enters \(\mathit{serve}\) one step
after a request arrives in \(\mathit{wait}\). The initial state is
\(s_0=\mathit{wait}\). Let the effect be the recurrence
property
\[
  E = GF\,grant
\]
(``grants occur infinitely often''), with actual trace the lasso
\[
  \tau=\bigl(\{req\},\{grant\}\bigr)^\omega ,
\]
which alternates \(\mathit{wait}\xrightarrow{req}\mathit{serve}\xrightarrow{}
\mathit{wait}\) and indeed grants infinitely often.

\paragraph{An environment-fairness assumption is required.} Crucially, no
controller can guarantee \(GF\,grant\) against \emph{every} environment: an
environment that issues only finitely many requests traps the machine in
\(\mathit{wait}\), where \(grant\) is never emitted. The well-posed strategy-level
question therefore carries the standard environment-fairness background
\[
  B_{\Frame} = GF\,req
\]
(``requests keep arriving''), and the guaranteed objective is the assume-guarantee
property
\[
  B_{\Frame}\rightarrow E \;=\; GF\,req \rightarrow GF\,grant ,
\]
which is again \(\omega\)-regular and yields a parity game. Without \(B_{\Frame}\)
the winning region for \(E\) is empty and every candidate is vacuously
``necessary,'' so stating \(B_{\Frame}\) is what makes the example meaningful ---
exactly the overdetermination/vacuity discipline of
Sections~\ref{sec:formal}--\ref{sec:overdetermination}.

The effect monitor \(A_E\) for \(GF\,grant\) is a two-priority parity (B\"uchi)
automaton under the min-even convention: priority \(0\) on \(grant\)-valuations
and priority \(1\) elsewhere, so a run is accepting iff priority \(0\) recurs
infinitely often. (A single priority cannot encode ``infinitely often''; two are
needed.) The product \(P=M\times A_E\) is accepting iff the play visits the
\(\mathit{serve}\)-emitting region infinitely often --- a genuine parity
(recurrence) acceptance, not reachability: a counterfactual that grants a million
times and then stops still \emph{fails} \(E\).

\paragraph{The strategy-level cause is the recurrence commitment.} Expose the
state and consider the candidate
\[
  C = GF\,\mathit{at\_serve}
\]
--- ``the controller returns to the serving state infinitely often.'' (This is
effect-disjoint by Assumption~\ref{ass:anticirc}: \(C\) names \(\mathit{at\_serve}\),
not \(grant\), and \(\Lang(GF\,\mathit{at\_serve})\not\subseteq\Lang(GF\,grant)\)
as a matter of pure logic, since the two coincide only through \(M\).) Under the
standard strategy frame \(\Frame_{\mathrm{strat}}\) with \(V_{\Frame}\) the
\(\mathit{wait}\)-vertices and background \(B_{\Frame}=GF\,req\), we apply
Proposition~\ref{prop:rpg-sound}. The witness search asks for a play in
\(\mathcal{A}[W_\exists]\) satisfying \(GF\,req\wedge GF\,grant\wedge\neg C\),
i.e.\ \(GF\,req\wedge GF\,grant\wedge FG\,\mathit{at\_wait}\). But \(grant\) is
emitted only in \(\mathit{serve}\), so \(GF\,grant\) entails
\(GF\,\mathit{at\_serve}\), which contradicts \(FG\,\mathit{at\_wait}\); no such
play exists, so \(C\) is \emph{strategy-necessary} (N2). (This uses the
\(B_{\Frame}\)-relativized N2: an unfair environment with finitely many requests
would, against the frozen \(\sigma\), yield a play with
\(FG\,\mathit{at\_wait}=\neg C\), but that play violates \(B_{\Frame}=GF\,req\) and
is therefore not an admissible deviation witness --- exactly why clause~(ii) is
relativized to \(B_{\Frame}\).) The extracted \(D_C\) is the set of
\(\mathit{wait}\)-vertices from which \(\sigma\) moves to \(\mathit{serve}\) on a
request --- the accepting-recurrence-preserving commitment, the causal dominion of
Section~\ref{sec:dominion-translation}: removing the transitions \(T_C\) lets the
environment force the play to remain in \(\mathit{wait}\), out of the recurring
accepting region.

By contrast, the \emph{stronger} candidate ``serve \emph{every} request,''
\(G((\mathit{at\_wait}\wedge req)\rightarrow X\,\mathit{at\_serve})\), is
\emph{not} strategy-necessary: a deviation may skip some requests yet serve
infinitely many others, securing \(GF\,req\rightarrow GF\,grant\) on a
\(B_{\Frame}\)-fair play while violating the invariant, so it admits a \(\neg C\)
play and fails the test. This is the
necessity discipline at work --- the necessary commitment is the recurrence
\(GF\,\mathit{at\_serve}\), not the stronger ``every request.'' And a reachability
candidate such as ``\(grant\) at time \(1\)'' is likewise not necessary: shifting
the first grant still leaves infinitely many. Only a recurrence-preserving
commitment passes, which is why the parity (recurrence) condition, not a
reachability check, is the right acceptance here.

\section{Decidability and Complexity}
\label{sec:complexity}

We record the cost of the two basic operations --- checking a candidate and
synthesizing a minimal cause --- in each regime. Let \(|M|\) be the number of
machine states, \(|A_E|,|A_{\neg C}|\) the monitor sizes, and \(d\) the number of
priorities of the relevant parity condition.

We state these as conditional bounds, not as a fully proved complexity theorem;
the hypotheses (monitor sizes, and how \(\preceq_\tau\) is handled) are made
explicit.

\paragraph{Input-only / output-intervention (one-player).} Checking a candidate
\(C\) has two parts (Step~\ref{step:check}): non-vacuity is non-emptiness of
\(\Adm_M^{\Frame}(\tau)\cap\Lang(\neg C)\) restricted to finite-change removals,
and effect failure asks that the closest such removal preserve \(\neg E\). We do
\emph{not} assume the closeness \emph{relation} \(\preceq_\tau\) is recognized by
a fixed automaton (its primary key compares unbounded change counts, which is not
a finite-state relation in general). Instead we use the constructive fact of
Remark~\ref{rem:lasso-close}: for lasso \(\tau\) and finite-state \(M\), the
closest removal is a \emph{shortest-edit} (minimum-count) path in the finite
product \(M\times A_{\neg C}\times A_E\), computable by a shortest-path/BFS search
over that product of size \(O(|M|\cdot|A_{\neg C}|\cdot|A_E|)\); the effect-failure
check then evaluates \(E\) on that witness. This is polynomial in the product
(plus the parity-solving cost \(|V|^{O(d)}\), quasi-polynomial, if \(E\) is a
genuine parity rather than reachability/B\"uchi effect). When \(C\) ranges over
LTL formulas the \emph{deterministic} parity monitors used in the explicit product
are worst-case \emph{doubly} exponential in \(|C|\) (LTL${}\to{}$deterministic
parity is doubly exponential), so the explicit product route is doubly exponential
in the formula size; the underlying language checks (non-emptiness and
effect-failure) are themselves no harder than LTL model checking ---
\textsc{PSpace} in the formula size --- when carried out via \emph{nondeterministic}
monitors rather than explicit determinization. We do \emph{not} claim a uniform
\textsc{PSpace} bound for the full minimum-count validity quantification. The certificate
\(C^{\flat}\) (Proposition~\ref{prop:weakest}) is \(\omega\)-regular and
computable from \(\uparrow_{\preceq_\tau}(F)\) \emph{when} that upward closure is
itself \(\omega\)-regular --- which we obtain for the count-first order via the
shortest-edit construction on lassos, and otherwise must posit as an explicit
automaticity hypothesis on \(\preceq_\tau\).

\paragraph{Minimal-cause search.} Over a finite template family
\(\mathcal{K}_{\Frame}\) (Step~\ref{step:search}), a \(\triangleleft\)-minimal
cause is found by checking each template, i.e.\ \(|\mathcal{K}_{\Frame}|\) checks.
Over the open LTL-sketch grammar the search is unbounded; restricting to formulas
of bounded size keeps it decidable (finitely many candidates) but exponential in
the bound, as in syntax-guided synthesis.

\paragraph{Strategy-structure (two-player).} Deciding strategy-guarantee
necessity (N2) for a \emph{fixed} \(V_{\Frame}\) is the procedure of
Proposition~\ref{prop:rpg-sound}: solve the parity game
\(G^{\mathrm{rev}}_{M,\Frame,C,E}\) (objective \(B_{\Frame}\rightarrow E\)) of size
\(O(|V|\cdot|A_B|\cdot|A_E|)\) with \(O(d)\) priorities to obtain the winning
region \(W_\exists\), then run a one-player \(\omega\)-regular non-emptiness check
for \(B_{\Frame}\wedge E\wedge\neg C\) on \(\mathcal{A}[W_\exists]\) (product with
\(A_B,A_E,A_{\neg C}\); an accepting-lasso search). Parity solving is in
\(\mathrm{UP}\cap\mathrm{co\text{-}UP}\) and quasi-polynomial time, and the
non-emptiness check is polynomial, so the decision is quasi-polynomial. (We do
\emph{not} invoke a permissive-strategy computation; the winning region plus the
re-imposed objective suffices.) Finding a \(\subseteq\)-minimal \(V_{\Frame}\) that
makes \(C\) necessary --- the ``minimal causal dominion'' \(D_C\) --- is a
\emph{separate} optimization over subsets of \(V_0\): naively \(2^{|V_0|}\)
candidate sets, each tested by the quasi-polynomial decision above. We do not
prove a tight class, but record an upper bound: checking that a given
\(V_{\Frame}\) induces necessity is one quasi-polynomial parity decision, and
checking minimality is a co-\(\mathrm{NP}\)-style quantification over its proper
subsets (each a parity decision); so the minimal-dominion decision sits within a
polynomial-hierarchy level relative to a parity oracle (a \(\Sigma_2\)-type
search), naturally encoded as SAT/SMT/MaxSAT over vertex-selection variables with
the parity decision as oracle. Exact completeness is left open.

\paragraph{Finite-horizon TempoBench case.} Here the candidate space
\(\mathcal{K}_{\Frame_{\mathrm{TB}}}\) is finite but \emph{exponential}: there are
at most \(2|I|(k+1)\) timestep-indexed literals, hence up to
\(2^{2|I|(k+1)}\) candidate conjunctions (constraint sets \(\Gamma\)). Checking
one candidate is a finite reachability check over \(A_{\mathrm{HOA}}\) restricted
to the window \(0,\ldots,k\), polynomial in \(|A_{\mathrm{HOA}}|\) and \(k\).
Finding a subset-minimal valid cause is therefore a search over an exponential
space; validity is not in general monotone in \(\Gamma\)
(Appendix~\ref{app:tempobench-alignment}), so greedy deletion is not sound without
an extra monotonicity hypothesis, and the minimal cause is computed by exhaustive
search, hitting-set, or CEGIS over the finite constraints. Each individual check
remains polynomial.

\section{Benchmark Architecture}

The construction naturally suggests benchmark tasks for temporal causal
reasoning:
\begin{enumerate}[label=\arabic*.]
  \item choose an LTL specification \(\varphi\);
  \item synthesize or hand-design a Mealy machine \(M\);
  \item choose an actual lasso trace \(\tau\);
  \item choose an effect \(E\), usually a focused subproperty of \(\varphi\);
  \item compute or manually design the intended minimal cause \(C\);
  \item ask a solver to identify \(C\) from plausible distractors.
\end{enumerate}

A task could ask: given this Mealy machine and the effect \(F\,grant\), which
temporal fact is the minimal cause? Candidate answers might include
\[
  req,\qquad
  F\,req,\qquad
  req\land X\,grant,\qquad
  G(req\rightarrow F\,grant),\qquad
  GF\,grant .
\]
The correct answer depends on whether the task asks for a trace-level cause or
a strategy-level cause. Making that distinction explicit is part of the point
of the benchmark.

\section{Main Pitfall}

The phrase ``minimal cause within the underlying LTL specification'' can
mean several different things:
\begin{enumerate}[label=\arabic*.]
  \item a minimal syntactic subformula of \(\varphi\);
  \item a minimal trace property that explains \(E\) on \(\tau\);
  \item a minimal strategy invariant of \(M\) that explains realization of
        \(\varphi\).
\end{enumerate}
These are not equivalent. For instance, if
\[
  \varphi =
  G(req\rightarrow F\,grant)
  \land
  G(cancel\rightarrow X\,\neg grant)
\]
and the observed effect is \(F\,grant\), then the relevant cause may depend on
the absence of \(cancel\), the current machine state, and the recurrence
structure of the product automaton. It may not be a literal subformula of
\(\varphi\).

For this reason, the cause should not be restricted to syntactic subformulas of
the original specification. A better approach is to let \(C\) be any
\(\omega\)-regular property in the chosen candidate class, then translate or
approximate it as readable LTL when needed.

\section{Certificate View}
\label{sec:certificate}

A useful output is not only a formula \(C\), but a certificate
\[
  (C,W,\Pi).
\]
Here \(C\) is the temporal cause and \(W\) is the accepting or winning region
that the cause preserves: in the one-player regimes \(W\) is the set of product
states of \(P=M\times A_E\) from which the effect is still enforced (the accepting
region witnessing \(\tau\models E\)); in the strategy-structure regime \(W\) is
the \(\exists\)-winning region of \(G^{\mathrm{rev}}\) (Section~\ref{sec:rpg}). We
define \(\Pi\) precisely as a finite set of \emph{counterfactual witnesses}, one
per rejected weaker (or otherwise competing) candidate \(C'\): each element of
\(\Pi\) is a lasso trace
\[
  \pi_{C'}=u_{C'}\,v_{C'}^{\omega}\ \in\
  \Close_{\tau,M}^{\Frame}(\neg C')\cap\Lang(E_{\Frame}),
\]
that is, a closest admissible \(\neg C'\)-removal that \emph{preserves} the
effect, which is exactly the counterexample showing \(C'\notin\mathsf{Causes}\).
Such a lasso exists whenever \(C'\) fails the test and is extracted, in the
one-player regimes, from the accepting cycle found by the non-emptiness check of
Step~\ref{step:check}, and in the strategy-structure regime from the
\(\exists\)-winning play of Proposition~\ref{prop:rpg-sound}. Thus \(\Pi\) is a
finite, checkable record of why no preferred competitor is a valid cause.

This certificate separates three ideas:
\begin{itemize}
  \item what happened on the observed execution,
  \item what had to keep happening over time,
  \item what would break if the proposed cause were removed.
\end{itemize}

The concise form of the construction, stated precisely, is the following.
\begin{center}
\fbox{\begin{minipage}{0.9\linewidth}
\textbf{Arena.} Form the frozen product \(P=M\times A_{E_{\Frame}}\), treating
\(M\) as a fixed strategy. For input-only and output-intervention frames this is
the one-player arena used throughout; for a strategy-structure frame the arena is
instead the original game \(G_\varphi\) with \(\sigma\) partially re-opened
(Step~\ref{step:search}).
\end{minipage}}
\end{center}
\begin{center}
\fbox{\begin{minipage}{0.9\linewidth}
\textbf{Counterfactual test.} \emph{Trace-existential necessity} (N1, the core
\(\mathsf{Causes}\) definition): \(C\) passes when the closest admissible
\(\neg C\)-traces are non-empty and all fail \(E_{\Frame}\)
(Section~\ref{sec:formal}); the adversary searches for an admissible nearest
\(\neg C\)-trace that preserves \(E_{\Frame}\), and finding none is a pass. This
is a one-player shortest-edit search plus non-emptiness check in \emph{all} regimes
(the system has no free move once \(M\) is frozen). \emph{Strategy-guarantee necessity} (N2,
strategy-level): the genuinely two-player reverse parity game of
Section~\ref{sec:rpg} asks whether the re-opened system can drop \(C\) and still
force \(E\) against every environment. The two-player parity machinery is the
object of N2, not of N1.
\end{minipage}}
\end{center}
\begin{center}
\fbox{\begin{minipage}{0.9\linewidth}
\textbf{Selection.} A minimal cause is any \(\triangleleft\)-minimal valid cause;
when \(\triangleleft\) is partial, several incomparable minimal causes may exist
(Definition~\ref{def:mincause}). Separately, the nearest-failure-anchored cause \(C^{\flat}\)
(Proposition~\ref{prop:weakest}) is a canonical \emph{certificate} --- the
maximal failure-free neighborhood of \(\tau\) --- not generally a readable
\(\triangleleft\)-minimal formula (Remark~\ref{rem:weakest-degenerate}).
\end{minipage}}
\end{center}

\section{Limitations and Conclusion}
\label{sec:conclusion}

\paragraph{What is established.} We have given a frame-relative, counterfactual
notion of temporal cause for synthesized controllers, decided in the common
regimes by shortest-edit search plus automaton non-emptiness
(Steps~\ref{step:check}--\ref{step:search}) and in
the strategy-structure regime by a two-player reverse parity game
(Section~\ref{sec:rpg}). Proposition~\ref{prop:weakest} characterizes the
nearest-failure certificate \(C^{\flat}\); Proposition~\ref{prop:rpg-sound}
relates the parity game to strategy-guarantee necessity; and
Proposition~\ref{prop:tb-reduction} shows the finite construction is the literal
specialization of the general definition, aligning it with TempoBench TCE.

\paragraph{What is conjectural or scoped.} Several claims are deliberately scoped.
The complexity statements of Section~\ref{sec:complexity} are stated at the level
of construction sizes and the standard solving bounds, not as a fully proved
complexity theorem with hypotheses; they rely on the automaticity of
\(\preceq_\tau\), which we establish for the canonical order on lasso
counterfactuals (Remark~\ref{rem:lasso-close}) but assume in the general
\(\omega\)-regular case. The reverse parity game decides \emph{strategy-guarantee}
necessity (N2), which is a distinct, type-level notion from the trace-existential
necessity (N1) of the core definition (Section~\ref{sec:two-notions}); we do not
claim a single game decides N1, nor that a single game also performs the
edited-vertex minimization. The relationship to Halpern--Pearl is a contrast, not
a proved separation (Section~\ref{sec:overdetermination}).

\paragraph{Limitations.} The method reports frame-relative \emph{necessary
conditions}, not full actual causes: it returns disjunctions for overdetermined
effects, does not distinguish preempting from preempted causes, and returns no
cause when the necessary condition lies outside the declared
\(\mathcal{K}_{\Frame}\). Its outputs are determinate only after the eligible
vocabulary is declared; vocabulary choice can still influence the answer
(Section~\ref{sec:canonical-frames}). The evidence here is two small worked
examples plus the TempoBench alignment, now supplemented at the trace level: a
full end-to-end TempoBench TCE instance exists (the publicly available sample,
a six-state synthesized arbiter, whose computed \(\Gamma^{\star}\) matches
TempoBench's per-cell ground truth exactly), alongside machine-checked
reproductions of five worked trace-level instances from HOA encodings. What
remains before broader empirical claims is a corpus-scale comparison (pending
the dataset's re-release) and the strategy-level case study --- ideally
extracting a strategy-level cause from a realistic synthesized arbiter. The
trace-level checker is additionally profiled (sub-second full-key computation
on small windows; the actual-causation object is the measured exponential
mode), and on a 107-instance random-transducer corpus the per-cell
actual-cause and determining-set unions coincided on every instance.

\paragraph{Conclusion.} The core idea --- frame-relative counterfactual necessity
for temporal effects over a frozen synthesized strategy, with a parity-game
generalization for strategy-level causes --- is, we believe, a useful organizing
frame for temporal causal explanation. The present document establishes the
definitions, the two regimes, and their soundness relationships; closing the
remaining gaps (a proved complexity theorem, the edited-vertex minimization
procedure, and an empirical evaluation) would turn this framework into a complete
method.

\appendix

\section{Symbol and Operation Dictionary}
\label{app:symbols}

This appendix lists the mathematical objects used in the document. When a
symbol has several common conventions in the literature, the convention used in
this document is stated explicitly.

\subsection*{Logical and Temporal Notation}

\begin{description}[leftmargin=2.8cm,style=nextline]
  \item[\(\LTL\)]
  Linear Temporal Logic. It is a logic for writing properties of infinite
  traces.

  \item[\(\varphi,\varphi_{\LTL}\)]
  The original LTL formula or specification. Its form is a formula over
  \(I\cup O\), so \(\varphi\in\LTL(I\cup O)\), with semantics
  \(\Lang(\varphi)\subseteq(2^{I\cup O})^\omega\). It is not the
  reverse-causal frame.

  \item[\(\Frame\)]
  The reverse-causal frame. It records the scoping information used by the
  explanation procedure, including the selected effect, background constraints,
  relevant outputs, candidate-cause space, allowed interventions, similarity
  relation, and minimality order. This is an explanatory scope object and is
  unrelated to a Kripke frame in modal logic; the term ``frame'' is used here in
  the causal-context sense.

  \item[\(\Frame_\varphi\)]
  A frame induced by \(\varphi\) and additional design choices \(\Design\).
  This notation means \(\Frame_\varphi=\operatorname{Frame}(\varphi;\Design)\),
  not that \(\varphi\) itself is the frame.

  \item[\(\Design\)]
  Additional design or scoping choices used to construct a frame, such as target
  selection, allowed counterfactuals, cause grammar, similarity metric, and
  minimality order.

  \item[\(\Frame_{\mathrm{scope}}\)]
  A scope-setting causal frame. It tells the reverse procedure which part of a
  larger effect should be treated as the target, which parts should be treated
  as background, and which parts may be ignored.

  \item[\(\Frame_{\mathrm{TB}}\)]
  A TempoBench-compatible causal frame. It restricts the effect to a positive
  output occurrence at one fixed timestep.

  \item[\(\AP(\varphi)\)]
  The set of atomic propositions mentioned by \(\varphi\).

  \item[AP]
  Atomic proposition. In TempoBench prompts, the AP list is the finite list of
  Boolean symbols used by the automaton, including both inputs and outputs.

  \item[\(E\)]
  The effect to be explained: a \emph{single} (possibly conjunctive)
  \(\omega\)-regular trace property. Its form is usually an LTL formula, a bounded
  temporal formula, or an equivalent language of traces. When convenient, this
  document uses \(E\) both for the formula and for the set of traces satisfying
  it; \(\Lang(E)\) is the explicit language notation. The \emph{space} of
  candidate effects one might choose is a separate object, written \(\Effects\);
  \(E\in\Effects\).

  \item[\(\Effects\)]
  The space (class) of candidate effects from which a particular effect \(E\) is
  drawn. It is a set of trace properties, not a single effect; the
  TempoBench-compatible effect class \(\mathcal{E}_{TB}(\Frame_{\mathrm{TB}})\) is
  one such space.

  \item[\(E_1,\ldots,E_m\)]
  Components of a larger effect \(E\). For example, if
  \(E=E_1\land E_2\land E_3\), then each \(E_r\) is one conjunctive obligation
  or sub-effect.

  \item[\(E_{\Frame}\)]
  The part of \(E\) selected by \(\Frame\) as the target effect. If no
  selection is made, then \(E_{\Frame}=E\).

  \item[\(B_{\Frame}\)]
  Background constraints supplied by \(\Frame\). Counterfactual traces may be
  required to satisfy \(B_{\Frame}\), so those constraints are treated as the
  setting of the causal question rather than as part of the cause.

  \item[\(S_{\Frame}\)]
  The set of selected effect indices. If \(E=E_1\land\cdots\land E_m\), then
  \(S_{\Frame}\subseteq\{1,\ldots,m\}\) identifies which \(E_r\)'s form the
  scoped target \(E_{\Frame}\).

  \item[\(O_{\Frame}\)]
  The output propositions selected as relevant by \(\Frame\). Its form is a
  subset \(O_{\Frame}\subseteq O\).

  \item[\(A_{\mathrm{env}}\)]
  Environment assumptions. These are constraints on admissible environment
  behavior, usually written as an LTL formula.

  \item[\(G_{\mathrm{sys}}\)]
  System guarantees. These are obligations the system is intended to satisfy,
  usually written as an LTL formula.

  \item[\(C\)]
  A candidate temporal cause. Its form is usually an LTL formula, a parity or
  B\"uchi automaton, or any representation of an \(\omega\)-regular trace
  property.

  \item[\(\mathcal{K}_{\Frame}\)]
  The candidate-cause space allowed by \(\Frame\). Its elements are possible
  explanations before they have been checked for validity or minimality.

  \item[\(\mathcal{I}_{\Frame}\)]
  The allowed intervention model supplied by \(\Frame\). It determines which
  counterfactual changes may be made to inputs, outputs, transitions,
  priorities, or strategy commitments. Note the deliberate typographic
  distinction from the input set \(I\): \(\mathcal{I}_{\Frame}\) (calligraphic,
  frame-subscripted) is the intervention model, whereas \(I\) (upright) is the
  set of input atomic propositions; they are unrelated objects.

  \item[\(\mathsf{Causes}_{M,\Frame,\tau}(E_{\Frame})\)]
  The set of all valid causes of selected effect \(E_{\Frame}\), relative to
  Mealy machine \(M\), causal frame \(\Frame\), and actual trace \(\tau\).
  Membership in this set means that actuality and counterfactual dependence
  hold.

  \item[\(C_{\mathrm{valid}}\)]
  An arbitrary valid cause. Its form is an element of
  \(\mathsf{Causes}_{M,\Frame,\tau}(E_{\Frame})\).

  \item[\(C^\star\)]
  A minimal cause (one of possibly several incomparable ones). Its form is a
  valid cause selected from \(\mathsf{Causes}_{M,\Frame,\tau}(E_{\Frame})\) such
  that no preferred valid cause is available under the ordering \(\triangleleft\).

  \item[\(C_\tau\)]
  A trace-level cause for the actual trace \(\tau\). It explains one observed
  execution.

  \item[\(C_M\)]
  A strategy-level cause for the Mealy machine \(M\). It explains a rule or
  invariant of the controller across its executions.

  \item[\(p,a,b,g,r,req,grant,ack,cancel,pending,stable,safe,ready\)]
  Atomic propositions used in examples. Each is a Boolean-valued proposition at
  each time step. The names \(g,r,req,grant,ack,cancel,pending,stable,safe,ready\)
  are mnemonic examples, not globally fixed symbols. Symbols such as
  \(in_0,in_1,\ldots\) are example input proposition names.

  \item[\(\neg C\)]
  Negation. The formula \(\neg C\) holds exactly on traces where \(C\) does not
  hold.

  \item[\(C_1\land C_2\)]
  Conjunction. It holds on traces where both \(C_1\) and \(C_2\) hold.

  \item[\(C_1\lor C_2\)]
  Disjunction. It holds on traces where at least one of \(C_1\) or \(C_2\)
  holds. This operation is not central in the construction but is part of
  ordinary LTL.

  \item[\(C_1\rightarrow C_2\)]
  Implication. It is false only when \(C_1\) holds and \(C_2\) does not hold.

  \item[\(X C\)]
  The next-time temporal operation. It holds at time \(t\) when \(C\) holds at
  time \(t+1\).

  \item[\(X^k C\)]
  The next-time operation repeated \(k\) times. It holds at the current time
  when \(C\) holds exactly \(k\) steps later. For example, \(X^3 grant\) means
  that \(grant\) holds three steps after the current position.

  \item[\(F C\)]
  The eventually temporal operation. It holds at time \(t\) when \(C\) holds at
  some time \(t'\geq t\).

  \item[\(G C\)]
  The always temporal operation. It holds at time \(t\) when \(C\) holds at
  every time \(t'\geq t\).

  \item[\(C_1\,U\,C_2\)]
  The until temporal operation. It holds when \(C_2\) eventually holds and
  \(C_1\) holds at every earlier time until then.

  \item[\(GF\,C\)]
  Infinitely often \(C\). This abbreviates \(G(F C)\). It means that from every
  point in time, \(C\) occurs again later.

  \item[\(FG\,C\)]
  Eventually always \(C\). This abbreviates \(F(G C)\). It means that from some
  future point onward, \(C\) always holds.

  \item[\(\models\)]
  Satisfaction. The expression \(\tau\models C\) means that trace \(\tau\)
  satisfies the temporal property \(C\).

  \item[\(\ell\)]
  An input literal. Its form is either \(i\) or \(\neg i\), where \(i\in I\) is
  one input atomic proposition.

  \item[\(c_j\)]
  A conjunction of input literals assigned to timestep \(j\). For example,
  \(c_j\) could be \(in_1\land \neg in_3\).

  \item[\(C_E\)]
  A finite-horizon cause for effect \(E\). In the TempoBench-compatible setting
  below, it has the form
  \[
    C_E=\bigwedge_{j=0}^{k} X^j c_j .
  \]
\end{description}

\subsection*{Sets, Traces, and Languages}

\begin{description}[leftmargin=2.8cm,style=nextline]
  \item[\(I\)]
  The set of input atomic propositions. Each element is controlled by the
  environment.

  \item[\(O\)]
  The set of output atomic propositions. Each element is controlled by the
  system.

  \item[\(I/O\text{ partition}\)]
  The division of atomic propositions into input propositions \(I\) and output
  propositions \(O\). It determines which player controls each proposition.

  \item[\(I\cup O\)]
  The set of all visible atomic propositions.

  \item[\(2^I\)]
  The powerset of \(I\). Its elements are sets of input propositions that are
  true at a single time step.

  \item[\(2^O\)]
  The powerset of \(O\). Its elements are sets of output propositions that are
  true at a single time step.

  \item[\(2^{I\cup O}\)]
  The alphabet of complete input/output valuations. Each letter is a set of
  propositions true at one time step.

  \item[\(i_t\)]
  The input valuation at time \(t\). Its form is a set \(i_t\in 2^I\).

  \item[\(o_t\)]
  The output valuation at time \(t\). Its form is a set \(o_t\in 2^O\).

  \item[\(a_t\)]
  The complete valuation at time \(t\), defined by \(a_t=i_t\cup o_t\). Its form
  is a set \(a_t\in 2^{I\cup O}\).

  \item[\(\alpha_t\)]
  A complete valuation at time \(t\) in a finite observed trace. It has the same
  form as \(a_t\): a set \(\alpha_t\subseteq I\cup O\).

  \item[\(t\)]
  A time index. Its form is a natural number \(t\in\Nat\).

  \item[\(j,k,n,m,r\)]
  Indices. The index \(k\) is usually the timestep at which the target output
  effect occurs, \(j\) ranges over earlier or equal timesteps, \(n\) is the last
  index of a finite trace, \(m\) is the number of components in a decomposed
  effect, and \(r\) indexes one selected component \(E_r\).

  \item[\(\Nat\)]
  The set of natural numbers.

  \item[\(\tau\)]
  The actual trace being explained. Its form is an infinite word
  \(a_0a_1a_2\cdots\) over alphabet \(2^{I\cup O}\).

  \item[\(\tau_{\le n}\)]
  A finite observed trace prefix. Its form is
  \(\alpha_0\alpha_1\cdots\alpha_n\), where each
  \(\alpha_t\subseteq I\cup O\).

  \item[\(\tau_{\le k}\)]
  The finite prefix of \(\tau_{\le n}\) ending at timestep \(k\), also written
  \(\tau_{\le n}[0..k]\). It has the form \(\alpha_0\alpha_1\cdots\alpha_k\).

  \item[\(|\tau_{\le n}|\)]
  The length of the finite trace prefix \(\tau_{\le n}\). If
  \(\tau_{\le n}=\alpha_0\cdots\alpha_n\), then
  \(|\tau_{\le n}|=n+1\).

  \item[\(\alpha_k\)]
  The valuation at timestep \(k\) of the finite observed trace
  \(\tau_{\le n}=\alpha_0\alpha_1\cdots\alpha_n\), using ordinary sequence
  indexing. We write \(\alpha_k\) (equivalently the bracket form
  \(\tau_{\le n}[k]\)) rather than the heavier \((\tau_{\le n})_k\).

  \item[\(\tau'\), \(\tau''\)]
  Counterfactual or comparison traces. Their form is also an infinite word over
  \(2^{I\cup O}\).

  \item[\(uv^\omega\)]
  A lasso representation of an infinite trace. Here \(u\) and \(v\) are finite
  words, and \(v^\omega\) means that \(v\) repeats forever.

  \item[\(\omega\)]
  The first infinite ordinal. In this document it is used informally to denote
  infinite repetition or infinite-word behavior.

  \item[\(\omega\)-regular property]
  A property of infinite traces recognizable by automata such as B\"uchi,
  Rabin, Streett, or parity automata.

  \item[\(\Lang(M)\)]
  The language generated by Mealy machine \(M\). Its form is a set of infinite
  traces over \(2^{I\cup O}\).

  \item[\(\Lang(E)\)]
  The language of traces satisfying effect \(E\). Its form is a set of infinite
  traces.

  \item[\(\Lang(C)\)]
  The language of traces satisfying candidate cause \(C\). Its form is a set of
  infinite traces.

  \item[\(\emptyset\)]
  The empty set. In an example trace, it means that no listed proposition is
  true at that time step.

  \item[\(\subseteq\)]
  Subset relation. \(A\subseteq B\) means every element of \(A\) is also an
  element of \(B\).

  \item[\(\subsetneq\)]
  Strict subset relation. \(A\subsetneq B\) means \(A\subseteq B\) and
  \(A\neq B\).

  \item[\(\cap\)]
  Set intersection. \(A\cap B\) contains exactly the elements that are in both
  \(A\) and \(B\).

  \item[\(\cup\)]
  Set union. \(A\cup B\) contains the elements that are in \(A\), in \(B\), or
  in both.

  \item[\(\overline{X}\)]
  Complement of a set or language \(X\), relative to the ambient universe of
  traces under discussion.
\end{description}

\subsection*{Mealy Machines}

\begin{description}[leftmargin=2.8cm,style=nextline]
  \item[\(M\)]
  A Mealy machine. Its form is the tuple
  \(M=(S,s_0,I,O,\delta,\lambda)\).

  \item[\(S\)]
  The finite set of Mealy-machine states.

  \item[\(s_0\)]
  The initial Mealy-machine state. Its form is an element \(s_0\in S\).

  \item[\(s_t\)]
  The Mealy-machine state at time \(t\). Its form is an element \(s_t\in S\).

  \item[\(\delta\)]
  The Mealy-machine transition function. Its form is
  \[
    \delta:S\times 2^I\to S.
  \]
  It maps the current state and current input valuation to the next state.

  \item[\(\lambda\)]
  The Mealy-machine output function. Its form is
  \[
    \lambda:S\times 2^I\to 2^O.
  \]
  It maps the current state and current input valuation to the current output
  valuation.

  \item[\(idle,pending\)]
  Example state names. In the toy example, they are elements of \(S\).

  \item[\(state=s\)]
  A notation for saying that the Mealy machine is currently in internal state
  \(s\). This is not automatically visible in ordinary LTL over \(I\cup O\)
  unless the state is exposed as an atomic proposition.

  \item[\(\mathit{at\_s}\)]
  An atomic proposition used to expose the internal state \(s\) to an LTL
  formula. For example, \(\mathit{at\_s5}\) means that the machine is in state
  \(s_5\). We use the prefix \(\mathit{at\_}\) rather than \(\mathit{in\_}\) to
  avoid collision with input proposition names such as \(in_0,in_1,\ldots\).
\end{description}

\subsection*{Automata and Parity Games}

\begin{description}[leftmargin=2.8cm,style=nextline]
  \item[\(A_\varphi\)]
  An automaton recognizing the traces that satisfy specification \(\varphi\).

  \item[\(A_E\)]
  A deterministic parity automaton for effect \(E\). Its form is the tuple
  \[
    A_E=(Q,q_0,2^{I\cup O},\delta_E,\Omega).
  \]

  \item[\(A_{\mathrm{HOA}}\)]
  The automaton supplied in HOA format in a TempoBench-style instance. It is an
  \emph{acceptor} over complete valuations \(2^{I\cup O}\) (inputs and outputs
  together), whose language is taken to be the set of admissible system
  behaviors; it thus plays the role of \(\Lang(M)\), not of the effect monitor
  \(A_E\), and is not a Mealy transducer. See Remark~\ref{rem:hoa} for the
  faithfulness assumption.

  \item[HOA]
  Hanoi Omega-Automata format. In this document it means the textual automaton
  representation used by TempoBench prompts.

  \item[\(Q\)]
  The finite set of automaton states.

  \item[\(q_0\)]
  The initial automaton state. Its form is an element \(q_0\in Q\).

  \item[\(q_t\)]
  The automaton state at time \(t\). Its form is an element \(q_t\in Q\).

  \item[\(q_t\xrightarrow{\alpha_t}q_{t+1}\)]
  A transition from state \(q_t\) to state \(q_{t+1}\) under valuation
  \(\alpha_t\).

  \item[\(\delta_E\)]
  The transition function of the automaton for \(E\). Its form is
  \[
    \delta_E:Q\times 2^{I\cup O}\to Q.
  \]
  It updates the automaton state after reading one complete valuation.

  \item[\(\Omega\)]
  The priority function of a parity automaton or parity game. In this document
  its form is a map from states or vertices to natural numbers. The symbol
  \(\Omega\) is reused for the game \(G\), the effect automaton \(A_E\), and the
  synthesis game \(G_\varphi\); in all of them the \emph{same} convention is in
  force --- a run/play is winning iff the least priority seen infinitely often is
  even (Sections~\ref{sec:background-parity} and~\ref{step:freeze}).

  \item[\(G_\varphi\)]
  The parity game built from specification \(\varphi\). It encodes the
  environment choices, system choices, and automaton progress needed to decide
  whether the system can force \(\varphi\).

  \item[\(\sigma\)]
  A system strategy. It maps the information available to the system to an
  output choice or game move. In finite-memory form it can be represented by a
  Mealy machine.

  \item[\(\rho\)]
  An environment behavior. Its form is an infinite sequence of environment
  choices, usually input valuations over \(I\).

  \item[\(\mathrm{Play}(\sigma,\rho)\)]
  The infinite play induced by system strategy \(\sigma\) and environment
  behavior \(\rho\). It is the trace or game path that results when the system
  follows \(\sigma\) while the environment follows \(\rho\).

  \item[\(P=M\times A_E\)]
  The synchronous product of the Mealy machine and the automaton for \(E\).
  Its states have the form \((s,q)\in S\times Q\).

  \item[\((s,q)\)]
  A product state, consisting of a Mealy-machine state \(s\) and an automaton
  state \(q\).

  \item[\(V\)]
  A set of vertices in a parity game.

  \item[\(T\)]
  A set of transitions or edges in a game graph.

  \item[\(D_C\)]
  A causal dominion or causal region associated with cause \(C\). Its form is a
  set of game vertices, \(D_C\subseteq V\).

  \item[\(T_C\)]
  A set of causal strategy commitments. Its form is a set of transitions,
  \(T_C\subseteq T\).

  \item[\(W\)]
  A winning or accepting region preserved by a cause. Its form depends on the
  representation, but usually it is a subset of product states or game vertices.

  \item[\(\mathrm{Acc}_{\mathrm{parity}}\)]
  The parity acceptance condition. It is the property that the infinite play's
  recurring priority pattern is winning under the parity convention chosen for
  the game or automaton.

  \item[SCC]
  A strongly connected component of a directed graph. It is a set of vertices
  where every vertex in the set can reach every other vertex in the set. In
  parity reasoning, an accepting SCC is a recurring region whose infinite
  behavior satisfies the parity acceptance condition.
\end{description}

\subsection*{Counterfactual and Minimality Notation}

\begin{description}[leftmargin=2.8cm,style=nextline]
  \item[\(\preceq_\tau\)]
  A similarity order relative to actual trace \(\tau\). The expression
  \(\tau'\preceq_\tau\tau''\) means that \(\tau'\) is at least as close to
  \(\tau\) as \(\tau''\) is.   The canonical instance used in this document is the count-first
  (minimal-intervention) \emph{lexicographic total order} of
  Step~\ref{step:similarity}, which is antisymmetric with \(\tau\) as unique
  minimum; the first-difference/Baire order and the pure change-set inclusion
  order are coarser alternatives not used as primary.

  \item[\(\prec_\tau\)]
  The strict part of \(\preceq_\tau\). It means strictly closer to \(\tau\).

  \item[\(d_I(\tau,\tau')\)]
  An input-distance measure between traces. In the finite-prefix case used as
  an example, it counts the number of positions where the input valuations of
  \(\tau\) and \(\tau'\) differ.

  \item[\(\Close_{\tau,M}(\neg C)\)]
  Shorthand for the input-only special case of \(\Close_{\tau,M}^{\Frame}\),
  obtained by taking \(B_{\Frame}=\True\) and the input-only intervention model,
  so that \(\Adm_M^{\Frame}(\tau)=\Lang(M)\). It is the set of
  \(\preceq_\tau\)-minimal elements of
  \(\{\tau'\in\Lang(M)\mid \tau'\models\neg C\}\). The frame-relative form below
  is the primary definition.

  \item[\(\Close_{\tau,M}^{\Frame}(\neg C)\)]
  The set of closest admissible counterfactual traces allowed by \(\Frame\) that
  violate \(C\). Formally, it is the set of \(\preceq_\tau\)-minimal elements of
  \[
    \{\tau'\in\Adm_M^{\Frame}(\tau)\mid \tau'\models\neg C\}.
  \]

  \item[\(\Adm_M^{\Frame}(\tau)\)]
  The admissible counterfactual traces under the frame. In the ordinary
  input-only case,
  \[
    \Adm_M^{\Frame}(\tau)=\Lang(M)\cap\Lang(B_{\Frame}).
  \]

  \item[\(\min_{\preceq_\tau} X\)]
  The set of minimal elements of \(X\) under preorder \(\preceq_\tau\). An
  element \(x\in X\) is minimal when no element of \(X\) is strictly closer
  according to the preorder.

  \item[\(\uparrow_{\preceq_\tau}(X)\)]
  The upward closure of set \(X\) under preorder \(\preceq_\tau\):
  \(\{\tau''\mid \exists\,\tau'\in X.\ \tau'\preceq_\tau\tau''\}\). Because
  \(\tau'\preceq_\tau\tau''\) means \(\tau'\) is \emph{closer} to the actual
  trace \(\tau\), ``above'' here means \emph{further away from} \(\tau\), not
  closer or better. Thus \(\uparrow_{\preceq_\tau}(\neg E)\) is the set of traces
  that are further from \(\tau\) than some effect-violating trace.

  \item[\(\triangleleft\)]
  A strict preference order on candidate causes. \(C'\triangleleft C\) means that
  \(C'\) is preferred to \(C\). The canonical instance
  (Section~\ref{sec:formal}) is lexicographic: smaller formula size, then smaller
  temporal depth, then more permissive (larger \(\Lang(C)\)).

  \item[\(\Pi\)]
  A finite set of counterfactual witnesses (Section~\ref{sec:certificate}). For
  each rejected competing candidate \(C'\), \(\Pi\) contains a lasso trace in
  \(\Close_{\tau,M}^{\Frame}(\neg C')\cap\Lang(E_{\Frame})\) --- a closest
  admissible \(\neg C'\)-removal that preserves the effect --- witnessing that
  \(C'\) is not a valid cause.

  \item[\((C,W,\Pi)\)]
  A certificate of explanation. \(C\) is the cause, \(W\) is the winning or
  accepting region it preserves, and \(\Pi\) records counterfactual evidence
  supporting minimality or rejecting alternatives.
\end{description}

\subsection*{TempoBench-Compatible Notation}

\begin{description}[leftmargin=2.8cm,style=nextline]
  \item[TTE]
  Temporal trace evaluation. This TempoBench task asks whether a finite trace is
  accepted by a supplied automaton.

  \item[TCE]
  Temporal causal evaluation. This TempoBench task asks for the minimal
  timestep-indexed input constraints needed for a selected output effect.

  \item[TempoBench-style TCE]
  Temporal causal evaluation in the restricted benchmark sense used in
  Appendix~\ref{app:effect}. The input is a finite-state machine, a finite trace generated by
  it, and a target output selected by \(\Frame_{\mathrm{TB}}\) that occurs at a particular
  timestep. The output is a set of timestep-indexed input conditions needed for
  that target output.

  \item[\(\texttt{XXX }g\)]
  TempoBench prompt notation for \(X^3 g\). More generally, \(k\) copies of
  \(\texttt{X}\) before an output name denote \(X^k\) applied to that output.

  \item[AP-level score]
  A TempoBench causality score that gives partial credit for correctly predicted
  atomic propositions at a timestep.

  \item[TS-level score]
  A TempoBench causality score that gives credit for a timestep only when the
  full predicted set of constraints at that timestep matches the ground truth.

  \item[\(\mathcal{E}_{TB}(\Frame_{\mathrm{TB}})\)]
  The TempoBench-compatible effect class after \(\Frame_{\mathrm{TB}}\) has
  selected the relevant outputs. In this document it is
  \[
    \mathcal{E}_{TB}(\Frame_{\mathrm{TB}})
    =
    \{\,X^k o \mid
    o\in O_{\Frame},\ 0\leq k<|\tau_{\le n}|,\ 
    o\in\alpha_k\,\}.
  \]
  Its elements are positive selected-output occurrences at fixed timesteps.

  \item[\(o\)]
  A single output atomic proposition, with \(o\in O\). This differs from
  \(o_t\), which is the whole set of outputs true at time \(t\). Symbols such
  as \(o_1\) and \(o_2\) are example output proposition names.

  \item[\(X^k o\)]
  A point output effect. It means that output proposition \(o\) is true exactly
  \(k\) steps after the start of the finite trace.

  \item[\(\Gamma_E\)]
  A set representation of a finite-horizon cause for effect \(E\). Its elements
  are pairs \((t,\ell)\), where \(t\) is a timestep and \(\ell\) is an input
  literal required at that timestep.

  \item[\(\Gamma_E^\star\)]
  The minimal finite-horizon cause set for effect \(E\). Its elements are
  the timestep-indexed input literals selected as necessary.

  \item[\((t,\ell)\)]
  A timestep-indexed input constraint. It means that input literal \(\ell\) must
  hold at timestep \(t\).
\end{description}

\subsection*{Common Operations in the Construction}

\begin{description}[leftmargin=2.8cm,style=nextline]
  \item[Compilation to automata]
  The operation that converts an LTL formula, such as \(E\) or \(\varphi\), into
  an automaton recognizing exactly the traces that satisfy it.

  \item[Synchronous product]
  The operation \(M\times A_E\) that runs the Mealy machine and the automaton in
  lockstep on the same input/output valuation at each time step.

  \item[Freezing the strategy]
  The operation of treating the Mealy machine \(M\) as fixed, so the system no
  longer chooses outputs freely during explanation.

  \item[Counterfactual intervention]
  An allowed change to inputs, outputs, transitions, priorities, or strategy
  commitments, depending on the chosen counterfactual model.

  \item[Cause checking]
  The operation that tests whether
  \[
    \Close_{\tau,M}^{\Frame}(\neg C)\neq\emptyset
    \quad\text{and}\quad
    \Close_{\tau,M}^{\Frame}(\neg C)\cap\Lang(E)=\emptyset .
  \]
  The first condition is non-vacuity (a closest removal exists); when it holds
  and the intersection is empty, the closest removals of \(C\) do not preserve
  the effect.

  \item[Cause synthesis]
  The operation of searching for a valid cause \(C\), rather than merely
  checking a proposed one.

  \item[Semantic minimality]
  A minimality criterion based on the trace languages denoted by candidate
  causes, rather than only the syntactic size of formulas.

  \item[Formula-size minimality]
  A minimality criterion based on the number of symbols or syntax-tree nodes in
  a candidate LTL formula.

  \item[Automaton-size minimality]
  A minimality criterion based on the number of states, transitions, or
  priorities in an automaton representing the candidate cause.

  \item[Temporal-depth minimality]
  A minimality criterion based on the nesting depth of temporal operators such
  as \(X\), \(F\), \(G\), and \(U\).
\end{description}

\section{The Original Specification \texorpdfstring{\(\varphi\)}{phi} and the Reverse Causal Frame \texorpdfstring{\(\Frame\)}{F}}
\label{app:phi}

In ordinary synthesis, \(\varphi\) is an LTL formula over \(I\cup O\):
\[
  \varphi\in\LTL(I\cup O),
  \qquad
  \Lang(\varphi)\subseteq(2^{I\cup O})^\omega.
\]
The reverse-causal setting does not change the type of \(\varphi\). Instead, it
introduces a new object, the causal frame \(\Frame\). The frame may be induced
by an original specification plus additional choices:
\[
  \Frame=\operatorname{Frame}(\varphi;\Design).
\]
If no original \(\varphi\) is available, \(\Frame\) can be specified directly.

The Mealy machine \(M\) tells us what behavior is produced: for each state and
input valuation, it tells us which output valuation appears next. By itself,
however, the machine does not say which behavior is intended, which assumptions
define the intended environment, which outputs matter, or which facts should
count as explanatory rather than accidental. The role of \(\Frame\) is to
provide that missing reverse-causal frame. A useful slogan is:
\[
  M \text{ gives behavior;}
  \qquad
  \Frame \text{ gives explanatory scope.}
\]
Thus the reverse-causal answer should be read as:
\[
  \text{cause of } E
  \text{ among behaviors of } M
  \text{ relative to causal frame } \Frame.
\]
If no original \(\varphi\) is supplied, the answer can still be frame-relative:
the frame is then chosen directly rather than induced from a specification.

\subsection{What \texorpdfstring{\(\Frame\)}{F} Can Represent}

In ordinary synthesis, \(\varphi\) is often the original LTL specification. In
the reverse-causal setting, \(\Frame\) records the additional information needed
to ask an explanatory question:
\begin{enumerate}[label=\arabic*.]
  \item \textbf{Original specification.} It can state the temporal property the
        controller was designed to satisfy, when such a \(\varphi\) is
        available.
  \item \textbf{Environment assumptions.} It can say which input behaviors are
        part of the intended operating environment.
  \item \textbf{System guarantees.} It can say which output behaviors the
        system is expected to provide under those assumptions.
  \item \textbf{Design intent.} It can identify which parts of the machine
        behavior are meaningful for the explanation.
  \item \textbf{Effect scope.} It can select which part of a larger effect is
        the target of explanation.
  \item \textbf{Vocabulary restriction.} It can restrict explanations to
        specification-level facts, rather than arbitrary implementation details.
\end{enumerate}

This means \(\Frame\) does not have to be treated as something already fully
encoded inside \(M\), and it is not the same object as \(\varphi\). It may
include assumptions, relevance choices, intervention policies, or abstractions
that are external to the machine. The important requirement is honesty: if a
frame is used, the result must be reported as relative to \(\Frame\).

\subsection{Restricting Admissible Counterfactuals}

The first major use of \(\Frame\) is to restrict which counterfactual traces
are allowed. Without background constraints, the reverse game may consider every trace
generated by \(M\):
\[
  \tau' \in \Lang(M).
\]
With background constraints \(B_{\Frame}\), the reverse game can instead require
\[
  \tau' \in \Lang(M)\cap\Lang(B_{\Frame}).
\]
This says that counterfactual traces must still be machine behaviors, but they
must also respect the intended background assumptions.

For example, suppose the machine has inputs \(req\) and \(cancel\), and output
\(grant\). Suppose the machine grants exactly when a request occurs and no
cancel occurs in the same step. If the actual step has \(req\), no \(cancel\),
and \(grant\), then a machine-relative cause of \(grant\) is
\[
  req\land \neg cancel.
\]
Both parts matter if arbitrary inputs are allowed.

Now suppose the background specification includes the environment assumption
\[
  A_{\mathrm{env}} = G(req\rightarrow \neg cancel).
\]
This says that the intended environment never cancels at the same time as a
request. If the reverse question is evaluated relative to this assumption, then
the counterfactual where \(req\) and \(cancel\) happen together is not
admissible. The minimal cause can become
\[
  req.
\]
The no-cancel fact has not vanished. It has moved from the cause into the
background setting.

This is the same distinction used in ordinary explanation. If a lamp turns on
when someone flips a switch, we usually say the switch caused the lamp to turn
on. We do not list every background condition, such as the wire being intact and
the house having electricity. In this construction, \(\Frame\) tells the
reverse game which facts are background conditions and which facts are candidate
causes.

\subsection{Selecting Which Part of \texorpdfstring{\(E\)}{E} to Explain}

The second major use of \(\Frame\) is effect selection. A large effect may
contain several obligations:
\[
  E = E_1\land E_2\land \cdots \land E_m.
\]
The reverse procedure may not always need to explain all of them. The scope
provided by \(\Frame\) can select a subset
\[
  S_{\Frame}\subseteq \{1,\ldots,m\}.
\]
The selected effect is then
\[
  E_{\Frame} =
  \bigwedge_{r\in S_{\Frame}} E_r.
\]
The reverse game computes a cause for \(E_{\Frame}\), not necessarily for all of
\(E\).

For example, suppose
\[
  E =
  G(req\rightarrow F\,grant)
  \land
  G(error\rightarrow F\,reset)
  \land
  G\neg(grant\land reset).
\]
This combines request handling, error handling, and mutual exclusion. If the
observed behavior contains a request followed by a grant, the explanation may
only need the request-grant obligation. A scope-setting specification can say:
\[
  \Frame_{\mathrm{scope}}
  =
  \text{only explain the request-grant obligation.}
\]
Then the selected target is
\[
  E_{\Frame} = G(req\rightarrow F\,grant).
\]
The reset behavior and mutual exclusion behavior need not appear in the cause
unless the frame says to keep them as background constraints.

There are two common choices:
\begin{enumerate}[label=\arabic*.]
  \item \textbf{Ignore the rest.} Explain \(E_{\Frame}\) and drop the unselected
        obligations from the reverse game.
  \item \textbf{Hold the rest fixed.} Explain \(E_{\Frame}\), but require the
        counterfactual traces to keep satisfying the unselected background
        obligations.
\end{enumerate}
The second choice is more like a controlled experiment: one target is varied
while the rest of the contract remains stable.

\subsection{Constraining the Vocabulary of Causes}

The third major use of \(\Frame\) is to restrict the language in which causes
are stated. A raw machine-level explanation might mention an internal state:
\[
  state=s_3.
\]
A specification-level explanation might instead mention a temporal fact:
\[
  req
  \qquad\text{or}\qquad
  G(req\rightarrow F\,grant).
\]
Both may be correct at different levels. The first describes implementation
state. The second describes the behavior the specification made important.

This matters when the machine contains outputs or states that are irrelevant to
the intended property. Suppose a controller both grants requests and toggles a
heartbeat output every step. If the target effect is a grant, a purely
machine-relative search may find accidental correlations involving heartbeat.
If
\[
  \varphi = G(req\rightarrow F\,grant),
\]
then the causal vocabulary can focus on \(req\), \(grant\), and related
request-response facts. The heartbeat behavior can be ignored unless
\(\Frame\) declares it relevant.

\subsection{Avoiding Vacuous Explanations}

The formula
\[
  G(req\rightarrow F\,grant)
\]
is true on a trace with no requests. It is true because no request obligation
was ever activated. This is called vacuous satisfaction. It is logically valid,
but it may not be the effect one wants to explain.

If the reverse question asks why the whole formula held on a no-request trace,
one possible answer is
\[
  G\neg req.
\]
That answer is technically meaningful: no request was left unanswered because
no request occurred. But if the benchmark is meant to test response reasoning,
this is not the intended target.

The frame \(\Frame\) can prevent this confusion by selecting only activated
obligations. If a request occurred at time \(0\), then the selected effect may
be the particular obligation that this request eventually gets a grant. If no
request occurred, the procedure can report that no non-vacuous request-grant
effect was activated.

\subsection{Modular Parity-Game Use}

The role of \(\Frame\) also makes the parity-game construction more flexible.
If a large effect decomposes as
\[
  E=E_1\land E_2\land E_3,
\]
then one can build separate monitors for the components instead of treating the
effect as one monolithic automaton:
\[
  A_E = A_{E_1}\times A_{E_2}\times A_{E_3}.
\]
The frame \(\Frame\) can then select which monitor is the target, which
monitors are background constraints, and which monitors are irrelevant for the
current explanation.

The reverse game is therefore not only checking whether a candidate cause is
necessary for a large undifferentiated property. It can ask a scoped question:
\[
  \text{Is } C \text{ necessary for the selected effect } E_{\Frame}
  \text{ under the background } B_{\Frame}?
\]
This makes the procedure more configurable because it can explain one meaningful
temporal obligation while controlling or ignoring the rest of the specification.

\subsection{Restricting \texorpdfstring{\(E\)}{E} to Match TempoBench}

The broad reverse-parity framework allows \(E\) to be any temporal effect. To
match a TempoBench-style task, \(\Frame_{\mathrm{TB}}\) restricts that broad
space to a much smaller effect class.

Let \(O_{\Frame}\subseteq O\) be the output propositions selected by the
TempoBench-compatible frame. Then the allowed effects are
\[
  E_{\Frame} \in \mathcal{E}_{TB}(\Frame_{\mathrm{TB}}),
\]
where
\[
  \mathcal{E}_{TB}(\Frame_{\mathrm{TB}})
  =
  \{\,X^k o \mid
  o\in O_{\Frame},\ 0\leq k<|\tau_{\le n}|,\ 
  o\in\alpha_k\,\}.
\]
In plain language, \(\Frame_{\mathrm{TB}}\) says:
\[
  \text{only explain positive selected-output occurrences at fixed timesteps.}
\]

This is how \(\Frame_{\mathrm{TB}}\) restricts the broad \(E\) appendix to the
TempoBench-compatible case. The unrestricted framework may explain state
reachability, liveness, recurrence, or parity acceptance. The TempoBench frame
selects only effects of the form
\[
  E_{\Frame} = X^k o.
\]
The cause is then a finite-horizon set of input constraints needed for that
output at that timestep.

\subsection{Why This Makes the Procedure More Configurable}

Adding \(\Frame\) makes the reverse procedure applicable to a wider range of
explanatory questions, for four reasons.

First, it prevents illegal or irrelevant counterfactuals from controlling the
answer. Environment assumptions can be treated as background, not as causes.

Second, it lets the procedure explain the right part of a large effect. This is
important when a specification contains many obligations but the question is
about only one of them.

Third, it separates implementation artifacts from specification-level causes.
The Mealy machine may contain states and outputs that are real but not relevant
to the causal question.

Fourth, it supports benchmark design. By changing \(\Frame\), one can create a
spectrum of tasks: machine-relative causes, specification-relative causes,
scoped temporal causes, and TempoBench-compatible point-output causes.

The final answer should therefore always state its frame:
\[
  C \text{ is a cause of } E_{\Frame}
  \text{ relative to } M \text{ and } \Frame.
\]
Without that statement, it is unclear whether the cause explains the whole
machine behavior, the whole effect, or only the scoped effect selected by the
causal frame.

\section{The Effect \texorpdfstring{\(E\)}{E} and the TempoBench-Compatible Restriction}
\label{app:effect}

This appendix explains the role of \(E\) more explicitly. The main document
uses \(E\) to mean ``the effect we want to explain,'' a single trace property.
The kind of property is deliberately broad: an effect can encode a state, an
output, a transition, a temporal pattern, or a parity-game success condition.
The space of such candidate effects --- the menu one might choose \(E\) from ---
is denoted \(\Effects\); each subsection below describes one family of members of
\(\Effects\). Appendix~\ref{app:phi} explains how \(\Frame\) narrows a chosen
\(E\) to a scoped target. When \(\Frame\) selects a target, the reverse game
explains \(E_{\Frame}\), the selected part of \(E\).

\subsection{The Broad Meaning of \texorpdfstring{\(E\)}{E}}

The cleanest general definition is:
\[
  E = \text{an effect property over the generated play or trace.}
\]
In language form, this can be written as
\[
  E \subseteq (2^{I\cup O})^\omega .
\]
This means that \(E\) is a set of infinite traces. If the actual trace belongs
to that set, then the effect occurred:
\[
  \tau \in E .
\]

We follow the convention that \(\models\) takes a formula on its right and
\(\in\) takes a set/language; sometimes \(E\) is written as a formula, such as
\(F\,grant\), in which case \(\tau\models E\) and \(\tau\in\Lang(E)\) coincide.
When the distinction matters, the language is written explicitly as
\(\Lang(E)\). If a causal frame \(\Frame\)
selects only part of \(E\), then the operative effect is \(E_{\Frame}\). If no
selection is made, then \(E_{\Frame}=E\).

\subsection{A State Can Be an Effect}

A specific reached state is a valid special case. For example, suppose the
Mealy machine eventually reaches internal state \(s_5\). We can define
\[
  E = F(state=s_5).
\]
This means that the machine eventually reaches state \(s_5\). The reverse
causal question is then:
\[
  \text{What minimum temporal condition } C
  \text{ caused the machine to reach } s_5?
\]

There is one important technical detail. Ordinary LTL specifications usually
talk about visible inputs and outputs, not hidden Mealy-machine states. If the
effect mentions an internal state, then either:
\begin{enumerate}[label=\arabic*.]
  \item the state must be exposed as an atomic proposition, such as
        \(\mathit{at\_s5}\), or
  \item the analysis must reason directly over the Mealy-machine or product-game
        states, where \(s_5\) is visible to the analysis procedure.
\end{enumerate}

So a state effect is legitimate, but it is not automatically an ordinary LTL
property over only \(I\cup O\) unless the state is made visible.

\subsection{An Output Can Be an Effect}

An effect can also be an output occurrence. For example,
\[
  E = F\,grant
\]
means that the output \(grant\) eventually occurs. A more precise finite-time
effect is
\[
  E = X^k grant,
\]
which means that \(grant\) occurs exactly \(k\) steps after the start of the
trace.

This asks:
\[
  \text{What caused this output to appear at this time?}
\]
Possible causes might be local facts about the trace, such as
\[
  req
  \qquad\text{or}\qquad
  req \land X\,grant,
\]
or a strategy-level rule such as
\[
  G(req\rightarrow F\,grant).
\]
The right kind of answer depends on the causal frame. A trace-level question
looks for what caused this output on this execution. A strategy-level question
looks for what rule in the controller makes this output behavior reliable.

\subsection{A Transition Can Be an Effect}

The effect can be a transition rather than a state. For example,
\[
  E = F\bigl((state=s_2)\land X(state=s_3)\bigr)
\]
means that the machine eventually moves from state \(s_2\) to state \(s_3\).
This is useful when the interesting causal question is:
\[
  \text{Why did the controller choose this branch rather than another branch?}
\]
Transition effects are richer than simple reachability effects because they
focus on a decision made by the controller.

\subsection{A Temporal Pattern Can Be an Effect}

The most useful effects are often temporal patterns. Suppose the original
specification is
\[
  \varphi = G(req\rightarrow F\,grant).
\]
Several different effects could be chosen:
\[
  F\,grant,
\]
meaning that a grant eventually occurred;
\[
  req\rightarrow F\,grant,
\]
meaning that a request was answered;
\[
  G(req\rightarrow F\,grant),
\]
meaning that all requests are always eventually answered; or
\[
  GF\,grant,
\]
meaning that grants occur infinitely often.

These effects are all meaningful, but they ask different causal questions. If
\(E=F\,grant\), the cause may be local. If
\(E=G(req\rightarrow F\,grant)\), the cause is probably a recurring response
rule. If \(E=GF\,grant\), the cause may involve an infinite loop or recurrence
structure.

\subsection{A Parity-Game Success Condition Can Be an Effect}

In the parity-game view, \(E\) need not be a single state or output. It can be
the success condition of an infinite play:
\[
  E = \mathrm{Acc}_{\mathrm{parity}}.
\]
This means that the play satisfies the parity acceptance condition. More
concretely, \(E\) could mean that the play eventually remains in an accepting
SCC, or that the priorities recurring infinitely often form a winning pattern.
Strictly speaking \(\mathrm{Acc}_{\mathrm{parity}}\) is a property of the
\emph{run} (the priority sequence of the product \(P=M\times A_E\)), not directly
of the \(I\cup O\) trace. We make the trace property well-defined by tying it to
the admissible traces, on which \(P\) really does have runs. Because \(A_E\) is
deterministic, each admissible trace \(\tau'\in\Lang(M)\) has a \emph{unique} run
in \(P\); define the parity-acceptance effect, relative to admissible traces, by
\[
  \Lang(E)=\{\,\tau'\in\Lang(M)\mid \text{the unique run of } \tau' \text{ in } P
  \text{ is accepting}\,\},
\]
which is equally the projection onto \(I\cup O\) of the accepting runs of \(P\)
(every run of \(P\) sits over an \(M\)-behavior, so the projection lands inside
\(\Lang(M)\)). This is an \(\omega\)-regular trace property, so \(E\) remains the
same kind of object as the bounded and liveness effects above; for traces outside
\(\Lang(M)\) there is no run and the effect is simply not evaluated, which is
harmless since the admissible set is always \(\subseteq\Lang(M)\) in the
regimes that use this effect.

This matters because parity games are about infinite behavior. Reaching one
state is not always enough. A play can temporarily visit a state that looks bad
and still win later. It can also visit a state that looks good and later lose.
The decisive fact is the infinite recurring pattern.

For parity reasoning, a natural effect is therefore:
\[
  \text{the play reaches and remains in a region whose infinite behavior is
  accepting.}
\]
This is why the broad reverse-parity construction treats \(E\) as a temporal
property, not merely as a reached state.

\subsection{Hierarchy of Possible Effects}

The candidate effects can be arranged in an informal order of \emph{increasing
richness} (this is a heuristic reading-aid, not a strict expressiveness ordering;
a state effect is not in general less expressive than an output effect):
\[
\begin{array}{c}
  \text{state effect}
  \ \rightsquigarrow\ 
  \text{output effect}
  \ \rightsquigarrow\ 
  \text{transition effect}
  \\[2pt]
  \rightsquigarrow\ 
  \text{bounded temporal effect}
  \ \rightsquigarrow\ 
  \text{unbounded liveness effect}
  \\[2pt]
  \rightsquigarrow\ 
  \text{parity acceptance effect}.
\end{array}
\]
Here \(\rightsquigarrow\) reads ``progresses to a typically richer kind,'' and
carries no claim of strict containment between consecutive entries.

Examples are:
\[
\begin{array}{lll}
  \text{state effect} & F(state=s_5) & \text{reachability causality},\\
  \text{output effect} & F\,grant & \text{output causality},\\
  \text{transition effect} &
    F\bigl((state=s_2)\land X(state=s_3)\bigr) &
    \text{decision causality},\\
  \text{response effect} &
    req\rightarrow F\,grant &
    \text{temporal dependency},\\
  \text{invariant effect} &
    G\,safe &
    \text{safety reasoning},\\
  \text{recurrence effect} &
    GF\,ready &
    \text{infinite-time reasoning},\\
  \text{parity effect} &
    \mathrm{Acc}_{\mathrm{parity}} &
    \text{game-theoretic temporal reasoning}.
\end{array}
\]

The most general lesson is:
\[
\boxed{
  E \text{ is whatever temporal fact the reverse game is being asked to
  explain.}
}
\]
A state can be that fact, but it is only one special case.

\subsection{TempoBench-Style Temporal Causality}

The broad definition of \(E\) is useful for reverse-parity explanation, but it
is too broad for direct comparison with TempoBench-style Temporal Causality
Evaluation. To make the setting TempoBench-compatible, the causal frame
\(\Frame_{\mathrm{TB}}\) restricts the selected effect \(E_{\Frame}\). In the
restricted TempoBench-style setting used here, the task is:
\begin{enumerate}[label=\arabic*.]
  \item start with a finite-state machine;
  \item observe a finite trace generated by that machine;
  \item choose an output that occurs at a particular timestep;
  \item identify the input facts at each timestep that were necessary for that
        output to occur.
\end{enumerate}

Thus the selected effect should not initially be an arbitrary temporal formula.
To match this setting, \(E_{\Frame}\) should be observable, output-based,
finite-horizon, and timestep-indexed.

\subsection{Finite Trace Notation}

The body of the document uses \(\tau\) for an infinite trace. For the
TempoBench-compatible restriction, we use a finite observed prefix:
\[
  \tau_{\le n}=\alpha_0,\alpha_1,\ldots,\alpha_n,
\]
where each
\[
  \alpha_t\subseteq I\cup O
\]
is the complete input/output valuation at timestep \(t\).

The length of this finite trace is
\[
  |\tau_{\le n}|=n+1.
\]
Using ordinary sequence indexing, the valuation at timestep \(k\) is written
\(\alpha_k\) (equivalently \(\tau_{\le n}[k]\)), and the prefix ending at
timestep \(k\) is \(\tau_{\le k}=\tau_{\le n}[0..k]=\alpha_0\cdots\alpha_k\).

For a single atomic proposition \(o\) and a valuation \(\alpha_t\subseteq I\cup O\)
(a set of propositions), satisfaction at that position is set membership: \(o\)
holds at time \(t\) iff \(o\in\alpha_t\), and the literal \(\neg o\) holds iff
\(o\notin\alpha_t\). We use membership for atoms and reserve \(\models\) for
satisfaction of a temporal formula by a trace, e.g. \(\tau\models F\,grant\).

Temporal operators on these finite prefixes are read with \(\LTLf\) (LTL over
finite traces) semantics. Equivalently, since every effect considered here is
strictly bounded by the horizon \(k<|\tau_{\le n}|\), one may evaluate the
prefix under ordinary infinite-trace LTL by appending an arbitrary infinite
extension; the truth value of \(X^k o\) does not depend on the extension.

\subsection{The TempoBench-Compatible Effect Class}

The clean restricted effect class uses the outputs selected by
\(\Frame_{\mathrm{TB}}\). Let
\[
  O_{\Frame}\subseteq O
\]
be those selected outputs. Then:
\[
  \mathcal{E}_{TB}(\Frame_{\mathrm{TB}})
  =
  \{\,X^k o
  \mid
  o\in O_{\Frame},\ 0\leq k<|\tau_{\le n}|,\ 
  o\in\alpha_k\,\}.
\]

In plain English:
\[
\boxed{
  E \text{ should be ``this output occurred at this exact timestep.''}
}
\]

Here \(o\in O\) is a single output atomic proposition, and \(k\) is the timestep
where that output is observed. In the TempoBench-compatible case, \(o\) must
also be selected by \(\Frame_{\mathrm{TB}}\), so \(o\in O_{\Frame}\). For
example,
\[
  E_{\Frame} = X^3 grant
\]
means that \(grant\) occurs at timestep \(3\).

This is a special case of the broader effect definition. The broad definition
allows arbitrary temporal facts. The TempoBench-compatible frame
\(\Frame_{\mathrm{TB}}\) keeps only positive selected-output occurrences at
fixed finite positions.

\subsection{The Corresponding Cause Form}

For a TempoBench-compatible effect
\[
  E_{\Frame}=X^k o,
\]
the cause is finite-horizon. It should describe input conditions at timesteps
\[
  0,1,\ldots,k
\]
that are necessary for output \(o\) to occur at timestep \(k\).

One formula form is:
\[
  C_E =
  \bigwedge_{j=0}^{k} X^j c_j,
\]
where each \(c_j\) is a conjunction of input literals. An input literal is
either an input proposition or its negation. For example:
\[
  in_1,\qquad
  \neg in_3,\qquad
  in_1\land \neg in_3.
\]

An equivalent set form is:
\[
  \Gamma_E = \{(t,\ell)\mid 0\leq t\leq k
  \text{ and } \ell \text{ is a required input literal at time } t\}.
\]
For example,
\[
  \Gamma_E=\{(1,in_2),(3,\neg in_0)\}
\]
corresponds to the bounded temporal formula
\[
  C_E = X\,in_2 \land X^3\neg in_0.
\]

The reverse-causal test asks whether these input constraints are necessary. If
one of the required constraints is removed or flipped while the counterfactual
trace remains as close as possible to the actual trace, then the target output
should disappear. If the output still occurs, the proposed cause was not
necessary.

\subsection{What to Exclude for Direct TempoBench Alignment}

To stay directly aligned with the TempoBench-style task, the initial effect
class should exclude:
\[
  F\,grant,
\]
\[
  G(req\rightarrow F\,grant),
\]
\[
  GF\,ready,
\]
\[
  \mathrm{Acc}_{\mathrm{parity}},
\]
and
\[
  F(state=s_7).
\]

These are all meaningful effects for the broader reverse-parity project, but
they change the task. The first three are temporal effects that are not tied to
one output at one finite timestep. The parity effect is an infinite-play success
condition. The state effect is internal-reachability-focused unless the state is
explicitly exposed.

For direct TempoBench compatibility, the restriction should be:
\[
\boxed{
  E_{\Frame} = X^k o
  \quad\text{where } o\in O_{\Frame}
  \text{ and } o\in\alpha_k.
}
\]

It is also best to start with positive output occurrences only. That means
effects such as \(X^k o\), not \(X^k\neg o\). Asking why an output did not occur
is a non-event causality question, which is usually more ambiguous because many
different facts can prevent an output.

\subsection{Point-Effect and Property-Effect Causality}

It is useful to name the restricted and broad versions separately.

\paragraph{Point-effect temporal causality.}
This is the TempoBench-compatible version:
\[
  E_{\Frame}=X^k o.
\]
The effect is one positive output at one timestep. The cause is a finite set of
timestep-indexed input constraints.

\paragraph{Property-effect temporal causality.}
This is the broader reverse-parity version. The effect \(E\) may be a bounded
or unbounded temporal property, such as a response property, recurrence
property, state reachability property, transition property, or parity acceptance
property.

A simple progression is:
\begin{enumerate}[label=\arabic*.]
  \item Output at time: \(X^k o\). This is TempoBench-compatible.
  \item Output conjunction at time: \(X^k(o_1\land o_2)\). This is a small
        extension.
  \item Bounded output pattern: \(X^2 a\land X^5 b\). This is an extension.
  \item Bounded response: \(X^j req\rightarrow X^k grant\). This is an
        extension.
  \item State reachability: \(F(state=s)\). This is a different task.
  \item Full LTL effect: \(G(req\rightarrow F\,grant)\). This is not
        TempoBench-compatible.
  \item Parity acceptance: \(\mathrm{Acc}_{\mathrm{parity}}\). This is not
        TempoBench-compatible, but it is reverse-parity-specific.
\end{enumerate}

Thus the broad reverse-parity idea can explain many kinds of effects, but the
TempoBench-compatible frame \(\Frame_{\mathrm{TB}}\) selects the narrow class
\[
\boxed{
  \mathcal{E}_{TB}(\Frame_{\mathrm{TB}}) = \{\,X^k o : o\in O_{\Frame}\,\},
}
\]
with \(k\) fixed by the observed finite trace and with the additional
requirement that the output actually occurs:
\[
  o\in\alpha_k.
\]

Under this restriction, the reverse-parity method becomes an alternate formal
way to solve a TempoBench-style causal task: compute the smallest input
constraints whose removal prevents the target output from occurring at the
chosen timestep.

\section{TempoBench Alignment from the Reverse Question Onward}
\label{app:tempobench-alignment}

This appendix explains how the reverse-parity procedure should be changed when
the goal is to match a TempoBench-style temporal causal evaluation (TCE)
task. The purpose is not to replace the general construction. The purpose is to
specialize it so that the solver receives the same kind of information
TempoBench gives, solves the same kind of finite-horizon causal problem, and
returns a \emph{necessity}-based answer of the same shape. One caveat throughout:
TempoBench labels TCE with CORP, a Halpern--Pearl \emph{actual-causation} tool
(in Halpern's \emph{modified} variant, with contingencies that reset outputs to
their actual-trace values, verified against the primary sources in the
companion verification suite),
whereas this construction computes counterfactual \emph{necessity}; the two differ
on overdetermined effects (Section~\ref{sec:overdetermination}), so what follows is
an \emph{alternate}, necessity-based account of the TCE task, not a reproduction of
CORP's answer key. Any corpus-scale alignment claim must be stated in CORP's own
(modified) discipline.

This appendix is based on the TempoBench paper and public implementation.
TempoBench has two tasks: temporal trace evaluation (TTE), which asks whether a
finite trace is accepted by a supplied automaton, and temporal causal
evaluation (TCE), which asks for the minimal input constraints needed for a
specified output at a specified timestep.\footnote{TempoBench paper:
\url{https://arxiv.org/abs/2510.27544}. Public implementation:
\url{https://github.com/nik-hz/tempobench}.}

\subsection{What TempoBench Gives the Solver}

A general reverse-parity problem may require choosing an effect, choosing an
actual trace, compiling monitors, and deciding what counterfactuals are allowed.
A TempoBench-style instance fixes most of this information in advance.

The solver is given:
\begin{enumerate}[label=\arabic*.]
  \item an automaton \(A_{\mathrm{HOA}}\) in HOA format;
  \item the list of APs used by the automaton;
  \item a finite observed trace \(\tau_{\le n}\);
  \item one or more effects to analyze, usually written in the form
        \(\texttt{XXX }g\), meaning \(X^3 g\);
  \item enough information in the prompt or dataset to know which output is the
        target effect and which APs are eligible input causes.
\end{enumerate}

The solver is not asked to synthesize a controller, choose a trace, choose an
effect, discover the original LTL specification, or explain an infinite parity
condition. It is asked to perform finite-horizon causal credit assignment over
the provided execution.

\subsection{Specializing the Reverse Question (Section~\ref{sec:reverse-question})}

The general reverse question is
\[
  M+\varphi+\Frame+E+\tau
  \longrightarrow
  C^\star .
\]
For TempoBench-style TCE, this becomes
\[
  A_{\mathrm{HOA}}
  +
  \Frame_{\mathrm{TB}}
  +
  \tau_{\le n}
  +
  E_{\Frame}=X^k o
  \longrightarrow
  C^\star .
\]

The meaning of each component is restricted:
\begin{description}[leftmargin=2.8cm,style=nextline]
  \item[\(A_{\mathrm{HOA}}\)]
  The supplied finite-state automaton. This replaces the need to synthesize or
  reconstruct a controller.

  \item[\(\Frame_{\mathrm{TB}}\)]
  The benchmark frame. It says that the only target effects are positive output
  occurrences at fixed finite timesteps, and that the causes should be
  timestep-indexed input constraints.

  \item[\(\tau_{\le n}\)]
  The supplied finite trace. The solver does not choose it.

  \item[\(E_{\Frame}=X^k o\)]
  The selected point effect. It says that output \(o\) occurs at timestep \(k\).

  \item[\(C^\star\)]
  The minimal cause, represented as a finite set of timestep-indexed
  input literals.
\end{description}

So the TempoBench-compatible reverse question is:
\[
  \text{Which input facts in the provided prefix were necessary for output }
  o \text{ at timestep } k?
\]

\subsection{Specializing Step~\ref{step:freeze}: Freeze the System Strategy}

In the general construction, this step freezes the system strategy \(M\) and
forms the product \(P=M\times A_E\) with a deterministic parity automaton for
the effect. In a TempoBench-style instance the correspondence is concrete, but
it requires us to be precise about what kind of object \(A_{\mathrm{HOA}}\) is.

\begin{remark}[What \(A_{\mathrm{HOA}}\) is, and the faithfulness assumption]
\label{rem:hoa}
TempoBench presents the system as a finite-state \emph{acceptor}
\(A_{\mathrm{HOA}}=(Q,2^{I\cup O},\delta,q_0,F)\) over the alphabet of complete
valuations \(2^{I\cup O}\): it reads valuation sequences
\(\alpha_0\alpha_1\cdots\) (each \(\alpha_t\subseteq I\cup O\), inputs \emph{and}
outputs together) and accepts a finite trace when its run from \(q_0\) ends in an
accepting state of \(F\) (this is exactly TempoBench's trace-acceptance, or TTE,
semantics). It is presented as an acceptor over \(I\cup O\) rather than as an
input/output transducer: it does not read inputs and emit outputs. Its language
is a set of valuation sequences, and we take \(\Lang(A_{\mathrm{HOA}})\) to be the
set of admissible system behaviors. In this sense \(A_{\mathrm{HOA}}\) plays the
role of \(\Lang(M)\) in the general construction (it supplies the admissible
behaviors), \emph{not} the role of the effect monitor \(A_E\).

For the reduction to be faithful we make two assumptions, both checkable against
the benchmark artifact. \textbf{(F1) Behavior model.}
\(\Lang(A_{\mathrm{HOA}})\) is exactly the set of executions of the system under
analysis, so ``accepted by \(A_{\mathrm{HOA}}\)'' coincides with ``is a genuine
behavior''; if instead it were a pure specification monitor \(A_\varphi\) whose
language is the set of \emph{specification-satisfying} traces, identifying it with
the system model would conflate \(A_\varphi\) with \(M\), and we exclude that
reading. \textbf{(F2) Functionality.} Since TempoBench's automata are obtained by
reactive synthesis from a controller specification, the behavior set is
\emph{functional in the inputs}: each input sequence has at most one accepted
\(I\cup O\) completion (the automaton is input-deterministic). This is
\emph{consistent} with the acceptor framing --- an acceptor over \(I\cup O\) may
perfectly well recognize a functional input/output relation --- and it is exactly
the property the operational reading of Step~\ref{step:check} below uses to talk
about ``the'' output at each timestep. We state (F1)--(F2) as explicit,
benchmark-level assumptions rather than building them into the data type, so that
the reduction's correctness is hostage to stated facts about TempoBench rather
than to a silent convenience. The causal task itself (inputs as candidate causes,
outputs as effects) presupposes (F2). \textbf{We have not independently verified
(F1)--(F2) against the released TempoBench HOA artifacts}; whether those automata
are input-deterministic system models (as opposed to specification monitors) is a
checkable property of the benchmark that we leave to be confirmed against the
artifacts. Accordingly the alignment of Appendix~\ref{app:tempobench-alignment} is
stated \emph{conditionally on (F1)--(F2)}: if they hold, the reduction
(Proposition~\ref{prop:tb-reduction}) is exact; if the supplied automata are
specification monitors, (F1) fails and the alignment does not apply.
\end{remark}

Given Remark~\ref{rem:hoa}, the point effect \(E_{\Frame}=X^k o\) needs no
separate parity monitor at all, since it is a bounded membership check on the
trace at timestep \(k\) (namely \(o\in\alpha_k\)). There is therefore no product
to build and no parity acceptance to track: the ``semantic arena'' reduces to
\(A_{\mathrm{HOA}}\) together with the bounded effect, and the solver receives
\(A_{\mathrm{HOA}}\) directly.

Therefore this step becomes:
\begin{enumerate}[label=\arabic*.]
  \item parse the HOA automaton;
  \item parse the AP list;
  \item identify the target output \(o\) from the effect \(X^k o\);
  \item identify the input APs that may appear as causes;
  \item parse the supplied finite trace \(\tau_{\le n}\);
  \item if state information is provided, verify the finite transition sequence
        \[
          q_t\xrightarrow{\alpha_t}q_{t+1}
        \]
        for the relevant prefix.
\end{enumerate}

The important constraint is that there is no new synthesis problem here. The
automaton is treated as the fixed system model. The reverse procedure should use
the same information a TempoBench solver sees: the automaton, the trace, the AP
names, and the target effect.

\subsection{Specializing Step~\ref{step:trace}: Choose the Actual Trace}

In the general construction, an actual trace must be chosen. In a
TempoBench-style instance, this step is already fixed by the benchmark. The
solver is given a finite observed trace
\[
  \tau_{\le n}=\alpha_0,\ldots,\alpha_n
\]
generated by the automaton. Therefore, the solver does not choose \(\tau\).
Instead, it parses the provided trace, verifies the state sequence if needed,
and restricts attention to the prefix
\[
  \tau_{\le k}
\]
when the target effect has the form
\[
  E_{\Frame}=X^k o.
\]

The reason for this restriction is simple: later timesteps cannot be causes of
an output at timestep \(k\) in this finite-horizon task. The relevant causal
window is:
\[
  0,1,\ldots,k.
\]
If the trace contains a cycle marker, such as a \(\texttt{cycle}\) annotation
with a loop index, that marker belongs to the trace representation but is not
the target of the finite causal query. The causal query only needs the finite
prefix up to the effect.

\subsection{Specializing Step~\ref{step:counterfactuals}: Define the Allowed Counterfactuals}

In the general construction, one must decide whether counterfactuals may change
inputs, outputs, transitions, priorities, or strategy commitments. TempoBench
uses a narrower frame. This step instantiates
\(\mathcal{I}_{\Frame_{\mathrm{TB}}}\), the intervention component of the
TempoBench-compatible causal frame.

For TempoBench-style TCE, use finite input-only counterfactuals:
\begin{enumerate}[label=\arabic*.]
  \item the automaton \(A_{\mathrm{HOA}}\) is fixed;
  \item the target output \(o\) is not manually intervened on;
  \item transitions are not changed;
  \item priorities or acceptance conditions are not changed;
  \item strategy-level commitments are not changed;
  \item candidate causes are input literals at timesteps \(0,\ldots,k\).
\end{enumerate}

A cause is represented as a finite set
\[
  \Gamma_E=\{(t,\ell)\mid 0\leq t\leq k
  \text{ and } \ell \text{ is an input literal}\}.
\]
The corresponding bounded temporal formula is
\[
  C_E=\bigwedge_{(t,\ell)\in\Gamma_E} X^t\ell.
\]

The benchmark allows positive and negative constraints. For example, a cause
may require \(in_2\) at timestep \(1\), or it may require \(\neg in_0\) at
timestep \(3\). If no input constraint is needed at a timestep, the output
format may record that timestep as having no constraints.

\subsection{Specializing Step~\ref{step:similarity}: Define Similarity}

The general construction needs a similarity relation over infinite traces. A
TempoBench-style instance should not use that full infinite-trace machinery.
The trace is finite, and the target effect occurs at a fixed timestep.

The closest counterfactuals should therefore be finite-prefix perturbations of
\[
  \tau_{\le k}.
\]
They should change only the input constraints being tested and should keep the
rest of the prefix as close as possible to the provided trace.

For this setting, the most natural closeness rule is:
\[
  \text{fewer changed timestep-literal constraints means closer,}
\]
counted in changed \emph{input cells}---under (F2) the outputs are determined by
the inputs, so only input changes are counterfactual choices, and these (not the
full-valuation positions of the general order) are what we count.
Minimality is then set-theoretic. A cause is preferred when it uses fewer or
strictly necessary timestep-indexed literals. The benchmark does not ask for a
semantically weakest \(\omega\)-regular property, a smallest parity automaton,
or a shortest LTL formula.

\subsection{Specializing Step~\ref{step:check}: Check a Candidate Cause}

In the general reverse-parity game, an adversary tries to violate the candidate
cause while preserving the effect. In TempoBench-style TCE, the same idea is
finite and concrete.

Given a candidate \(\Gamma_E\), first check actuality:
\begin{enumerate}[label=\arabic*.]
  \item every listed literal \((t,\ell)\in\Gamma_E\) holds in the supplied trace
        at timestep \(t\);
  \item the target output holds at timestep \(k\):
        \[
          o\in\alpha_k.
        \]
\end{enumerate}

Then check counterfactual dependence. The adversary tries to remove or flip a
listed input constraint while keeping the rest of the relevant prefix close to
the original trace. The candidate is valid only if the target effect disappears:
\[
  X^k o
  \text{ no longer holds under the allowed counterfactual.}
\]

Operationally, this is a finite reachability check over \(A_{\mathrm{HOA}}\), not
an infinite parity acceptance check. Since \(A_{\mathrm{HOA}}\) is a finite-state
acceptor over complete valuations (Remark~\ref{rem:hoa}), the test is phrased as
the existence of an accepted valuation sequence: after imposing the perturbed
input constraints, does there exist a finite valuation sequence
\(\alpha_0'\cdots\alpha_n'\) such that (i) it is accepted by \(A_{\mathrm{HOA}}\)
in the sense made precise below (its run ends in \(F\)), (ii) its inputs match
the perturbed input constraints, (iii) it agrees with \(\tau_{\le n}\) elsewhere
as the closeness order requires, and (iv) \(o\in\alpha_k'\)? If such a sequence
exists, the effect survives the intervention and the proposed cause is \emph{not}
necessary; if no such sequence exists, the proposed cause passes the
counterfactual test. Under the functionality assumption~(F2) of
Remark~\ref{rem:hoa} (each input sequence admits at most one accepted completion),
this existence question coincides with the familiar reading ``simulate the
perturbed inputs and read off whether \(o\) appears at time \(k\)''.

\subsection{Specializing Step~\ref{step:search}: Synthesize the Minimal Cause}

The general construction offers several possible synthesis routes: full
automata-theoretic cause synthesis, CEGIS over LTL sketches, or causal dominion
extraction. For TempoBench-style TCE, the output should be much narrower.

The result should be a finite-horizon cause:
\[
  C^\star =
  \bigwedge_{(t,\ell)\in\Gamma_E^\star} X^t\ell,
\]
where \(\Gamma_E^\star\) is the minimum set of timestep-indexed input literals
needed for \(X^k o\).

TempoBench's public prompt style expects a JSON-like answer organized by effect
and timestep. For example, an effect written as \(\texttt{XXX }g\), meaning
\(X^3 g\), may have an answer shaped like:
\[
\begin{array}{ll}
  \text{timestep }0: & \text{no constraints},\\
  \text{timestep }1: & \text{no constraints},\\
  \text{timestep }2: & \text{no constraints},\\
  \text{timestep }3: & r.
\end{array}
\]
The mathematical object behind that answer is
\[
  \Gamma_E^\star=\{(3,r)\}.
\]

Thus, for TempoBench alignment, synthesis should not return a broad LTL
explanation such as \(G(req\rightarrow F\,grant)\), and it should not return a
strategy-level invariant. It should return the minimal finite set of input
constraints, placed at the correct timesteps.

\subsection{Trace-Level, Not Strategy-Level}

The trace-level versus strategy-level distinction is important. TempoBench TCE
is a trace-level point-effect task. The question is not:
\[
  \text{Why does the whole controller realize } \varphi?
\]
It is:
\[
  \text{Which provided input facts caused this output at this timestep?}
\]

Therefore, the TempoBench-compatible version should use:
\[
  E_{\Frame}=X^k o
\]
and should return a finite cause over \(\tau_{\le k}\). It should not try to
explain the whole LTL specification, a full parity-winning condition, or the
controller's entire strategy.

\subsection{Formal Specialization}

Rather than informally ``reading'' the general \(\omega\)-regular machinery as
finite, we give the finite-horizon objects explicitly, \emph{entirely in
finite-trace terms}; no infinitary acceptance condition for \(A_{\mathrm{HOA}}\)
is needed or used. Fix the supplied automaton \(A_{\mathrm{HOA}}\), the finite
trace \(\tau_{\le n}\), and the point effect \(E_{\Frame}=X^k o\) with \(k\le n\).
The candidate space \(\mathcal{K}_{\Frame_{\mathrm{TB}}}\) is the finite set of
conjunctions of timestep-indexed input literals over the window \(0,\ldots,k\),
identified with their constraint sets
\(\Gamma\subseteq\{0,\ldots,k\}\times \mathrm{Lit}(I)\), where
\(\mathrm{Lit}(I)=\{i,\neg i\mid i\in I\}\):
\[
  C_\Gamma=\bigwedge_{(t,\ell)\in\Gamma}X^t\ell .
\]
There are at most \(2|I|(k+1)\) such literals, hence
\(|\mathcal{K}_{\Frame_{\mathrm{TB}}}|\le 2^{2|I|(k+1)}\) candidate constraint
sets: the candidate space is finite but \emph{exponential} in the window, not
linear.

\paragraph{Finite-horizon admissible set.}
We make finite acceptance precise, since the supplied \(A_{\mathrm{HOA}}\) is a
finite-state acceptor with accepting set \(F\) (Remark~\ref{rem:hoa}). A finite
valuation sequence \(\beta=\beta_0\cdots\beta_n\) is \emph{accepted} iff it has a
run
\[
  q_0\xrightarrow{\beta_0}q_1\xrightarrow{\beta_1}\cdots
  \xrightarrow{\beta_n}q_{n+1}
  \quad\text{with } q_{n+1}\in F ,
\]
i.e.\ all transitions are defined and the run ends in an accepting state; this is
TempoBench's TTE acceptance and is the single, fixed meaning of ``accepted'' used
throughout this appendix. (Under the functionality assumption (F2) of
Remark~\ref{rem:hoa}, the run is unique once the inputs are fixed.) Writing
\(\Lang_{\le n}(A_{\mathrm{HOA}})\) for the set of accepted length-\((n{+}1)\)
sequences, the bounded admissible set anchored at \(\tau_{\le n}\) is
\[
  \Adm^{\le k}_{A,\Frame_{\mathrm{TB}}}
  =\Lang_{\le n}(A_{\mathrm{HOA}}).
\]
Because \(X^k o\) depends only on positions \(0,\ldots,k\), it suffices to track
prefixes through timestep \(k\); we keep the full length \(n\) only so that
``accepted'' refers to a complete run of the supplied trace.

\paragraph{Finite-horizon closeness and removals.}
For a candidate \(\Gamma\), a removal is an admissible \(\beta\) that violates
\(C_\Gamma\) (i.e. some \((t,\ell)\in\Gamma\) fails in \(\beta\)). The closeness
order counts changed timestep-literal constraints, and the closest removals are
\[
  \Close^{\le k}_{\tau_{\le n},A}(\neg C_\Gamma)
  =\min_{\#\text{changes}}
   \{\beta\in\Adm^{\le k}_{A,\Frame_{\mathrm{TB}}}\mid \beta\not\models C_\Gamma\}.
\]

\paragraph{Finite-horizon valid causes.}
\[
\begin{aligned}
  \mathsf{Causes}^{\le n}_{A,\Frame_{\mathrm{TB}},k}(X^k o)
  =\{\,\Gamma\mid{}&
  \tau_{\le n}\models C_\Gamma,\ o\in\alpha_k,\\
  &\Close^{\le k}_{\tau_{\le n},A}(\neg C_\Gamma)\neq\emptyset,\\
  &\forall\,\beta\in\Close^{\le k}_{\tau_{\le n},A}(\neg C_\Gamma).\ o\notin\beta_k
  \,\}.
\end{aligned}
\]
A minimal cause is \(C^\star=C_{\Gamma_E^\star}\), where \(\Gamma_E^\star\) is any
subset-minimal element of
\(\mathsf{Causes}^{\le n}_{A,\Frame_{\mathrm{TB}},k}(X^k o)\) (several incomparable
ones may exist); minimality is by set inclusion over \(\Gamma\), not by LTL
formula size, temporal depth, or automaton size.

\begin{remark}[Validity is not monotone; greedy deletion is unsound in general]
\label{rem:nonmonotone}
Membership in \(\mathsf{Causes}^{\le n}_{A,\Frame_{\mathrm{TB}},k}\) is not
monotone under \(\subseteq\) on \(\Gamma\): enlarging \(\Gamma\) can make a
candidate valid (a removal now has to flip more, and the cheapest removal may stop
preserving the effect) or invalid (the closest \(\neg C_\Gamma\)-removal may flip
an \emph{irrelevant} added literal while keeping the target output). Hence the
naive greedy procedure ``delete literals one at a time while validity holds'' is
not guaranteed to reach a subset-minimal valid cause, and computing
\(\Gamma_E^\star\) requires genuine search over the exponential space --- by
exhaustive enumeration, a hitting-set/MaxSAT encoding, or CEGIS over the finite
constraints. Each individual validity check is polynomial (finite reachability
over \(A_{\mathrm{HOA}}\) on the window \(0,\ldots,k\)); it is the \emph{search}
that is hard. If one additionally posits that validity is upward- or
downward-closed for a particular automaton, greedy deletion becomes sound, but we
do not assume this in general.
\end{remark}

\begin{proposition}[Finite specialization --- a definitional correspondence]
\label{prop:tb-reduction}
This is a \emph{definitional} restatement, not a deep theorem: the finite-horizon
objects above were \emph{defined} to mirror the general clauses, and the content is
that the mirroring is exact and needs no infinitary acceptance condition on
\(A_{\mathrm{HOA}}\); we record it as a proposition only to state that
correspondence precisely. Let
\(\Frame_{\mathrm{TB}}\) be the input-only frame with candidate space
\(\mathcal{K}_{\Frame_{\mathrm{TB}}}\), background \(B_{\Frame}=\True\), bounded
admissible set \(\Adm^{\le k}_{A,\Frame_{\mathrm{TB}}}=\Lang_{\le n}(A_{\mathrm{HOA}})\)
(finite acceptance as defined above), the change-counting closeness order on the
window \(0,\ldots,k\), and subset minimality, and assume (F1)--(F2) of
Remark~\ref{rem:hoa}. Interpret the general clauses of \(\mathsf{Causes}\)
(Section~\ref{sec:formal}) over finite traces with \(\LTL_f\) semantics on the
window. Then for every \(\Gamma\),
\(\Gamma\in\mathsf{Causes}^{\le n}_{A,\Frame_{\mathrm{TB}},k}(X^k o)\) iff
\(C_\Gamma\) satisfies the four valid-cause clauses on \(\tau_{\le n}\) under
\(\Frame_{\mathrm{TB}}\). In particular, the finite construction \emph{is} the
general definition restricted to the point effect; no claim is made about, and
nothing depends on, infinite extensions of \(\tau_{\le n}\) or an
\(\omega\)-language of \(A_{\mathrm{HOA}}\).
\end{proposition}

\begin{proof}
Each of the four clauses of \(\mathsf{Causes}\) refers only to positions
\(0,\ldots,k\): actuality is \(\tau_{\le n}\models C_\Gamma\) and \(o\in\alpha_k\);
non-vacuity is non-emptiness of the admissible \(\neg C_\Gamma\)-removals; and
effect-failure quantifies over the closest such removals, asking \(o\notin\beta_k\).
The effect \(X^k o\) and every literal \(X^t\ell\) (\(t\le k\)) is determined by
the length-\((k{+}1)\) prefix under \(\LTL_f\) semantics, so each clause is a
finite statement about prefixes in \(\Lang_{\le n}(A_{\mathrm{HOA}})\). These are
verbatim the clauses of \(\mathsf{Causes}^{\le n}_{A,\Frame_{\mathrm{TB}},k}\):
the admissible sets coincide by definition, the closeness order is in both cases
the count-only coarsening of \(\preceq_\tau\) on the window (the change-counting
preorder of Step~\ref{step:similarity}, so that \(\Close\) is the full set of
minimum-count removals and the effect-failure clause quantifies over all of them),
and (F2) makes ``\(o\) at position \(k\)'' a well-defined function of the
admissible inputs. Hence the two memberships are literally the same predicate. No
infinite object appears in any clause, which is why the reduction needs no
\(\omega\)-acceptance for \(A_{\mathrm{HOA}}\).
\end{proof}

\subsection{Evaluation Compatibility}

TempoBench evaluates causal answers at two granularities.
\begin{description}[leftmargin=2.8cm,style=nextline]
  \item[AP-level]
  A prediction can receive partial credit for correctly identifying individual
  atomic propositions at a timestep.

  \item[TS-level]
  A timestep receives credit only when the whole set of constraints predicted
  for that timestep matches the ground truth.
\end{description}

This means the reverse-pipeline output should preserve timestep structure. A
flat formula such as
\[
  X\,in_2\land X^3\neg in_0
\]
is mathematically fine, but a TempoBench-compatible answer should also preserve
the grouped representation:
\[
\begin{array}{ll}
  1: & in_2,\\
  3: & \neg in_0.
\end{array}
\]
That grouping matches how AP-level and TS-level scoring are computed.

TempoBench also controls difficulty with finite, observable features: effect
depth \(k\), number of automaton states, transition count, trace length, causal
input count, and number of unique inputs in the trace. The pipeline of this
appendix should preserve those features. It should not add hidden information from the
original LTL specification or from unrelated files, because doing so would make
the adapted task easier than the benchmark task.

\subsection{Checklist for TempoBench Alignment}

To make the reverse-parity pipeline match TempoBench-style TCE, use the
following restrictions:
\begin{enumerate}[label=\arabic*.]
  \item Use the provided HOA automaton as \(A_{\mathrm{HOA}}\).
  \item Use the provided finite trace \(\tau_{\le n}\).
  \item Use the provided effect \(E_{\Frame}=X^k o\).
  \item Restrict causal attention to \(\tau_{\le k}\).
  \item Allow only finite input-literal counterfactuals.
  \item Do not intervene directly on outputs.
  \item Do not change transitions, priorities, or strategy commitments.
  \item Minimize by necessary timestep-indexed input constraints.
  \item Return causes grouped by timestep.
  \item Treat broader effects, liveness, parity acceptance, and strategy-level
        explanations as extensions beyond TempoBench.
\end{enumerate}

With these restrictions, the reverse-parity framework becomes a formal
explanation procedure for the same kind of problem TempoBench tests: determining
which earlier input facts were necessary for a later output in a provided finite
automaton trace.

\end{document}
